\section{Introduction}\label{sec:intro}
Magnetic Resonance Spectroscopy (MRS) is a powerful non-invasive tool used to measure metabolite concentrations \textit{in vivo}, providing insights into various biochemical processes in health and disease. Recently, artificial intelligence (\AI) and deep learning (\DL) approaches have shown promise for metabolite quantification \cite{Das2017, Das2018a, Olliverre2018, Iqbal2021, Schmid2023, Shamaei2023} and denoising tasks \cite{Lei2021, Chen2024, Lam2020, Li2022}, with the potential to complement or replace traditional iterative methods. As these methods mature, the central evaluation question shifts from whether they can produce accurate estimates to whether their behavior is sufficiently characterized to support trust within a defined operating range. Despite their rapid development, the evaluation and reporting of such \DL models remain highly heterogeneous. Many studies rely on incomplete, poorly justified, or inconsistent statistical analyses, limiting reproducibility, comparability, and the interpretation of reported performance. Similar concerns have been highlighted more broadly in \AI for digital medicine \cite{Muller2026}, where methodological shortcomings, inadequate validation strategies, and insufficient consideration of clinical utility can result in models that appear successful during development but fail to provide meaningful benefits in practice. Furthermore, predictive performance alone does not establish that a model is reliable across the conditions in which it will be used or sufficiently comparable to established methods for the intended application, emphasizing the need for rigorous characterization, transparent reporting, and evaluation frameworks that extend beyond conventional accuracy metrics.

As Markowetz highlights in \cite{Markowetz2024}, most published models are never used in clinical practice. Markowetz points out that "success in the clinic is hard earned" because healthcare providers both have an unequivocal responsibility to their patients and work in highly regulated environments. The regulatory approval processes of the U.S. Food and Drug Administration (FDA) and the European Medicines Agency (EMA) reflect the high standards for proving a drug, device, or algorithm is safe. The level of validation typically considered sufficient in academia often falls short of the rigorous standards for regulatory approval and adoption into clinical practice. This challenge is not unique to \MRS. Recently, Lekadir et al.~\cite{Lekadir2025} established the FUTURE-AI international consensus guidelines for trustworthy and deployable artificial intelligence in healthcare, identifying a broad set of technical, clinical, ethical, legal, and deployment considerations that must be addressed throughout the entire \AI lifecycle. Their recommendations span seven domains, including fairness, robustness, explainability, usability, and traceability. While FUTURE-AI provides a holistic framework for the development, validation, deployment, and monitoring of healthcare \AI systems, many of its recommendations must still be interpreted and implemented within the unique methodological constraints of specific fields such as \MRS.

The present work addresses this methodological gap by establishing a common framework at the intersection of \MRS, deep learning, and statistics for the Nuclear Magnetic Resonance (\NMR) and machine learning (\ML) communities. The framework is organized around increasingly consequential questions concerning data validity, model performance, robustness, uncertainty, agreement, and fitness for intended use. "Laying the Foundation" provides a comprehensive toolbox of practical methods and guidelines for the development, characterization, and evaluation of models for \DL-based model fitting and metabolite quantification. It establishes a shared background for \MR spectroscopy and \DL, with the goal of making terminology, assumptions, and current developments accessible to a broad interdisciplinary audience. It then addresses data considerations that directly affect model validation, including the use of synthetic data, data curation, annotation accuracy, sampling distributions, SNR definitions, out-of-distribution errors, dataset splitting, dataset size, and sample size requirements. These topics define the conditions under which \DL-based model fitting and metabolite quantification models can be trained, tested, and interpreted, and whether the underlying data support valid claims about model performance.

With these data requirements established, the framework turns to statistical considerations that are often overlooked in \DL-based \MRS studies, including assumptions of linearity and nonlinearity, qualitative and quantitative assessments, limitations of parametric metrics, resampling-based implementations, statistical significance, effect size, clinical relevance, and the role of distributions in \DL. Building on this statistical basis, it then introduces characterization metrics for evaluating error, dispersion, regression performance, model benchmarking, distributional similarity, and multivariate prediction behavior. These metrics are intended to describe both model-level and metabolite-level performance, while recognizing that no individual metric is sufficient to establish model trustworthiness.

Subsequent sections translate these metrics into practical analyses and visualizations, including correlation analysis, calibration, model- and component-level comparisons, SNR-dependent performance, nonparametric confidence bounds, and remaining multivariate fitting parameters. These analyses extend evaluation beyond pooled or average performance by examining how model behavior changes across relevant operating conditions. This work then addresses uncertainty and interpretability, with an emphasis on spectral modeling, predictive uncertainty analysis, decomposed uncertainties, component-level uncertainty, and the limitations of using Cram{\`e}r--Rao lower bounds as a reference for \DL models. Here, the relevant question is whether uncertainty estimates provide reliable information about prediction error rather than merely producing an uncertainty output. The final validation stage asks whether a \DL method agrees sufficiently with established methods for the intended use. Method comparison studies are therefore discussed through agreement analysis, regression-based comparisons, multivariate comparisons, and equivalence testing, followed by reporting and reproducibility recommendations that aim to improve comparability, transparency, and resource sharing across the \MRS and \DL communities.

Although the primary focus is on \DL-based model fitting, the presented framework is broadly applicable to other \MRS-related tasks. Because modifications of spectral data often influence downstream quantification results, the proposed methodologies can also be used to evaluate preprocessing methods, such as denoising and artifact removal, and reconstruction methods, for accelerated acquisitions, with respect to their effect on subsequent metabolite estimation. More broadly, the statistical principles, validation strategies, uncertainty analyses, and reporting recommendations described throughout this article are largely independent of the underlying model architecture or application domain. Consequently, many of the concepts presented herein can be adapted to a wide range of quantitative problems in \NMR, magnetic resonance, and quantitative machine learning.

\section{Background}\label{sec:background}
In recent years, deep learning (\DL) has become increasingly integrated into magnetic resonance spectroscopy (MRS) workflows. It is currently being applied to every part of \MRS protocols, including quantification \cite{Shamaei2023, Chen2023a, Wu2026}, fitting \cite{Shamaei2023, Shamaei2025}, denoising \cite{Wang2023}, artifact removal \cite{Gudmundson2023a, Kyathanahally2018, Tapper2021, Merkofer2025}, and spectral quality classification \cite{Kyathanahally2018a,Jang2021}. Detailed reviews of the use of \AI throughout \MRS pipelines have been provided by van de Sande et al. \cite{VandeSande2023} and LaMaster et al. \cite{LaMaster2026}. The need for and potential of \DL applications, particularly where traditional processing and quantification methods are limited, shape the current development and exploratory trends. The most common current trends include: 1) accelerating and automating routines for data processing, nuisance signal removal, and model fitting; 2) improving \DL fitting models by incorporating physics-informed priors; and 3) fast and automatic quality assessment for \MRS data, important for determining not simply whether a method performs well, but whether its outputs can be trusted under the conditions in which it is intended to operate. 

Currently, \MRS is also trending towards increasingly higher resolution \MRSI, which means that the amount of data and the number of spectra to be processed are skyrocketing. Most such software packages process data individually, sequentially, and iteratively, exacerbating the processing time demands. \MRSI acquisitions commonly acquire thousands of spectra, which is even more expensive both in terms of time and computation. \DL techniques have successfully reduced these processing times from hours to seconds, improving the clinical feasibility \MRSI. This work addresses a complementary obstacle: how to characterize these \DL fitting models rigorously enough to determine whether, where, and under what conditions their outputs can be trusted for clinical use.

This work addresses the intersection of deep learning and \MRS. There are differences in terminology between these two fields. To address these differences for spectroscopists, a dictionary of important terms from the \DL realm is included in the Supporting Information (SI) Section 1. For deep learning researchers in the field of \MRS, readers are strongly encouraged to consult the following series of consensus papers and guidelines issued by multiple \MRS expert groups of the International Society for Magnetic Resonance in Medicine (ISMRM). These papers provide a comprehensive reference library for understanding, planning, acquiring, processing, analyzing, and publishing \MRS studies. The series includes an in-depth reference for \MRS terminology \cite{Kreis2020} and volumes on \MRS clinical method development for the brain \cite{Wilson2019}, B0 shimming \cite{Juchem2021}, advanced spectroscopic imaging \cite{Maudsley2020}, advanced single voxel techniques in humans \cite{Oz2021}, preclinical \MRS for the rodent brain \cite{Lanz2021}, pre-processing, analysis and quantification \cite{Near2021}, water and lipid suppression techniques \cite{Tkac2021}, $^{31}$P and proton \MRS in skeletal muscle \cite{Meyerspeer2020,Krssak2021}, diffusion weighted \MRS \cite{Ligneul2023}, contribution of macromolecules to brain \MRS \cite{Cudalbu2021}, and description of the minimum reporting standard for publishing \cite{Lin2021}. Additionally, the ISMRM's Spectroscopy Study Group's Committee on Code \& Data Sharing released an educational resource \cite{Fleischer2022} called "Getting Started in MRS" to help people new to \MRS. 

While these consensus papers establish methodological standards across many topics in \MRS, a comparable consensus addressing the integration of deep learning and \MRS has, to date, not been presented. To this end, this work aims to bridge terminology, and methodological considerations across both fields while providing a shared basis for the rigorous evaluation of \DL-based \MRS methods.

\subsection{Current State of the Field}\label{sec:background:subsec:SOTF}
The evaluation of deep learning methods for \MRS metabolite quantification has been carried out using a wide variety of metrics, analyses, and visualizations, with little standardization across the literature. Importantly, these approaches characterize different properties of model behavior and therefore do not necessarily support equivalent claims about model performance. Early works \cite{Hiltunen2002,Bhat2006} relied primarily on correlation-based metrics, such as Pearson correlation coefficients and normalized $L_2$ error, alongside visual tools such as Bland-Altman plots and three-dimensional mesh plots of error as a function of noise and phase distortion. These foundational approaches established the convention of reporting per-metabolite performance, a practice that has persisted throughout the literature, though the choice of metric and included metabolites has varied considerably.

Absolute and relative error metrics have been the most commonly reported measures of quantification accuracy. Mean absolute error (\MAE) and mean absolute percent error (\MAPE) appear frequently \cite{Chandler2019, Lee2019b, Lee2020, Long2024, Wu2024}, while symmetric MAPE (\SMAPE) has been favored in several studies \cite{Hatami2018, Shamaei2020, Shamaei2021}. Mean squared error (MSE) and root mean squared error (\RMSE) have also been employed \cite{Das2017, Hiltunen2002, Rizzo2023}, and some works have proposed composite or domain-specific scores, such as the fit quality number (\FQN) \cite{Near2021} \cite{Jeong2024, Wu2024} or a multi-component peak-picking score \cite{Schmid2023}. Regression-based metrics, including slope, intercept, $R^2$, and standard error derived from predicted-versus-ground-truth scatter plots, have been reported in a subset of studies \cite{Chandler2019, Iqbal2021, Rizzo2022, Rizzo2023}, often presented alongside the scatter plots themselves as a means of assessing both accuracy and linearity of response.

Robustness analyses have been used to probe method performance under varying acquisition conditions, most commonly noise level or signal-to-noise ratio (\SNR) \cite{Bhat2006, Hatami2018, Shamaei2021, Lee2022, Long2024}, spectral linewidth \cite{Lee2022, Long2024}, phase distortion \cite{Bhat2006, Shamaei2021}, and macromolecule baseline variation \cite{Hatami2018, Shamaei2021}. These analyses are typically presented graphically, as line or bar plots of a chosen error metric as a function of the varied parameter, though few studies apply formal statistical tests to assess the significance of observed trends. Among the exceptions, \cite{Das2017} employed out-of-bag error analysis, \cite{Songeon2023} applied $t$-tests with Bonferroni correction, and \cite{Jeong2024} reported Mann-Whitney $U$ tests with Bonferroni correction on correlation and error metrics.

Visualization of spectral fits, such as displaying the input spectrum, estimated fit, and residual, has been nearly universal across the literature as a qualitative demonstration of model performance \cite{Gurbani2019, Lee2019b, Shamaei2023, Zhang2023, Wu2024}. Bland-Altman plots comparing \DL estimates to a reference method have been used in a smaller number of works \cite{Bhat2006, Songeon2023, Shamaei2025}, while error box plots or violin plots stratified by metabolite or method have appeared in \cite{Das2017, Jeong2024}. Correlation heatmaps of inter-metabolite estimated concentrations have also been proposed as a means of assessing the plausibility of joint estimates \cite{Shamaei2023}.

Uncertainty quantification has been approached inconsistently across the literature. The majority of reviewed works report no uncertainty estimate for individual predictions \cite{Hatami2018, Chandler2019, Shamaei2020, Iqbal2021, Shamaei2021, Zhang2023, Wu2024}. Several studies include Cramér-Rao lower bounds (CRLBs) solely as a reference baseline from conventional model-fitting methods rather than as an output of the \DL model itself \cite{Rizzo2023, Dziadosz2023, Long2024}. A smaller number of works have introduced \DL-native uncertainty estimates. For example, \cite{Lee2020} trained a separate network to predict concentration error. Subsequently, \cite{Lee2022} employed Monte Carlo dropout to decompose aleatoric and epistemic uncertainty while \cite{Rizzo2022} compared network-derived aleatoric and epistemic uncertainties against \CRLBs using reliability diagrams and calibration error. Finally, \cite{Shamaei2025} reported all three, \CRLB, aleatoric, and epistemic uncertainty, projected into brain space alongside metabolite concentration maps, though the computational details of each were not fully elaborated.

Taken together, the characterization landscape across \DL-based \MRS quantification studies is notably heterogeneous. Metrics, analyses, and visualizations differ substantially between groups, making direct cross-study comparison difficult. More fundamentally, the presence of a reported metric does not by itself establish the broader property required for trustworthiness. Correlation does not establish agreement, pooled error does not establish robustness across relevant operating conditions, and an uncertainty estimate does not establish that the reported uncertainty is calibrated. This n the claims that different forms of evidence can support, motivates the development of a common evaluation framework, which is addressed in the present work.

\section{Data Considerations}
\label{sec:data model aggregate}
Machine learning (\ML) and \DL techniques are inherently data-dependent. Data validity is therefore the first condition for meaningful characterization of model performance. This dependency implies that 1) large datasets must be considered holistically alongside the \DL model to form an integrated system, and 2) these \AI methods require rigorous and systematic data curation. The data and \DL model must therefore be treated as an integrated system when interpreting model performance.

Neural networks learn from their training data. Therefore, to ensure reliable performance, the training data (simulated or measured \textit{in vivo}) must align closely with the data the model will encounter when deployed. This requirement distinguishes \DL and \ML models from traditional estimation techniques that optimize fitting parameters on a per-sample basis at run time. \DL models use features learned from the training data to characterize new samples. The degree to which these learned feature representations match the feature space of the unseen data directly impacts the accuracy and precision of the results. Accordingly, estimates of model performance are only meaningful to the extent that the training, validation, and test data support evaluation across the conditions in which the model is intended to operate.

\subsection{Synthetic Data in Model Validation}
Machine learning and deep learning models are biased estimators \cite{Mehrabi2019, Rosenberg2024}. Rigorous characterization therefore benefits from validation on simulated datasets with known ground truth. Although testing on \textit{in vivo} data is essential, it is insufficient to fully characterize model robustness and reliability. Synthetic datasets offer control over biases and parameter space coverage, making them particularly valuable when pathological diversity or class balance is limited. Accordingly, synthetic data should complement \textit{in vivo} datasets, especially prior to real-world deployment. This combined strategy improves generalization while preserving clinical relevance.

\hl{Theoretical }\MRS and \MRSI signals arise from well-characterized quantum mechanical spin dynamics and chemical exchange processes and are further shaped by hardware-specific acquisition parameters\hl{, sequence design, local micro- and macroscopic $B_0$ field inhomogeneities, and artifacts caused by non-ideal measurement conditions}. Because the signal generation process can be described by an explicit forward model grounded in physics, realistic synthetic datasets can be generated using established simulation frameworks. This enables systematic exploration of parameter spaces and controlled evaluation of model performance under diverse experimental conditions.

Physics-based simulation frameworks in \MRS offer several advantages. First, they enable a comprehensive evaluation of \DL model performance by uniformly and independently spanning the entire data parameter space, allowing models to identify independent features. Second, they may modulate (i.e., by inclusion with different degrees of freedom, or exclusion) natural correlations and implicit biases, such as those related to age, sex, or pathology, inherent in \textit{in vivo} data. This enables more reliable and flexible training and validation, as well as allowing models to learn features with progressive levels of independence.  Lastly, they enable stress testing under extreme but plausible conditions. Combined synthetic and \textit{in vivo} validation improves both robustness and clinical relevance. Synthetic data therefore provides a controlled basis for identifying performance boundaries and failure modes, while \textit{in vivo} data remain necessary to establish relevance to the target population and intended use.

\subsection{Data Curation and Annotation}\label{sec:5:subsec:data curation}
Robust \DL model performance depends on rigorous, comprehensive data curation and annotation, appropriate sampling distributions and class balance, and bias mitigation. Together, these factors support model reproducibility, reliability, and ultimately generalizability.

\subsubsection{Physiological and Pathological Data}
\MRS is a medical \MR modality is used to study a wide range of physiologies and pathologies. In \MRS, when not referring to pathological data, the term \textit{healthy controls} is commonly used. Because the meaning of \textit{healthy} can change based on age and sex, this work prefers the term \textit{physiological} which encompasses non-pathological biological variability across populations. In \MRS, the term \textit{pathological} therefore refers to data from tissue voxels with abnormal metabolism associated with \hl{neurological, psychiatric, oncological, metabolic, and systemic disorders, including both chronic and acute conditions}. Because of this variety, it is particularly important to clearly define the scope within which a model is characterized and expected to perform, together with the requirements that the training and testing data must satisfy to support that scope. The intended physiological and pathological scope therefore determines which sources of biological variability must be represented during model development and explicitly evaluated during testing.

The nature of metabolic profiling is strongly multivariate, being based on evaluating multiple metabolites, and these profiles change across different physiological and pathological conditions. This means model performance should be characterized for metabolites individually and jointly. Individual analysis identifies metabolite-specific failure modes, while joint characterization across the metabolite profile captures whether the model preserves the multivariate metabolic patterns associated with the physiological and pathological states represented in its intended scope. Joint characterization also provides a model-level representation that facilitates comparison between \DL fitting methods.

\subsubsection{Class Balance, Sampling Distributions, and Bias Mitigation}
Bias mitigation strategies, such as balanced sampling or data augmentation, ensure sufficient representation across the parameter space, improve model generalization, and reduce dependence on dominant data classes. Datasets with imbalanced class representations can result in models mischaracterizing underrepresented classes by interpolating features from dominant classes, which degrades performance \cite{Geirhos2020}. Consider, for instance, if a rare pathology like choroid plexus tumors is underrepresented in a tumor classification or segmentation dataset, a model may inaccurately classify these cases using features learned from more common pathologies such as gliomas. Synthetic data should be designed to ensure sufficient representation across all parameter ranges, especially for rare or sparsely sampled parameters. This balance is typically achieved by using data sampled from uniform distributions, or distributions that provide adequate coverage of the predefined parameter ranges, rather than distributions concentrated around their central values. Parameter dropout schemes during simulation can improve model generalization if the omissions are balanced and maintain the independence of parameter distributions. When a parameter is omitted during data simulation, associated spectral components must also be removed to prevent the introduction of erroneous features during training. If, for example, a metabolite is dropped, its lineshape and frequency shift parameters should also be dropped because the parent component is missing.

\subsubsection{Defining the Signal-to-Noise Ratio (SNR)}
\label{sec:dataModel:subsec:curation:subsubsec:SNR}
The performance of \DL models depends on data quality. Generally, when artifacts are excluded and biophysical properties (e.g., lipid contamination or macromolecular background signal) are comparably characterized across a dataset, \SNR is a major determinant of signal quality in \MRS. Characterizing performance as a function of \SNR is therefore necessary to determine the range of signal quality over which a model can be considered reliable. Noisy spectra or metabolite spectral components that are, according to a given measure, indistinguishable from the noise floor should be excluded from the datasets used in training and testing, to prevent \DL models from learning spurious features. Models trained on such filtered annotations perform more reliably and avoid degradation in performance due to biased feature representations. 

A clear definition of \SNR is therefore necessary to establish such an exclusion criterion. The terminology consensus paper \cite{Kreis2020} identifies several methods for calculating and referencing \SNR, each with its own benefits and drawbacks. A measure of global spectral \SNR (\gSNR) is useful but limited to determining the inclusion or exclusion of a spectrum as a whole. By contrast, a metabolite-specific \SNR (\mSNR) is much more informative, because it can be used to assess specific metabolic components within a given spectrum. The definition of \mSNR relies on the specific metabolite pattern, which is readily available either for simulated cases or after model fitting for \textit{in vivo} spectra. A quantity representing specific metabolite pattern "visibility" with respect to the noise floor is required. In the frequency-domain, spectral visibility can be represented by 1) the maximum value of a single spectral peak (either real or absolute component) or 2) the total area under the metabolite pattern curve (AUC). The first method, shown in Eqn.~\ref{eqn:spectral snr}, is the standard spectral \SNR definition. 
\begin{equation}\label{eqn:spectral snr}
    \textnormal{\textit{mSNR\textsubscript{S}}}(n) = \frac{1}{\sigma_{noise}}\cdot \max(\Re(x_n(\omega)))
\end{equation}
where $mSNR_S$ is the spectral \SNR of the given metabolite $n$, $\sigma_{noise}$ is the standard deviation of noise, $\mathcal{R}(x)$ represents the real component of a spectrum, $x$, and $x(\omega)$ is the broadened and scaled spectral basis function of metabolite $n$. The second method, shown in Eqn.~\ref{eqn:power snr} and referred to as the power \SNR, is the preferred definition because it is agnostic to shimming properties and accounts for metabolite \textit{J}-coupling patterns and proton counts by taking the integral of the signal. The definition of the noise floor is then identified as the standard deviation of pure noise, as extracted from the end of the \FID in the time-domain or a metabolite-free region in the frequency-domain.
\begin{equation}\label{eqn:power snr}
    \begin{aligned}
        \textit{mSNR}&\textsubscript{\textit{P}}(n) = \\&\sqrt{\frac{1}{\sigma_{noise}^2}\cdot \frac{1}{Ns}\sum_{n=1}^{Ns}\left|y\doteq A_n x_n(t) e^{-t(d_n-gt)} \right|^2}\\
         &= \sqrt{\frac{1}{\sigma_{noise}^2}\cdot \frac{1}{Ns}\sum_{n=1}^{Ns}\sqrt{\Re(y)^2 + \Im(y)^2}}
     \end{aligned}
\end{equation}
where $mSNR_P$ is the square root of the power \SNR of the given metabolite $n$, $y$ is the metabolite \FID, $A_n$ is the amplitude scaling factor, $x_n(t)$ is the \FID basis function of metabolite $n$, $d_n$ is the Lorentzian line broadening term, \textcolor{orange}{$g$} is the Gaussian line broadening term \hl{for when a Voigt lineshape is used}, $t$ is the time vector, \hl{$Ns$ is the number of data points in the }\FID \hl{or spectrum}, and $\mathcal{I}$ is the imaginary vector component of the data. The square root is used to bring the \SNR values closer to the numerical ranges of the current definitions by converting the metric to units of amplitude instead of true power. When calculating this in the frequency domain and the Fourier transform does not normalize the data with $\nicefrac{1}{\sqrt{Ns}}$, then $Ns$ should be squared, $Ns^2$, due to Parseval's theorem. For calculation in the time domain, the power can be calculated by taking the absolute value of the first \FID point. 

\subsubsection{Out-of-Distribution Errors}
\label{sec:dataModel:subsec:curation:subsubsec:OOD}
The deterioration in \DL performance due to feature underrepresentation often occurs at the tails, or edges, of the parameter space distribution, such as samples with either very high or very low metabolite concentrations, $C$. To mitigate these edge-limitations, parameter ranges in the training datasets are often expanded. Inflated parameter ranges, such as $[0.9 \times \text{C}_{min}, 1.1 \times \text{C}_{max}]$ \cite{LaMaster2024} or $[0.2 \times \text{C}_{max}]$ \cite{Rizzo2023}, can improve model robustness. Data augmentation techniques, such as parameter dropout or manipulation of sampling density at the tails of the parameter space, may also address these issues \cite{Rizzo2023}. However, care must be taken to avoid introducing class imbalances, which can degrade performance for other parameters. Careful reporting of model limitations in these regions is necessary for reliable characterization. These analyses help define the model's operating range by identifying conditions under which predictions become unreliable or should not be trusted.

\subsubsection{Annotation Accuracy and Ground Truth} 
\label{sec:dataModel:subsec:curation:subsubsec:Annotations}
Simulation parameters provide exact ground-truth labels by construction, whereas traditional model fitting results exhibit inconsistencies across software implementations \cite{Zollner2021, Marjanska2022}. The accuracy of the reference annotations therefore constrains the strength of any claim about model accuracy derived from them. When using \textit{in vivo} data without ground truth, semi-supervised or unsupervised \DL methods paired with robust signal models may provide alternatives.

For synthetic data, however, annotation accuracy also concerns whether formally exact ground-truth labels correspond to spectral components that remain recoverable under the simulated measurement conditions. When using synthetic training data, this work recommends a two-stage training and characterization procedure. In the first stage, the \DL model is trained using the original synthetic dataset and its unmodified ground-truth annotations. The errors should be plotted against the \mSNR values described above on an independent characterization set to determine whether a transition in prediction performance emerges as metabolite visibility decreases. Different estimators may retain useful predictive performance to different \mSNR levels, so this transition should be regarded as a model-specific effective recoverability threshold rather than a change in the physical noise floor of the acquisition. Below this threshold, the model no longer recovers the corresponding component with the prespecified level of accuracy or reliability.

Because \mSNR does not fully characterize intrinsic estimation difficulty, the stability of the recoverability threshold should also be examined across the information conditions characterized in Section~\ref{sec:Analy&Visual:subsec:infoStand}. The information-standardized analysis is used here to determine whether observations with similar \mSNR exhibit materially different recoverability because of factors such as spectral overlap, linewidth, or parameter covariance. It therefore provides a means of validating the adequacy of an \mSNR-based threshold rather than an additional filtering criterion. If the transition remains sufficiently stable across information conditions, a single \mSNR threshold provides a practical model-specific criterion. If substantial condition-dependent variation remains, the limitations of an \mSNR-only threshold should be reported and the operating range characterized using the relevant additional conditions.

Once the model-specific recoverability criterion has been established, it can be used to revise the annotations for the second training stage. Metabolite components that fall below the prespecified criterion can be removed from the explicit supervised targets, for example by setting the corresponding concentration target to zero, or reassigned to the nuisance or baseline representation when the training formulation permits this distinction. Metabolite-specific auxiliary parameters associated exclusively with a component that has been removed from the explicit target should be masked or zeroed consistently so that the network is not trained to infer parameters for a component designated as unrecoverable. Shared parameters, such as lineshape terms that jointly describe multiple spectral components, should not be removed solely because one metabolite falls below the recoverability criterion. When the corresponding modulated basis function is reassigned to the baseline or another nuisance component, its contribution remains represented in the simulated spectrum while the model is no longer required to attribute that contribution to an explicitly quantified metabolite. This preserves the spectral complexity presented to the model while aligning the supervised annotation with the component that the model is expected to recover.

The purpose of this procedure is to avoid explicitly training a model to recover parameters from spectral components for which its first-stage characterization provides insufficient evidence of reliable identifiability. Because this threshold is estimator-specific, it should be estimated separately for each model architecture or training strategy and fixed before the second training stage and final evaluation. The characterization data used to determine the threshold should remain independent of the final test set to prevent evaluation information from influencing the annotation procedure.

\subsection{Dataset Requirements for Model Characterization}
\label{sec:dataModel:subsec:data requirements}
In deep learning, datasets are commonly partitioned via randomly sampled into training, validation, and test sets. Typically, test sets constitute at least 5\% or 10\% of the total dataset, depending on the total dataset size, and serve a crucial purpose: they are used exclusively at the end to evaluate the model's performance on unseen data. For generalizability, the test set should represent the target population and operating conditions for which model performance is being claimed, which should closely approximate the training distribution. However, test datasets need not be confined to a subsample of the training dataset. Utilizing diverse and robust test sets is beneficial for thoroughly characterizing the model's performance and its limitations, as already discussed in Section~\ref{sec:5:subsec:data curation}. The design of the test set therefore determines which regions of the model's operating range can be evaluated and which claims about generalization are supported. This section will outline the size requirements and considerations for simulating or selecting data for test sets. 

\subsubsection{Dataset Splitting}
\label{sec:dataModel:subsec:data requirements:subsubsec:data splitting}
Random or na\"{i}ve splitting is often the default approach for partitioning the data into training, validation, and test splits. However, this can lead to an imbalanced and biased test set, which misrepresents the true data distribution, resulting in unreliable estimates of model performance and reduced generalizability. A common strategy to mitigate the limitations of a single random train-validation-test or train-test split is k-fold cross-validation (CV), where the data is repeatedly partitioned into $k$ training and validation subsets, providing more robust model estimates. When an independent test set is available, it should remain excluded from this process and be reserved for final model characterization. If not done properly, this strategy can introduce data leakage, which happens when information from a validation or test set becomes available during model training or selection, again leading to optimistically biased estimates of model performance \cite{Brookshire2024, Igwe2026}. 

Other strategies for creating balanced partitions involve optimizing split assignments based on dataset features, such as class distributions, sample sizes, and subject identities, as in multi-objective optimization (MOO) \cite{Igwe2026}. Generally, depending on the complexity of the data, at least one feature can be identified to guide more balanced split assignments, especially for longitudinal data. Subject-level splitting \cite{Rumala2023} and class-balanced splitting \cite{Xu2018, Khan2022} are examples of strategies used to create more representative and balanced test sets.

\subsubsection{Dataset Size Analysis}
\label{sec:dataModel:subsec:data requirements:subsubsec:dataset size analysis}
Sample size plays two distinct roles in deep learning. First, the training set must be sufficiently large to enable the model to learn the underlying mapping without overfitting or relying on spurious correlations. Second, the validation and test sets must be large enough to provide reliable and statistically meaningful estimates of model performance.

\hl{After partitioning the available data, the} training set is used to optimize model parameters, the validation set guides hyperparameter tuning and model selection, and the test set provides an unbiased estimate of generalization performance. The size of each split should therefore be chosen not only with respect to total dataset availability, but also to ensure stable training and statistically robust evaluation.

The size of the training dataset required to sufficiently train a \DL model depends on several factors including the number of weights in the neural network and the model architecture, the complexity of the learning task, and the complexity of the data. A practical way to estimate the necessary training sample size is to perform a dataset size analysis. Small datasets may fail to learn the underlying task adequately or may rely on spurious features. Using extremely large datasets is computationally expensive. This analysis will help approximate the minimum required number of training samples to successfully train a given network. 

This analysis involves training a model several times using different sizes of training datasets and then plotting the training and validation results to evaluate the curve of performance versus dataset size, \hl{as shown in Fig.}~\ref{fig:dataset size analysis}. The next portion of the curve is called the power-law region \cite{Muhlenstadt2024}. This region is where additional data provides additional information to the model and performance will increase. Past a certain point, this curve will saturate and form a plateau, known as the “irreducible error region” \cite{Muhlenstadt2024}, which can be seen on the right-hand side of Fig.~\ref{fig:dataset size analysis}. This work recommends a methodical approach: start with \hl{$10^4$} samples and increase the order of magnitude, i.e. the exponent, by 0.25 for every new trial. Hatami et al. \cite{Hatami2018}, in Fig.~\ref{fig:dataset size analysis}, used a coarse scale for the analysis, which works for some shallow or small model architectures, however, a more widely applicable range needs a somewhat finer scale, as proposed here. Most experiments will likely only need to use \hl{$10^4$ to $10^6$}, but models like transformers will need much more. 

After each training trial, \hl{the cumulative validation performance is plotted as a function of dataset size}. The held-out test set should remain untouched until the final evaluation. As the dataset grows, the curve will typically progress through the expected learning phases before approaching a plateau. Once performance has stabilized across several successive trials, the minimum sufficient training size can be estimated by identifying the transition point between the power-law improvement region and the saturation (irreducible error) region. In practice, this corresponds to selecting the dataset size closest to this transition, rounded to the nearest 10,000 samples.

\begin{figure}
    \centering
    \includegraphics[width=0.4\textwidth,keepaspectratio=True]{images/Hatami_01.png}
    \caption{Example of a dataset size analysis from  \cite{Hatami2018}. Beyond a variable threshold, the gains from additional data plateau.}
    \label{fig:dataset size analysis}
\end{figure}

\subsubsection{Validation and Test Dataset Size}
\label{sec:dataModel:subsec:data requirements:subsubsec:test set size}
For test-set analyses based on Pearson's $r$ and related metrics, this work recommends a minimum sample size of 500 independent samples. Armstrong et al. \cite{Armstrong2019} found that when presenting \hl{Pearson's $r$, and consequently any $r$-derived metric,} 500 samples are required to justify reporting 2 significant digits and $10^6$ samples are necessary for 3 significant digits. They recommended always reporting at least 2 digits, which this work also adopts. However, they emphasize the need to understand the limitations based on sample size. As mentioned above, larger sample sizes generally increase statistical power and may yield smaller \textit{p}-values for a given effect size, but the \textit{p}-value does not quantify confidence in the statistic itself. Therefore, this work also endorses using larger sample sizes, when available. For \textit{in vivo} data, when this minimum sample size is not available, the uncertainty and precision of the reported estimates should be characterized and reported appropriately. 

In small datasets, where minimum sample size requirements are not met, confidence intervals (CIs) can be estimated through bootstrapping, further discussed in Sec.~\ref{sec:perf:subsec:bootstrap}. Bootstrapping is a non-parametric resampling with replacement method that is used to estimate the sampling distribution of a statistic from the original data, enabling the calculation of CIs \cite{Haukoos2005, Good2005}. Similarly, permutation tests can be useful for estimating p-values in small samples where parametric assumptions may not be appropriate \cite{Good2005}. Resampling cannot compensate for nonindependence or an unrepresentative test set. The resampling scheme must also preserve any dependence structure present in the data.

\subsection{Sample Size Considerations}
\label{sec:Reg:subsec:comparison:subsubsec:samplesize}
\hl{Sample size requirements differ between the }\emph{in vivo} and \emph{in silico} evaluation regimes and must account for the stratified structure of the test sets. Accordingly, the nominal number of spectra should not be conflated with the number of statistically independent observations.

\subsubsection{In silico}
In the \textit{in silico} regime, simulation parameters are sampled to be \hl{independent and identically distributed (i.i.d.)} and each spectrum constitutes a fully independent observation, so standard sample size formulae for the Bland-Altman (BA) method apply directly~\cite{Lu2016}. The required sample size is governed by the desired width of the confidence interval on the limits of agreement and the expected standard deviation of differences. \hl{Lu et al.}~\cite{Lu2016} derive sample-size requirements for agreement analyses based on limits of agreement from the desired precision and power of the limits-of-agreement procedure, the expected distribution of pairwise differences, and the prespecified agreement limits. The required sample size therefore depends on the study design. More flexible formulations that allow bias or dispersion to vary across the measurement range may require larger samples because additional model parameters must be estimated reliably. The stratified evaluation structure defined in Section~\ref{sec:Reg:subsec:comparison:subsubsec:stratification} requires that the \hl{physiological and pathological} test sets each independently satisfy the minimum sample size criterion, so that stratified limits of agreement are estimated with comparable precision across both subsets. As proposed in Section~\ref{sec:dataModel:subsec:data requirements:subsubsec:test set size}, this work adopts a broader recommended minimum of at least 500 independent samples, particularly in the \textit{in silico} regime. 

\subsubsection{In vivo}
In the \textit{in vivo} regime, intra-subject correlation among spectra from the same subject reduces the effective sample size relative to the total spectrum count. The effective sample size $n_\textnormal{eff}$ must be estimated from the intraclass correlation (ICC) among intra-subject spectra before sample size adequacy can be assessed. Treating each spectrum as an independent observation in this setting inflates $n$ and produces confidence intervals on the limits of agreement that are artificially narrow. Consequently, as stated above, the nominal number of spectra should not be interpreted as the number of independent observations. Stratified \textit{in vivo} analysis further subdivides an already reduced effective sample, making adequate characterization of the pathological subset a practical constraint: the number of subjects with pathological concentrations must be sufficient to support per-stratum inference independently of the physiological set. This is a design constraint that should be addressed prospectively in the data collection strategy rather than accommodated post hoc by pooling strata.

\subsubsection{Correlated Data Samples}
\label{sec:5:subsec:correlated samples}
\textit{In vivo} \MRS data can experience two types of correlation that reduce the effective sample size. The first arises in \MRSI, where the truncated k-space sampling introduces a sinc-like point-spread function (PSF) determined by the k-space trajectory and reconstruction window. The side lobes of the PSF cause signal from spatially adjacent voxels to leak into each target voxel, correlating the fitting errors of neighboring voxels. This correlation is short-range and concentrated in the immediate spatial neighborhood of each voxel. The second source of correlation is physiological and applies to both \MRSI and SVS data. All spectra acquired from the same subject share a common metabolic state, \hl{system-level} disease burden, and scanner field conditions, making spectra from different brain regions within the same subject correlated regardless of their spatial separation. This intra-subject correlation is global and does not decay with voxel distance. As a result, the effective sample size in the \textit{in vivo} setting is bounded below by the number of subjects and above by the number of spectra, with the true value depending on the degree of intra-subject correlation for each metabolite. For the purposes of statistical inference, the number of subjects is treated as the conservative effective sample size. In the \textit{in silico} regime, provided that simulation parameters are sampled i.i.d., each spectrum constitutes a fully independent observation and the effective sample size is the number of simulated spectra, which is one reason the \textit{in silico} regime affords more reliable method comparison than the \textit{in vivo} regime. This statistical advantage does not replace \textit{in vivo} validation, which remains necessary to establish performance under real biological and acquisition variability. 

\subsubsection{Effective Sample Size Calculations}
\label{subsubsec:effective_sample_size}
When spectra from the same subject are correlated, the total number of spectra overstates the amount of independent information in the dataset. The effective sample size $N_{\textnormal{eff}}$ quantifies the equivalent number of independent observations and is derived from the intraclass correlation coefficient (ICC)~\cite{Shrout1979} and the design effect~\cite{Kish1965}\hl{, as outlined in the following steps}. This distinction is critical because the precision of statistical inference depends on the amount of independent information rather than on the nominal number of spectra.

\textbf{Step 1: Fit a one-way random effects model.} For a given metabolite, fit:
\begin{equation}\label{eqn:random_effects_icc}
    d_{s,n} = \mu + w_s + \nu_{s,n}
\end{equation}
\noindent \hl{where $d_{s,n}$ is the metabolite value for subject $s$ and spectrum $n$}, $w_s \sim \mathcal{N}(0, \sigma^2_{\textnormal{between}})$ is the subject-level random effect, described by a Gaussian distribution, and $\nu_{s,n} \sim \mathcal{N}(0, \sigma^2_{\textnormal{within}})$ is the residual. The variance components $\sigma^2_{\textnormal{between}}$ and $\sigma^2_{\textnormal{within}}$ are estimated by restricted maximum likelihood (REML).

\textbf{Step 2: Compute the ICC.} The ICC \hl{($\rho$)} is the proportion of total variance attributable to inter-subject differences~\cite{Shrout1979}:
\begin{equation}\label{eqn:icc}
    \rho = \frac{\sigma^2_{\textnormal{between}}}
                {\sigma^2_{\textnormal{between}} +
                 \sigma^2_{\textnormal{within}}}
\end{equation}

A value of $\rho = 0$ indicates that spectra from the same subject are uncorrelated, or no more similar than, spectra from different subjects. A value of $\rho = 1$ indicates that all spectra from the same subject are identical and the subject is the sole independent unit.

\textbf{Step 3: Compute the design effect.} The design effect quantifies the inflation in variance relative to a simple random sample of the same size~\cite{Kish1965}:
\begin{equation}\label{eqn:deff}
    \textnormal{deff} = 1 + \rho\,(\bar{n} - 1)
\end{equation}
\noindent where $\bar{n}$ is the mean number of spectra per subject. When cluster sizes are unequal, $\bar{n} = N_{\textnormal{total}} / S$, where $S$ is the number of subjects, which is the standard
approximation for unbalanced designs.

\textbf{Step 4: Compute the effective sample size.} The effective sample size is the total number of spectra divided by the design effect~\cite{Kish1965}:
\begin{equation}\label{eqn:neff}
    N_{\textnormal{eff}} = \frac{N_{\textnormal{total}}}{\textnormal{deff}}
                     = \frac{S \cdot \bar{n}}
                            {1 + \rho\,(\bar{n} - 1)}
\end{equation}

Two practical considerations apply. First, the ICC is metabolite-specific: metabolites with large inter-subject biological variability relative to intra-subject spectral variability will yield higher $\rho$ than metabolites whose concentration is spatially uniform across the brain. $N_{\textnormal{eff}}$ should therefore be computed separately for each metabolite and the range reported, rather than as a single aggregate value. Second, a commonly used rule of thumb is that reliable cluster-robust inference requires at least 30 independent clusters \cite{Cameron2008}. In this context, this applies to the number of subjects $S$, not to $N_{\textnormal{eff}}$, and should be verified against the subject count in the dataset before the metabolite-level inferential statistics are interpreted. Thus, a large number of spectra does not by itself establish that the dataset provides sufficient independent evidence for a particular statistical claim.

\textbf{Multivariate Effective Sample Size.} Several methods below allow for multivariate-level characterization across the joint set of metabolites. Such methods also have minimum sample size requirements, but this \hl{is} more complicated to calculate due \hl{to} the correlations between metabolites and intra-subject spectra. When the datasets are pooled across subjects and metabolites, there are two cross-classified clustering dimensions simultaneously: each observation $d_{i,n}$ belongs to metabolite cluster $i$ and to spectrum cluster $n$ at the same time, with neither clustering dimension nested within the other~\cite{Raudenbush2002}. The single-level Kish formula of Equation~\ref{eqn:deff} does not apply to this structure, and the design effect must account for both sources of clustering jointly. The cross-classified design effect is:
\begin{equation}\label{eqn:deff_cc}
    \textnormal{deff}_{\textnormal{CC}} =
        1 + \rho_i\,(\bar{K} - 1) + \rho_n\,(\bar{N} - 1)
\end{equation}
\noindent where $\rho_i$ is the ICC associated with the metabolite clustering dimension, $\rho_n$ is the ICC associated with the subject clustering dimension, $\bar{K}$ is the mean number of metabolites per spectrum, and $\bar{N}$ is the mean number of spectra per subject. The model-level effective sample size is then:
\begin{equation}\label{eqn:neff_cc}
    N_{\textnormal{eff}} = \frac{K \cdot S \cdot \bar{N}}
                            {\textnormal{deff}_{\textnormal{CC}}}
\end{equation}

Both ICCs are estimated from the fitted cross-classified random effects model by extracting the three variance components ($\hat{\tau}^2_{\textnormal{metabolite}}$, $\hat{\tau}^2_{\textnormal{subject}}$, and $\hat{\sigma}^2_{\textnormal{residual}}$) and computing:
\begin{subequations}\label{eqn:icc_cc}
    \begin{align}
        \rho_i &= \frac{\hat{\tau}^2_{\textnormal{metabolite}}}
                      {\hat{\tau}^2_{\textnormal{metabolite}} +
                       \hat{\tau}^2_{\textnormal{subject}} +
                       \hat{\sigma}^2_{\textnormal{residual}}},\label{eqn:icc_cc met}\\[2pt]
        \rho_n &= \frac{\hat{\tau}^2_{\textnormal{subject}}}
                      {\hat{\tau}^2_{\textnormal{metabolite}} +
                       \hat{\tau}^2_{\textnormal{subject}} +
                       \hat{\sigma}^2_{\textnormal{residual}}}\label{eqn:icc_cc subject}
    \end{align}
\end{subequations}

As stated above, the rule of thumb for the minimum number of clusters is at least 30 independent clusters~\cite{Cameron2008}. Because this setting has two cluster dimensions, this requirement applies separately to both the metabolite and subject dimensions. Here, the metabolite dimension refers to the number of metabolite-specific clusters in the fitted model for a given metabolite, not the number of distinct metabolites in the panel. This should be verified against both $K$ and $S$ before the multivariate-level inferential statistics are interpreted. For the metabolite dimension, the range of the number of per-metabolite clusters should be reported.

Conditioning on prespecified covariates that explain systematic variation in estimation difficulty, such as metabolite-specific \mSNR may explain part of the systematic variability between metabolite clusters. If these covariates explain meaningful variability associated with the metabolite clustering dimension, the remaining metabolite-level random-effect variance $\hat{\tau}_{\textnormal{metabolite}}^2$ can decrease after adjustment. Because the ICC is the proportion of total variance attributable to the metabolite clustering dimension, a reduction in $\hat{\tau}_{\textnormal{metabolite}}^2$ typically lowers $\rho_i$. This can reduce the corresponding design effect and increase the model-based effective sample size relative to an unadjusted analysis, although the magnitude of the change should be verified empirically rather than assumed \textit{a priori}.

\section{Statistical Considerations}
Once the data provides an appropriate basis for evaluation, the next question is whether the statistical analysis accurately characterizes model behavior and supports the claim being made. Deep learning model evaluation requires careful consideration of metrics that accurately reflect performance across diverse conditions. Because prediction and error distributions from \DL models may deviate from normality, their distributions should be examined explicitly, with robust or nonparametric approaches used when the assumptions underlying parametric inference are not satisfied.

Depending on the underlying ground-truth distribution, training data, model architecture, and optimization procedure, predicted variables and their errors may exhibit a wide range of distributional shapes. Well-learned variables should reproduce the relevant characteristics of their target distributions, whereas poorly learned variables may show distorted, compressed, or otherwise mismatched distributions. Rizzo et al.~\cite{Rizzo2023} \hl{show} examples in Fig.~\ref{fig:violin plots} of the variety of distribution shapes that networks learn. When these distributions violate the assumptions required by a particular parametric analysis, robust or nonparametric alternatives should be considered. Median-based summaries, for example, can provide more robust descriptions of central tendency under skewed or heavy-tailed distributions. A detailed analysis of such metrics is reported in \hl{Sec.}~\ref{sec: qual and quant}.

The dependency on data distribution in machine learning and deep learning applications requires model validation and calibration analyses to ensure accurate and reliable performance. Deep learning models incorporate inductive biases through their architectures, while their learned weights define the resulting feature representations \cite{Geirhos2020}. This is true for all \DL models, including quantification and denoising methods. Calibration and validation analyses are necessary to characterize how the model will perform when deployed. Validation establishes whether predictions remain sufficiently accurate across the intended operating conditions, while calibration assesses whether model-derived confidence or uncertainty is consistent with the observed predictive behavior. Section~\ref{sec:Reg} offers an overview of such methods.

The same comprehensive analyses used to characterize metabolite quantification performance should also be used to characterize the predictive performance for the remaining model fitting variables. Some may consider these additional parameters as auxiliary to the quantification parameters, which are the primary focus of model fitting. However, these variables will also impact metabolite quantification results and the ability to reconstruct fitted spectra. These analyses should, at a minimum, confirm that the model performs as expected. \hl{They can also provide a more complete understanding of the capacity and limitations of the proposed approach.}

When reporting model performance, this work strongly encourages researchers to switch from the commonly used box-and-whisker plots to violin plots. These plots are used to visualize distributions of cohort measurements or errors from different methodologies. The nature of deep learning prevents making assumptions about a variable’s distribution. Therefore, violin plots should be used to add information about the shape of the distribution to the box-and-whisker plot. As shown in Fig.~\ref{fig:violin plots}, violin plots make it much easier to visually inspect distributions, especially given the complexities of testing for differences between these shapes of dependent samples.

\subsection{Assumptions of Linearity and Nonlinearity}
\label{sec: assumptions}
Similar to conventional MRS quantification methods, deep learning quantification and \DL-based spectral modeling are inherently nonlinear inverse problems. Given a measured NMR signal (either in time or spectral domain), whose origin can be described by a physically informed and \textit{a priori} defined “forward model”, the objective is to estimate the input parameters of the fitting model, based on the observed data. The solution to this problem involves inverting the forward model, and thus constitutes an inverse problem. Traditional model fitting typically uses linear-combination models in which concentrations satisfy linearity assumptions, while only a few parameters do not (e.g., spectral lineshapes, offsets, artifacts, etc.). In contrast, \DL-based models introduce additional complexity due to the nonlinear activation functions used as filter operators within \DL architectures. This nonlinearity forms the cornerstone of the mathematical complexity underlying the success of \DL models and is what enables \DL networks to model complex systems effectively.

Despite the nonlinearity in \DL quantification and spectral modeling, for unconstrained continuous parameter estimation, the parameter estimation step in the final output layer will generally use a linear activation function, such as an affine function, or no activation function, i.e. the identity function. Other output parameterizations may be appropriate when explicit constraints on the estimated parameters are required. While the complexity of nonlinearities is beneficial in feature extraction and feature embedding layers, a linear final layer without them can provide a linear regression from the embedded feature space to the output space. Last-layer (\LL) approaches, such as Bayesian last layer (\BLL) \cite{Weber2018, Brosse2020, Watson2021, Fiedler2023}, are widely used in statistics-based \DL. These approaches treat all layers preceding the final layer as point-estimates, which simplifies network optimization and facilitates the use of standard linear regression metrics to evaluate model performance.

\subsection{Qualitative and Quantitative Assessments} %222->192
\label{sec: qual and quant}
While some works have primarily focused on predicting metabolite concentrations \cite{Hatami2018, Shamaei2021, Shamaei2023}, others have focused on more comprehensive deep spectral modeling \cite{Gurbani2019, Iqbal2021, Rizzo2023, LaMaster2024, Shamaei2025}. \DL-based spectral modeling offers a distinct advantage over metabolite-only predictions: it combines the speed of \DL models with the ability to evaluate the fit quality that comes from traditional fitting methods. The fit quality is generally judged by the fit residual, which describes the discrepancy between the reconstructed fit and the original data. This evaluation is essential for model development and characterization for (1) assessing the appropriateness of the predictions and the underlying signal model, (2) estimating errors, (3) supporting out-of-distribution (\OOD) error detection, (4) supporting model interpretability, and (5) calculating the fit quality number (\FQN), \hl{as recommended by the consensus papers }\cite{Near2021, Kreis2020}. Therefore, the design of \DL models should include the prediction of all \hl{model fitting} parameters (i.e. concentrations, lineshapes, phase and frequency offsets, etc.). Alternatively, \DL models' design could rely on the prediction of concentrations only as starting points to then initialize traditional fitting algorithms to estimate the remaining parameters necessary for signal reconstruction. Notably, evaluating model fits is essential even when uncertainty measures for estimated parameters are reported. Uncertainty and fit quality provide complementary evidence and should not be interpreted as interchangeable indicators of prediction reliability.

Nevertheless, it should be noted that \DL-based spectral modeling does not necessarily entail visual fit evaluation. While we recommend visual inspection for model development and characterization phases, it would be infeasible in practice for high-volume data, like high-resolution \MRSI. This work endorses \DL-based model fitting since the \FQN acts as a proxy to error estimation and out-of-distribution (\OOD) detection. Furthermore, robust quantitative metrics are essential for evaluating and comparing various models. In deep learning applications, large-scale datasets involving thousands to millions of spectra are standard. Visual inspection of model fit quality becomes impractical at this scale, and manual inspection across multiple metabolites and large numbers of spectra becomes increasingly impractical. These challenges highlight the necessity to develop reliable quantitative metrics to evaluate and compare models and their effects efficiently. Qualitative inspection can therefore reveal failure modes during development, while scalable quantitative analyses are required to establish how frequently and under what conditions those failures occur.

\subsection{The Problem with Parametric Metrics}
\label{sec:perf:subsec:transitioning}
In statistical analysis, many commonly used summary statistics and parametric inferential procedures rely, to different degrees, on assumptions about the underlying data or sampling distributions. Mean-based summaries and variance are sensitive to outliers and skewness, while parametric inference for statistics such as Pearson's $r$ relies on additional distributional assumptions. In contrast, nonparametric summaries, such as the median ($\widetilde{\mu}$ or $\widetilde{x}$) and interquartile range ($IQR$), offer robust alternatives that provide able less outlier-sensitive insights into the data's central tendency and variability under non-normal conditions. 

In deep learning, normality cannot be assumed and is not guaranteed. Bootstrapping and permutat\-ion testing are well-documented approaches to \hl{reduce reliance} on parametric assumptions about sampling distributions and are discussed below in Sec.~\ref{sec:perf:subsec:bootstrap}. Their primary drawback is that they are very time- and compute-intensive. In the context of deep learning, these characterizations are often performed during training after each epoch. It can be burdensome, however, to use these techniques after every epoch. To address this during training, this work proposes using median-based variants of the characterization metrics, which will be presented throughout the rest of this work. While median-based variants do not replace full inferential analysis, they can provide robust and computationally efficient summaries during the development phase. For the final characterization of trained models, the bootstrapping and permutation-based inference should be considered where appropriate to the statistic, study design, and dependence structure.

\subsection{Analytical versus Bootstrapping \& Permutation Implementations}\label{sec:perf:subsec:bootstrap}
Many of the evaluation metrics discussed throughout this work have both analytical and resampling-based implementations. Analytical methods directly calculate statistics such as confidence intervals and \textit{p}-values from equations under specified statistical assumptions, while bootstrap and permutation methods empirically estimate these quantities from resampled or null distributions, requiring fewer assumptions about the underlying data distribution \cite{Mooney1993, Good2005}.

\paragraph{Bootstrapping.}
Bootstrapping estimates the sampling distribution of a statistic by repeatedly resampling the observed data with replacement. Given a test set:
\begin{subequations}\label{eqn:bootstrap_samples}
\begin{equation}
    Z=\{z_1,z_2,\ldots,z_n\},
\end{equation}
a bootstrap sample
\begin{equation}
    Z^{*(k)}=\{z^{*(k)}_1,z^{*(k)}_2,\ldots,z^{*(k)}_n\}
\end{equation}
\end{subequations}
\noindent where $Z^{*(k)}$ denotes the $k^{\text{th}}$ bootstrap sample and is generated by sampling $N$ observations with replacement from the observed dataset, $Z$. For a statistic of interest $T$, such as \RMSE, \MAPE, Lin's concordance coefficient, or a regression coefficient, the bootstrap estimate is:
\begin{subequations}\label{eqn:bootstrap_distribution}
\begin{equation}
    T^{*(k)} = T(Z^{*(k)}),
\end{equation}
for $k=1,\ldots,B$ bootstrap replicates. The empirical distribution of
\begin{equation}
    \{T^{*(1)},T^{*(2)},\ldots,T^{*(B)}\}
\end{equation}
\end{subequations}
\noindent which approximates the sampling distribution of $T$. Confidence intervals can then be obtained from the empirical quantiles:
\begin{equation}
    CI_{95\%} = \left[ Q_{0.025}(T^*), Q_{0.975}(T^*) \right].
\end{equation}

\paragraph{Permutation Testing.}
Permutation methods are used primarily for hypothesis testing. The observed value of a test statistic $T_{\mathrm{obs}}$ can be written as:
\begin{subequations}\label{eqn:permutation_obs}
    \begin{equation}
        T_{\mathrm{obs}} = T(Z).
    \end{equation}
    Under the null hypothesis the observations or labels, $Z$ are randomly permuted to break the observed association:
    \begin{equation}
        T^{(k)}_{\mathrm{perm}} = T(\pi_k(Z)),
    \end{equation}
\end{subequations}
\noindent where $\pi_k$ denotes the $k^{\mathrm{th}}$ random permutation of $Z$. Repeating this procedure produces an empirical null distribution against which the observed statistic is compared. The resulting permutation \textit{p}-value is:
\begin{equation}\label{eqn:permutation_pval}
    p = \frac{ 1+ \sum_{k=1}^{P} I\!\left(T^{(k)}_{\mathrm{perm}} 
    \ge 
    T_{\mathrm{obs}} \right) } {P+1},
\end{equation}
\noindent where $I(\cdot)$ is the indicator function.

Bootstrap and permutation methods serve complementary purposes. Bootstrapping is primarily used to estimate confidence intervals and uncertainty bounds, whereas permutation methods are used for significance testing. Because both procedures derive empirical reference distributions from the observed data, they can reduce reliance on homoscedastic and parametric distributional assumptions in \DL-based model fitting. However, the resampling scheme must preserve relevant dependence structures and be appropriate for the statistic and inferential question being addressed.

\subsection{Interpretation of Statistical Evidence}
\label{sec:perf:subsec:awareness}
Statistical concepts and quantities that appear straightforward, such as \textit{p}-values, confidence intervals, normality assumptions, or measures of agreement, can be easily misinterpreted when they are applied without attention to the structure of the data and the intended use of the model. This section highlights several recurring statistical ambiguities and fallacies that are particularly relevant for the validation, comparison, and reporting of deep learning-based fitting models. The central consideration is not simply which statistical quantity is reported, but what property it measures and what claim that evidence can support.


\subsubsection{Statistical Significance: p-values and critical values}
\label{sec:perf:subsec:awareness:subsubsed:p-values}
There has been considerable discussion in recent years regarding the interpretation and use of \textit{p}-values and statistical significance. In 2019, \textit{The American Statistician} and \textit{Nature} published several editorials and commentaries addressing common misuses of statistical significance and encouraging more nuanced approaches to statistical inference \cite{Amrhein2019, Nature2019, Wasserstein2019}. The points most relevant to this work are summarized below. However, readers are strongly encouraged to consult the original articles for a more complete discussion.

A central issue is that not every evaluation question requires a test of statistical significance. In many model validation settings, the primary objective is not to decide whether an effect exists, but to estimate the magnitude, uncertainty, and practical relevance of model error. For example, when evaluating a fitting model, it may be more informative to report the expected error, confidence intervals, limits of agreement (LoA), calibration behavior, or uncertainty estimates than to report whether a metric differs significantly from zero or from another method. The decision to report a \textit{p}-value should therefore be justified by the inferential question being asked.

When \textit{p}-values are reported, they should not be interpreted as measures of effect size, clinical relevance, model quality, or the probability that the null hypothesis is true. A \textit{p}-value is a function of both the observed effect and the sample size. Consequently, very small effects can become statistically significant when the number of samples is sufficiently large, whereas practically important effects may fail to reach statistical significance in small or underpowered studies. Gomez-de-Mariscal et al. \cite{Gomez-de-Mariscal2021} demonstrated this sample-size dependence in detail and provided a method to account for this effect in the \textit{p}-value calculations. This issue is especially relevant for deep learning studies, where validation and test sets may contain thousands or millions of spectra, particularly when spectra are simulated. In such settings, \textit{p}-values can become extremely small even when the corresponding difference is negligible for practical or clinical purposes.

For this reason, \textit{p}-values should not be used as the primary evidence that a deep learning-based fitting model is accurate, reliable, or clinically useful. In large validation datasets, statistical significance may be almost guaranteed and therefore uninformative unless it is accompanied by effect sizes, uncertainty intervals, and a clear interpretation of practical relevance. Conversely, when the goal is formal binary decision-making, such as rejecting or failing to reject a null hypothesis under a predefined testing procedure, the decision should be based on the pre-specified significance level, the corresponding critical value, and the assumptions of the statistical test. In all cases, statistical significance should be interpreted alongside the magnitude of the observed effect, the uncertainty of the estimate, the dependence structure of the samples, and the clinical or methodological context.

\subsubsection{Weighing Statistical Significance, Effect Size, and Clinical Relevance}
\hl{A statistically significant difference is not necessarily a clinically meaningful difference. This distinction is particularly important in} \DL-based \MRS \hl{model fitting, where large synthetic or augmented datasets can make even negligible numerical differences appear statistically significant. Conversely, a result that is not statistically significant may still be clinically relevant if the uncertainty interval includes errors large enough to affect interpretation, diagnosis, treatment monitoring, or downstream decision-making.} \todo[inline]{still a todo. -- Marked as repetitive by Roland!}

For this reason, model evaluation should emphasize the magnitude and clinical interpretation of errors rather than relying solely on statistical significance. In quantitative \MRS, the practical importance of an error depends on the metabolite, field strength, acquisition protocol, spectral quality, expected biological range, and intended clinical or research application. For example, a small absolute error may be negligible for a high-concentration metabolite but unacceptable for a low-concentration metabolite or for a metabolite used to distinguish between clinically relevant groups. Similarly, two fitting methods may differ statistically while remaining interchangeable for a specific clinical task, or they may show no statistically significant difference while still having limits of agreement that are too wide for clinical use.

Confidence intervals are therefore often more informative than \textit{p}-values alone. A confidence interval provides a range of values compatible with the observed data under the assumptions of the analysis and allows the reader to assess both the precision of the estimate and whether the range is compatible with practical or clinical requirements. For example, confidence intervals can be reported for mean absolute error, bias, regression coefficients, limits of agreement, calibration slopes, or differences between methods. These intervals can be obtained analytically when assumptions are appropriate or through bootstrap procedures when analytical assumptions are difficult to justify \cite{Hazra2017}. In either case, the interval should be interpreted in relation to a pre-specified threshold of acceptable performance whenever such a threshold is available.

This distinction also motivates the use of equivalence or non-inferiority testing in method-comparison studies. Traditional null-hypothesis testing asks wheth\-er two methods are statistically different, but this is often not the most relevant question for clinical translation. A more clinically meaningful question may be whether the difference between methods is small enough to be considered acceptable, or whether a \DL-based method performs no worse than an established fitting method within a predefined tolerance. Such analyses require explicit definition of acceptable error margins, which should be justified using prior literature, biological variability, phantom experiments, expert consensus, or task-specific clinical requirements.

Accordingly, performance reporting should separate three related but distinct concepts: statistical detectability, numerical effect size, and practical or clinical relevance. Statistical detectability concerns whether an observed difference is unlikely under a specified null model. Effect size describes the magnitude of the difference or error. Clinical relevance determines whether that magnitude matters for the intended use case. Robust evaluation of \DL-based fitting models requires all three to be considered together. This distinction is essential when moving from evidence that two methods differ to evidence that their agreement is sufficient for a defined application.

\subsubsection{Distributions in Deep Learning}
Many commonly used statistical procedures rely, either explicitly or implicitly, on distributional assumptions. In model evaluation, these assumptions are often simplified to the question of whether the data are normally distributed. However, in deep learning-based model fitting, several distinct distributions must be considered, and they do not necessarily have the same properties or the same relevance for evaluation. The relevant distribution is therefore determined by the statistical question and the assumptions of the analysis, rather than by a general requirement that model outputs be normally distributed.

In the context of this work, at least three distributions are important. \hl{The first is the distribution of the sampled target values used to generate or select the data.} In synthetic datasets, this distribution is determined by the simulation design and may be uniform, Gaussian, empirically derived, deliberately imbalanced, or constrained by physiological plausibility. In \textit{in vivo} datasets, this distribution reflects the population, acquisition protocol, pathology, preprocessing, and inclusion criteria. The second distribution is the distribution of the model predictions. A model may reproduce the target distribution well, compress extreme values, introduce bias, or generate implausible predictions outside the range of expected metabolite concentrations. The third distribution is the distribution of residuals, defined as the difference between predicted and reference values. Residuals are central to many performance metrics and statistical analyses, but their behavior may vary across the parameter space, SNR, metabolite concentration, linewidth, spectral quality, scanner protocol, or pathology.

These distinctions are important because normality of one distribution does not imply normality of the others. For example, normally sampled target concentrations do not guarantee normally distributed predictions or residuals. Similarly, non-normal target values do not necessarily imply that residual-based methods are invalid if the residuals themselves are approximately symmetric, homoscedastic, and well behaved under the assumptions of the chosen analysis. In practice, residuals from deep learning models may be skewed, heavy-tailed, heteroscedastic, or structured by SNR and metabolite concentration. Such behavior is particularly likely near the limits of detection, in spectra with overlapping resonances, or in out-of-distribution cases.

Therefore, distributional assumptions should be evaluated at the level required by the analysis being performed. If a method assumes normally distributed residuals, then the residuals, rather than the sampled metabolite concentrations or predictions alone, should be examined. If parametric confidence intervals, regression models, or hypothesis tests are used, their assumptions should be stated and assessed. If those assumptions are not reasonable, nonparametric methods, robust statistics, bootstrap confidence intervals, transformation strategies, stratified analyses, or distribution-free agreement measures may be more appropriate.

For interpretability and uncertainty analysis, the residual distribution is especially important. A model whose residuals are approximately centered, symmetric, and stable across the relevant parameter space is easier to characterize than one whose errors vary systematically with concentration, \SNR, or other acquisition-dependent factors. However, Gaussian residuals should not be treated as a universal requirement for a model to be useful. Rather, the relevant question is whether the residual structure is understood, reported, and compatible with the intended use of the model. Distributional analysis therefore supports a claim about how reliably performance can be characterized, not a claim of model quality by itself.

\section{Characterization Metrics}\label{sec:eval metrics}
Current practice in \MRS model evaluation is characterized by inconsistently defined metrics and non-standardized sets of reported metabolites, \hl{as shown in }Section~\ref{sec:background:subsec:SOTF}, limiting reproducibility and cross-study comparison. This section identifies metrics that effectively characterize model-level and metabolite-level performance across four dimensions: precision, accuracy, regression quality, and distributional behavior. Each metric should be interpreted according to the specific property it measures and the claim that this evidence can support. The metrics below are intended to support comprehensive model-level assessment. Where overall performance is insufficient, metabolite-level evaluation can provide diagnostic insight into the sources of model failure. Metabolite-level comparison is non-trivial, however, as there are roughly 30 metabolites that can appear in a single \MRS spectrum depending on tissue type and pathology, making it essential to establish criteria by which metabolites can be meaningfully grouped and compared.

\subsection{Error Metrics}
\label{sec:Eval:subsec:errorMetrics}
Error analysis is organized around two principal components: accuracy and precision. \textit{Accuracy} describes the closeness of estimated values to the true values (i.e., ground truth) and is affected by systematic bias whereas \textit{precision} describes the consistency of predictions around that bias, and is also referred to as dispersion. In deep learning, several assumptions underpinning classical evaluation and inferential methods, such as the normality of residuals, independence of errors, and homoscedasticity, may be violated. Metrics or inferential procedures that rely on these assumptions can therefore provide misleading summaries or uncertainty estimates when the assumptions are not satisfied. The following sections examine the implications of these violations and present appropriate alternatives.

\subsubsection{Centroids}
\label{sec:Eval:subsec:error:subsubsection:centroids}
In many deep learning applications, the normality assumption of the prediction errors may not be satisfied. This work proposes modifying calculations that use the mean ($\mu$ in Eqn.~\ref{eqn:mean}) to instead use the median ($\widetilde{\mu}$ in Eqn.~\ref{eqn:median}) when a robust summary of central tendency is required. Both metrics describe the centroid of a distribution, but the median is less sensitive to extreme outliers, which are common with \DL models. For approximately symmetric distributions, particularly those close to normality, the mean and median will generally be similar.
\begin{equation}\label{eqn:mean}
    \mu = \frac{1}{N}\sum_{n=1}^N x_n
\end{equation}
\begin{equation}\label{eqn:median}
    \widetilde{\mu} = \textnormal{med}(x) = \begin{cases}
        x_{\nicefrac{n+1}{2}} & \textnormal{if $n$ is odd}\\
        \frac{1}{2} \left(x_{\nicefrac{n}{2}}+x_{\nicefrac{n}{2}+1}\right) & \textnormal{if $n$ is even}
    \end{cases}
\end{equation}
\noindent where $x$ is the ordered sample of size $N$, $x_n$ denotes the $n$-th order statistic, $\mu$ is the arithmetic mean, and $\widetilde{\mu}$ is the median-based alternative adopted throughout this work. Note that the piecewise definition assumes the data are sorted in ascending order.

\subsubsection{Dispersion}\label{sec:Eval:subsec:error:subsubsection:dispersion}
\paragraph{Variance ($\sigma^2$).}
Variance (Eqn.~\ref{eqn:variance}) measures the average squared deviation of data points from the mean, penalizing deviations equally in both directions. As a mean-based estimator, it is susceptible to outliers, consistent with the limitations discussed in Section~\ref{sec:perf:subsec:transitioning}. A robust dispersion alternative (Eqn.~\ref{eqn:np variance}) is proposed, replacing both the mean centroid and the expectation operator with their median-based counterparts. This formulation is closely related to the Median Absolute Deviation (MAD) \cite{Rousseeuw1993}, a well-established robust dispersion estimator, and inherits much of its resistance to extreme values. The use of median-based dispersion propagates robustness through downstream metrics by ensuring that the underlying dispersion estimates are not distorted by outliers.
\begin{equation}\label{eqn:variance}
    \sigma^2 = \frac{1}{N-1} \sum_{n=1}^N \left( \mu - x_n \right)^2 
\end{equation}
\begin{equation}\label{eqn:np variance}
    \widetilde{\sigma}^2 = \textnormal{med} \left( (\widetilde{\mu} - x_n)^2  \right) 
\end{equation}
\noindent where $N$ is the sample size, $x_n$ is the $n$-th observation, $\mu$ is the arithmetic mean, $\widetilde{\mu}$ is the median, $\sigma^2$ is the standard sample variance with Bessel's correction, and $\widetilde{\sigma}^2$ is the median-based robust alternative. The nonparametric standard deviation $\widetilde{\sigma} = \sqrt{\widetilde{\sigma}^2}$ is used in subsequent metrics.

\paragraph{Interquartile Range (\textit{IQR}).}
The interquartile range ($IQR$) is a robust measure of dispersion that represents the range containing the central $50\%$ of the data distribution. Equation~\ref{eqn:iqr} shows that it is calculated from the sorted data as the difference between the third ($Q3$, 75th percentile) and first ($Q1$, 25th percentile) quartiles. %
\begin{subequations}\label{eqn:IQR}
    \begin{align}
        IQR &= Q_3 - Q_1 \label{eqn:iqr} \\[1ex]
        \text{Outlier Range} &:= [\ Q_1 - 1.5 \cdot IQR, \label{eqn:iqr_outlier} \\
        &\phantom{:= [}\ \ Q_3 + 1.5 \cdot IQR\ ] \nonumber
    \end{align}
\end{subequations}

Eqn.~\ref{eqn:iqr_outlier} defines Tukey's fence \cite{Tukey1977}, a standard criterion for identifying outliers. Because the $IQR$ is computed from the central portion of the distribution, it is relatively unaffected by extreme values compared with mean-based measures of dispersion. Unlike $RMSE$, which amplifies the influence of large errors through squaring and is therefore sensitive to outliers regardless of distributional shape, the $IQR$ provides a robust summary of variability for skewed distributions. This makes it particularly appropriate for characterizing \DL model outputs in model fitting, where outlying predictions are not uncommon.

\paragraph{Coefficient of Variation (\textit{CV}).}
The coefficient of variation is a scale-invariant, dimensionless measure of relative dispersion, defined as the ratio of a dispersion measure to a centroid measure. Traditionally, it is defined as the standard deviation ($\sigma$) divided by the mean ($\mu$). Its unit-free nature allows direct comparison across metabolites with different concentration ranges and units. The standard and robust formulations are:
\begin{subequations}
    \begin{align}
        &CV = \frac{\sigma}{\mu} \label{eqn:cv} \\ &\widetilde{CV} = \frac{\widetilde{\sigma}}{\widetilde{\mu}} \label{eqn:np cv}
    \end{align}
\end{subequations}
\noindent where $\mu$ and $\sigma$ are the mean and standard deviation, and $\widetilde{\mu}$ and $\widetilde{\sigma}$ are their median-based counterparts as defined above. The robust formulation 
$\widetilde{CV}$ may be preferred for \DL applications where distributions and residuals tend to be non-normal or outliers are present. Both formulations are undefined when the centroid is zero and become unreliable when it is small relative to the dispersion. In \MRS, this is a practical consideration for metabolites with concentrations near zero or with low signal-to-noise.

\subsubsection{Errors}\label{sec:Eval:subsec:error:subsubsection:errors}
\paragraph{Mean Percent Error ($MPE$).}
The mean percent error ($MPE$) quantifies the bias of a set of predictions as the mean signed relative error:%
\begin{equation}\label{eqn:MPE}
    MPE = \frac{100}{N}\sum_{n=1}^N \frac{x_n-y_n}{y_n}
\end{equation}
\noindent where $N$ is the total number of samples, $x_n$ are the predicted values, and $y_n$ are the ground-truth values. Under this sign convention, $MPE < 0$ indicates systematic underestimation and $MPE > 0$ indicates systematic overestimation. Values closer to zero indicate lower bias. Both $MPE$ and its nonparametric analogue $MePE$ (Eqn.~\ref{eqn:MePE}) are undefined when $y_n = 0$. In \MRS this is a practical consideration for metabolites with concentrations at or near zero, so zero or near-zero reference values require explicit handling. If a stabilizing constant such as $\epsilon=1$\textsc{E}$^{-6}$ is added to the denominator, its choice and its effect on the reported error should be justified. For non-normal or skewed error distributions, $MePE$ provides a more robust estimate of the central tendency of the error distribution, and for near-normally distributed errors it converges to $MPE$.
\begin{equation}\label{eqn:MePE}
    MePE = 100 \cdot \textnormal{med}\left( \frac{x_n - y_n}{y_n} \right)_{n=1}^N
\end{equation}

\paragraph{Mean Absolute Percent Error, ($MAPE$).}
The mean absolute percent error ($MAPE$) measures the average magnitude of relative errors, and is therefore a joint reflection of both bias and dispersion. A model with $MPE = 0\%$ can still exhibit large $MAPE$ if individual errors are large but cancel. Conversely, a model with nonzero $MPE$ will have its $MAPE$ inflated by the systematic bias even if precision is high. To isolate dispersion from bias, the percent errors can be centered by subtracting the estimated bias before taking absolute values:
\begin{subequations}\label{eqn:MAPE}
\begin{align}
    MAPE &= \frac{100}{N} \sum_{n=1}^N \frac{|x_n-y_n|}{y_n} \label{eqn:MAPE basic}\\ &\Rightarrow \frac{100}{N} \sum_{n=1}^N \left( \frac{|x_n-y_n|}{y_n} - \frac{MPE}{100} \right) \label{eqn:MAPE w/correction}
\end{align}
\end{subequations}
\begin{equation}\label{eqn:MeAPE}
\begin{aligned}
    MeAPE &= \frac{100}{N}\ \textnormal{med}\left( \frac{|x_n-y_n|}{y_n} \right) \\ &\Rightarrow \frac{100}{N}\ \textnormal{med} \left( \frac{|x_n-y_n|}{y} - \frac{MePE}{100} \right)
\end{aligned}
\end{equation}
\noindent Equation.~\ref{eqn:MAPE w/correction} computes the mean absolute deviation of percent errors around the estimated bias $MPE$, isolating precision from systematic offset. The nonparametric analogue $MeAPE$ applies the same correction using $MePE$ as the bias estimate, which is more appropriate when the error distribution is skewed. Both $MAPE$ and $MeAPE$ are in the range $[0, +\infty)$, with smaller values indicating greater precision. Their percentage basis and scale independence make them suitable for cross-metabolite and cross-dataset comparison, unlike scale-dependent metrics such as $RMSE$.


\paragraph{Root Mean Squared Error, (\textit{RMSE}).}
The root mean squared error ($RMSE$) is the quadratic mean of the differences between predicted and ground truth values:
\begin{equation}\label{eqn:rmse}
    RMSE = \sqrt{\frac{1}{N}\sum_{n=1}^N\left(x_n - y_n \right)^2} %\approx \sqrt{\textnormal{med}\left(\left(y - x \right)^2\right)}
\end{equation}
\begin{equation}\label{eqn:rmese}
    RMeSE = \sqrt{\textnormal{med}\left(\left(x_n - y_n \right)^2\right)}
\end{equation}
\noindent where $N$ is the sample size, $x_n$ are the predicted values, and $y_n$ are the ground truth values. $RMeSE$ (Root Median Squared Error) is a notation adopted in this work for the median-based analogue. $RMSE$ is scale-dependent and therefore not suitable for comparison across datasets or metabolites with different concentration ranges \cite{Hyndman2006}. As a mean-based metric, it penalizes large errors disproportionately due to the squaring operation, and is therefore sensitive to outliers regardless of distributional shape.

$RMSE$ is equivalent to the standard deviation of prediction errors only when the estimator is unbiased (i.e., the mean prediction error is zero). This equivalence motivates its use as a proxy for comparing \DL model performance against the Cram\'{e}r--Rao Lower Bound (CRLB) reported by traditional fitting algorithms, though this comparison is only valid to the extent that the \DL model is unbiased. Furthermore, $RMSE$ alone does not characterize asymmetry, tail behavior, or other distributional structure in non-Gaussian error distributions. For these reasons, $RMSE$ should not be used as the sole evaluation metric. $RMeSE$ is proposed here as a more robust alternative for non-normal error distributions, and percentage-based metrics such as $MAPE$ and $MeAPE$ are preferred where scale independence is required. Collectively, these error metrics characterize the magnitude, bias, and dispersion of prediction errors. They do not by themselves establish agreement with an existing fitting method or whether the observed errors are acceptable for the intended clinical or research use.

\subsection{Regression Metrics}
\label{sec:Eval:subsec:RegMetrics}
Capturing the statistical distributions of estimated parameters is essential. However, it is equally important to evaluate the linearity between model outputs and known reference values, independently of the broader complexity of the data. Particularly in preliminary \textit{in silico} settings, where ground truth values are readily available, regression analysis serves as a critical tool for evaluating quantification algorithms with clear, quantitative assessments of model fit and predictive accuracy. These metrics facilitate direct comparisons across network architectures, helping to identify models yielding more accurate predictions in controlled environments, and complement broader development strategies such as architecture search and data augmentation. However, regression metrics primarily characterize association and systematic response behavior and should not be interpreted as evidence of agreement or interchangeability.

Despite their utility, common regression metrics have limitations, such as relying on distributional or error-structure assumptions for particular inferential procedures and requiring minimum sample sizes. The metrics presented below are most informative for per-metabolite univariate analyses because na\"{i}ve multivariate summaries risk masking heterogeneous performance. These constraints limit their applicability for comprehensive model validation. Therefore, they should be supplemented with additional methods to ensure robust validation. The following sections present relevant linear regression metrics, along with their parametric and nonparametric formulations and considerations for using them.

\subsubsection{Pearson's Correlation Coefficient, \textit{r}}
\label{sec:Eval:subsec:RegMetrics:subsubsec:pearson}
The Pearson correlation coefficient, $r$, \cite{Pearson1895} quantifies the strength and direction of the linear relationship between two continuous variables. For descriptive purposes, $r$ can be computed without a normality assumption. However, classical parametric inference for $r$, including the conventional test of zero correlation and Fisher-transformed confidence intervals, is derived under bivariate normality. In sufficiently large samples, asymptotic or resampling procedures can support inference under weaker distributional assumptions, but large sample size alone does not guarantee robustness to substantial skewness or heavy tails. The inferential procedure, e.g. regular inference or bootstrapping, should therefore be selected according to the relevant distributional assumptions rather than assuming that non-normality is resolved by sample size alone. Furthermore, Pearson's $r$ summarizes the overall linear association between two variables and may therefore obscure heteroscedastic behavior. In such settings, strong overall correlation does not necessarily imply consistent performance throughout the parameter space, as poor predictions in some regions may be masked by better performance elsewhere.

Consistent with the discussion in Sec.~\ref{sec:perf:subsec:transitioning}, a robust median-based alternative, $\widetilde{r}$, is also defined below for use when outliers or distributional skew are a concern. Comparing the parametric and nonparametric versions will highlight the impact of failures and extreme outlier predictions. The two versions of the coefficient are defined as:
\begin{subequations}\label{eqn:r}
    \begin{align}
        r &= \frac{\sum(x_i - \bar{x})(y_i - \bar{y})}{\sqrt{\sum(x_i - \bar{x})^2\sum(y_i - \bar{y})^2}}\label{eqn:mean r}\\[2pt]  \widetilde{r} &= \frac{\sum(x_i - \tilde{x})(y_i - \tilde{y})}{\sqrt{\sum(x_i - \tilde{x})^2\sum(y_i - \tilde{y})^2}}\label{eqn:median r}
    \end{align}
\end{subequations}
\noindent where $x$ is the set of predicted values, $y$ is the corresponding set of ground truth values, $\bar{x}$ and $\bar{y}$ are their respective means, and $\tilde{x}$ and $\tilde{y}$ are their respective medians. Both coefficients are bounded to the range $[-1, 1]$, where $r = 1$ indicates a perfect positive linear relationship, $r = -1$ a perfect negative linear relationship, and $r = 0$ the absence of a linear relationship, though the latter does not preclude a nonlinear association. This bounded range facilitates comparison across datasets and models. In the multivariate setting of model fitting, both coefficients should \hl{be} computed and interpreted on a per-metabolite basis, where they provide a direct characterization of model performance for each quantity of interest. The per-metabolite coefficients can then be averaged into an aggregate metric.

\subsubsection{Coefficient of Determination, \textit{$r^2$}}\label{sec:Eval:subsec:RegMetrics:subsubsec:r2}
The coefficient of determination, $r^2$, quantifies the proportion of variance in the observed outcomes that is explained by the model. It is defined through the decomposition of total variance into explained and unexplained components, specifically, $SS_{tot} = SS_{reg} + SS_{res}$, and is given by:
\begin{equation}\label{eqn:r2}
\begin{aligned}
    r^2 &= 1 - \frac{SS_{res}}{SS_{tot}}
    \\ \textnormal{ given that } SS_{res} &\approx \sum_i(y_i - x_i)^2 
    \\ \textnormal{ and } SS_{tot} &= \sum_i (y_i - \Bar{y})
\end{aligned}
\end{equation}
\noindent where $x$ denotes the fitted values, $y$ the observed values, $\bar{y}$ the mean of the observed values, $SS_{res}$ the residual sum of squares, and $SS_{tot}$ the total sum of squares. Under standard conditions, i.e., an ordinary least-squares linear model fit with an intercept to the same data, $r^2 \in [0, 1]$ where $r^2 = 0$ indicates that the model explains none of the variance and $r^2 = 1$ indicates a perfect fit. When the same expression is applied directly to fixed predictions from a \DL model, however, $r^2$ can be negative if the predictions perform worse than the mean of the observed values. In simple linear regression with an intercept, $r^2$ is equivalent to the square of the Pearson correlation coefficient $r$, directly linking the strength of the linear association to the proportion of variance explained. This equivalence does not extend to multiple regression, where $r^2$ reflects the joint contribution of all predictors.

While $r^2$ provides an interpretable and widely used measure of goodness-of-fit, two limitations are worth noting in this context. First, $r^2$ does not penalize model complexity, i.e., it will increase, or at worst remain unchanged, as predictors are added, regardless of whether those predictors are meaningful. This motivates the use of adjusted $r^2$, discussed below. Second, a high $r^2$ does not establish that residuals are homoscedastic or free of systematic structure, which are fundamental assumptions of $r^2$. Nonlinear behavior, heteroscedasticity, or localized prediction failures can therefore remain hidden despite a high overall $r^2$, similar to Pearson's $r$. When used for univariate analysis, even from a multivariate regression, $r^2$ serves as a useful complement to $r$: where $r$ describes the direction and strength of the linear association, $r^2$ gives a direct, variance-based account of model adequacy.

\subsubsection{Adjusted Coefficient of Determination, \ensuremath{\bar{r}^2}}
\label{sec:Eval:subsec:RegMetrics:subsubsec:adj r2}
By construction, $r^2$ is non-decreasing as additional predictors are added to a model, making it an optimistic measure of fit when comparing models of different complexity. To counteract this, Wherry \cite{Wherry1931} proposed an adjusted coefficient of determination, $\bar{r}^2$, which penalizes model complexity through the degrees of freedom:
\begin{equation}\label{eqn:adjusted r2}
    \Bar{r}^2 = 1 - \frac{SS_{res}/df_{res}}{SS_{tot}/df_{tot}} = 1 - (1 - r^2)\frac{n-1}{n-p-1}
\end{equation}
\noindent where $n$ is the sample size, $p$ is the number of predictors, $df_{res} = n - p - 1$ and $df_{tot} = n - 1$ are the residual and total degrees of freedom, respectively. By construction, $\bar{r}^2 \leq r^2$, and $\bar{r}^2$ can become negative when the model fits worse than a horizontal mean line. This metric is well-suited to model selection in traditional linear combination fitting methods \cite{Olejnik2000,Pirenne2024}, where the number of predictors is small and well-defined.

For \DL models, however, $\bar{r}^2$ has limited utility as a model characterization metric. The formula assumes $p$ reflects the number of predictors in a linear model. Importantly, the number of trainable parameters in the \DL model does not define $p$. Instead, $p$ refers to the number of predictors, or input variables, in the regression used to evaluate the model outputs. For example, in a per-metabolite regression of reference values on \DL predictions, $p=1$ regardless of the number of trainable parameters in the network. Consequently, $\bar{r}^2$ adjusts for the complexity of the auxiliary regression rather than the complexity of the \DL model that generated the predictions. When the same regression structure is used across \DL models, $\bar{r}^2$ therefore provides little additional information beyond $r^2$ and should not be interpreted as a complexity-adjusted measure for comparing \DL architectures.

This distinction is particularly relevant for \DL-based spectral modeling. Explicitly predicting nuisance and fitting parameters, such as phase, frequency offsets, or lineshape parameters, may provide useful auxiliary supervision by requiring the learned representation to account for sources of spectral variation that would otherwise need to be implicitly ignored when predicting metabolite concentrations alone. Whether this improves quantification, however, is a property of the learned multi-output representation and should be established through controlled comparisons of models trained with different sets of output variables. In this setting, $\bar{r}^2$ may still be used for an auxiliary regression containing multiple predicted quantities as predictors, but it then characterizes the complexity and explanatory performance of that regression rather than the benefit of jointly modeling those quantities within the \DL network.

\subsubsection{Concise Score, $r^3$}\label{sec:Eval:subsec:RegMetrics:subsubsec:concise}
The concise score, $r^3$, was proposed by Rizzo et al.~\cite{Rizzo2023} to synthesize multiple regression metrics into a single scalar, with the original formulation defined as:
\begin{equation}\label{eqn: original concise}
    concise = a \cdot r^2 
\end{equation}
\noindent where $a$ is the slope of the linear regression of predicted values on ground truth values and $r^2$ is the coefficient of determination. A limitation of this formulation is that $a$ is unbounded, as it depends on the scale and units of both variables. This makes the original $concise$ score difficult to interpret in absolute terms and unreliable for comparison across metabolites, datasets, and models.

To address this, $a$ is replaced by the Pearson correlation coefficient $r$, which is the standardized slope of the regression line, related to $a$ by $r = a \cdot s_y/s_x$, where $s_x$ and $s_y$ are the standard deviations of the predicted and observed values respectively. This substitution carries the same directional information as $a$ while removing its scale dependence. The proposed metric is then:
\begin{equation}\label{eqn: new concise}
    r^3 = r \cdot \max(0, r^2) 
\end{equation}

The resulting score is denoted $r^3$ and is defined here as the product of the independently calculated Pearson correlation coefficient, $r$, and coefficient of determination, $r^2$. The notation should therefore not be interpreted as implying that $r^2$ is necessarily calculated as the algebraic square of $r$. Both $r$ and $r^2$ should be calculated directly from the predictions and ground truths. When the independently calculated coefficient of determination satisfies $0 \leq r^2 \leq 1$, the concise score is bounded to $[-1,1]$, as shown in Eqn.~\ref{eqn:concise breakdown}. If $r^2<0$, which is possible when $r^2$ is calculated directly from fixed model predictions, this bound no longer applies and the negative $r^2$ itself indicates performance worse than the mean-reference baseline. This is why Eqn.~\ref{eqn: new concise} bounds $r^2$ from below at zero. This bounded range is a deliberate property of the proposed formulation: it ensures $r^3$ is directly interpretable and comparable across metabolites, datasets, and models in a way the original $concise$ score is not.

The resulting metric jointly reflects the strength and direction of the linear association and the goodness of predictive fit captured by $r^2$. The sign of $r^3$ follows that of $r$. In model fitting, a negative $r^3$ would indicate an inverse relationship between predictions and ground truth, signaling a fundamental failure of the model. Because $r$ and $r^2$ are calculated separately, they can capture complementary aspects of prediction behavior. For example, a model may exhibit a strong linear association while still showing poor predictive agreement with the reference values due to systematic scaling or offset errors. 
As a composite metric, $r^3$ rewards models that simultaneously achieve high correlation and high variance explanation, and as with $r$ and $r^2$, it is most appropriately computed and interpreted on a per-metabolite basis in the multivariate setting of model fitting. As with $r$ and $r^2$ individually, heteroscedasticity in the data can bias the calculated metric.

One inherent limitation of any composite metric is information compression: a given value of $r^3$ can arise from different combinations of $r$ and $r^2$, and reporting $r^3$ alone sacrifices the diagnostic granularity provided by its components individually. It is therefore recommended that $r^3$ be reported alongside $r$ and $r^2$ rather than as a replacement for them, especially given the bounding constraint in Eqn.~\ref{eqn: new concise}.

\subsubsection{Theil-Sen Estimator}
\label{sec:Eval:subsec:RegMetrics:subsubsec:theil-sen}
The Theil-Sen estimator \cite{Theil1950} provides a robust, nonparametric approach for estimating the slope and intercept of a linear relationship between two variables. While $r$ and $r^2$ characterize the strength and explained variance of a linear association, neither provides a direct estimate of the regression slope or intercept. Consequently, a model can achieve $r \approx 1$ and $r^2 \approx 1$ while exhibiting a slope that deviates substantially from unity or a nonzero intercept, reflecting proportional or fixed bias that both metrics leave entirely undetected. The Theil-Sen estimator addresses this gap directly, providing explicit estimates of slope and intercept that reveal the nature and magnitude of systematic bias across the concentration range. Unlike ordinary least squares regression, which minimizes squared residuals and is sensitive to outliers, the Theil-Sen method estimates the slope as the median of all pairwise slopes between data points:
\begin{equation}\label{eqn:theil-sen}
    \begin{gathered}
        y = \hat{m}x + \hat{b}\text{, where}\\[1pt]
        \hat{m} = \text{median}\left(\left\{ \frac{y_j-y_i}{x_j-x_i}:1\leq i < j \leq n \right\}\right)\\[1pt]%\text{ where } x_i \neq x_j\\[1pt]
        \hat{b} = \text{median}\left(\left\{ y_i-\hat{m}x: 1\leq i \leq n \right\}\right)
    \end{gathered}
\end{equation}
\noindent where $\hat{m}$ is the slope computed from the set of pairwise ratios $(y_j - y_i)/(x_j - x_i)$ for $1 \leq i < j \leq n$ when $x_i\neq x_j$, and $\hat{b}$ is the intercept that is subsequently obtained as the median of the residual terms $y_i - \hat{m}x_i$. Here the slope quantifies proportional bias, or scaling across the concentration range, while the intercept quantifies the fixed offset. 

This median-based formulation confers robustness to outliers and skewed distributions, making the estimator well suited for robust estimation in the presence of heteroscedasticity, heavy-tailed errors, non-normal distributions, or outliers. In particular, extreme prediction errors, which can disproportionately influence least-squares estimates, have a limited effect on the Theil-Sen slope due to the use of medians. This property is consistent with the use of nonparametric correlation measures and facilitates a more reliable characterization of central trends in the presence of atypical observations.

Despite these advantages, the Theil-Sen estimator has limitations. The method assumes an underlying monotonic linear relationship and does not account for more complex nonlinear dependencies. Furthermore, its computational cost scales quadratically with the number of samples due to the pairwise slope calculations, which may become restrictive for large datasets. In practice, the estimator is most informative when used alongside parametric approaches, where discrepancies between estimates can indicate sensitivity to outliers or violations of model assumptions. For metabolite-level analyses, the Theil-Sen estimator provides a complementary, robust measure of linear association that is less affected by localized prediction failures.

\renewcommand{\arraystretch}{1.2}
\begin{table*}[T!]
    \centering
    \begin{tabular}{|l|c||c|c|c|c|c|}\hline
         & & \multicolumn{5}{c|}{Strength of Linear Association} \\\hline
        Context & Metric & Negligible & Weak & Moderate & Strong & 
        Very Strong \\\hline \hline
        \multirow{3}{*}{General \cite{Schober2018}}
        & $|r|$    & [0.00, 0.10) & [0.10, 0.40) & [0.40, 0.70) & 
                     [0.70, 0.90) & [0.90, 1.00] \\
        & $r^2$    & [0.00, 0.010) & [0.010, 0.16) & [0.16, 0.49) & 
                     [0.49, 0.81) & [0.81, 1.00] \\
        & $|r^3|$  & [0.00, 0.001) & [0.001, 0.064) & [0.064, 0.34) & 
                     [0.34, 0.73) & [0.73, 1.00] \\\hline
        \multirow{3}{*}{MRS Quantification}
        & $|r|$    & [0.00, 0.10) & [0.10, 0.50) & [0.50, 0.80) & 
                     [0.80, 0.95) & [0.95, 1.00] \\
        & $r^2$    & [0.00, 0.010) & [0.010, 0.25) & [0.25, 0.64) & 
                     [0.64, 0.90) & [0.90, 1.00] \\
        & $|r^3|$  & [0.00, 0.001) & [0.001, 0.125) & [0.125, 0.51) & 
                     [0.51, 0.86) & [0.86, 1.00] \\\hline
    \end{tabular}
    \caption{Strength of linear association for regression metrics. 
    General thresholds are adapted from Schober et al.~\cite{Schober2018}, with $r^2$ and $|r^3|$ boundaries derived by squaring and cubing the $|r|$ thresholds respectively. MRS quantification thresholds reflect the higher standard appropriate for algorithm comparison and repeatability assessment in this context.}
    \label{tab:correlation values}
\end{table*}
\renewcommand{\arraystretch}{1.0} % reset after
\subsubsection{Stratification References}
\label{sec:Eval:subsec:RegMetrics:subsubsec:stratification}
Table~\ref{tab:correlation values} provides stratifications of $r$, $r^2$, and $r^3$ derived from the thresholds proposed by Schober et al.~\cite{Schober2018} for $|r|$, with the $r^2$ and $|r^3|$ boundaries obtained by squaring and cubing the $|r|$ thresholds respectively. The transformed boundaries are intended as heuristic reference values that preserve an interpretable correspondence with the established $|r|$ scale, rather than as independently validated cutoffs for $r^2$ or $r^3$. These thresholds, originally adapted from Cohen's effect size conventions \cite{Cohen1988}, are widely used as a general reference but are explicitly context-dependent and should not be treated as universal standards.

In the \MRS context, the appropriate thresholds depend on the task. For methods that modify spectral content prior to fitting, such as denoising, lower correlations are expected and the general Schober thresholds provide a reasonable reference. For quantification tasks, whether comparing algorithms or assessing repeatability, a higher standard is both expected and required: an $r$ value of $0.50$ would not constitute meaningful association in this setting. The recommended  thresholds for \MRS quantification contexts are provided in the lower rows of Table~\ref{tab:correlation values}. Regardless of context, thresholds should ideally be defined prospectively for a given task rather than selected post hoc to suit observed results.

\subsection{Model Benchmarking Metrics}
\label{sec:Eval:subsec:benchmarking}
Effective comparison of models requires global metrics capable of capturing overall performance across all multivariate predicted outputs (e.g., metabolite concentrations, fitting parameters). Sections~\ref{sec:Eval:subsec:errorMetrics} and \ref{sec:Eval:subsec:RegMetrics} introduced component-level metrics. This section presents model-level metrics to enable systematic model benchmarking: 1) Lin's concordance correlation coefficient \cite{Lin1989} ($\rho_C$), which should be reported for point estimate and predictive models; 2) the negative log predictive density ($NLPD$), which applies to predictive models with uncertainty estimates (see Section~\ref{sec:Interpretability}); and 3) error consistency, which captures the degree to which prediction errors are reproducible across models. These metrics should be reported over the full set of metabolites to facilitate direct comparison across models. When values are aggregated across metabolites with different scales or concentration ranges, appropriate standardization or an explicitly defined metabolite-level aggregation procedure is required to prevent individual metabolites from disproportionately determining the model-level result. Computing and visualizing them per metabolite and as a function of metabolite-level \SNR provides additional diagnostic insight into model shortcomings and failure modes.

\subsubsection{Lin’s Concordance Correlation Coefficient, \textnormal{$\rho$\textsubscript{C}}} \label{sec:Eval:subsec:RegMetrics:subsubsec:Lin's rho}
Lin's $\rho_C$ \cite{Lin1989} was designed to assess the agreement between two methods measuring the same quantity. This is distinct from correlation: two methods can be highly correlated yet systematically disagree due to differences in mean or scale. $\rho_C$ addresses both simultaneously and in this work is used to compare the outputs of different fitting protocols applied to the same data.
\begin{subequations}
    \begin{align}
        \rho_C = r \cdot C_b &= \frac{r \cdot 2\sigma_x\sigma_y}{\left(\mu_x-\mu_y\right)^2 + \sigma_x^2 + \sigma_y^2} \label{eqn:lin's rho}\\ 
        \widetilde{\rho_C} &= \frac{r \cdot 2\widetilde{\sigma_x}\widetilde{\sigma_y}}{\left(\widetilde{\mu}_x-\widetilde{\mu}_y\right)^2 + \widetilde{\sigma_x}^2 + \widetilde{\sigma_y}^2}\label{eqn:lin's rho np}
    \end{align}
\end{subequations}
\noindent where $r$ is Pearson's $r$, $C_b$ is the distributional concordance coefficient (also called the bias correction factor), $\mu_x$, $\mu_y$ are the means, and $\sigma^2_x$, $\sigma^2_y$ are the variances computed without Bessel's correction. The metric decomposes as $\rho_C = r \cdot C_b$, separating precision from accuracy: $r$ captures the strength of the linear relationship between the two methods, while $C_b \in [0,1]$ penalizes deviation from the line of perfect agreement (the 45° line through the origin). When the means and standard deviations of the two samples are equal, i.e. $C_b = 1$, $\rho_C$ reduces to Pearson's $r$. When they differ, $r$ is scaled down accordingly. This decomposition directly mirrors the accuracy-precision framework introduced in Section~\ref{sec:Eval:subsec:errorMetrics} and makes $\rho_C$ more diagnostically informative than $r$ alone. The difference between $\rho_C$ and $r$ quantifies the additional loss of concordance attributable to departures in location and scale from the line of perfect agreement. Consistent with Lin's formulation, the variances are computed using the $1/N$ normalization rather than Bessel's correction.

Lin's original analysis \cite{Lin1989} demonstrated robustness of the estimator across normal, uniform, and Poisson distributions. However, as discussed in Section~\ref{sec:perf:subsec:transitioning}, the distributional form cannot be assumed \textit{a priori} for \DL model outputs, and mean- and variance-based estimators are non-robust in the presence of outliers or heavy-tailed distributions. Consistent with the framework applied throughout this work, a nonparametric analogue $\widetilde{\rho}_C$ is proposed by substituting $\widetilde{\mu}$ for $\mu$ and $\widetilde{\sigma}^2$ for $\sigma^2$ in Eqn.~\ref{eqn:lin's rho} (see Sections~\ref{sec:Eval:subsec:error:subsubsection:centroids} and \ref{sec:Eval:subsec:error:subsubsection:dispersion}). This modification has not been formally validated in the literature to the same extent as Lin's original formulation and should be interpreted accordingly.

Table~\ref{tab:SoA rho_c values} provides stratification thresholds for $\rho_C$ according to McBride \cite{McBride2005}, for both continuous and discrete variables. These thresholds apply equally to $\widetilde{\rho}_C$ as a general reference, subject to the same contextual caveats discussed in Section~\ref{sec:Eval:subsec:RegMetrics:subsubsec:stratification}.
\begin{table*}[ht]
    \centering
    \begin{tabular}{|c||c|c|c|c|c|}\hline
        & \multicolumn{5}{c|}{Strength-of-Agreement} \\\hline
        & Negligible & Weak & Moderate & Strong & Very Strong \\\hline \hline
        Continuous & [0.00, 0.10) & [0.10, 0.90) & [0.90, 0.95) & [0.95, 0.99) & [0.99, 1.00] \\
        Discrete & [0.00, 0.10) & [0.10, 0.65) & [0.65, 0.80) & [0.80, 0.90) & [0.90, 1.00] \\
        \hline 
    \end{tabular}
    \caption{Strength of agreement for Lin's $\rho_C$, adapted from McBride \cite{McBride2005}, for continuous and discrete variables.}
    \label{tab:SoA rho_c values}
\end{table*}

\subsubsection{Negative Log Predictive Density, \textnormal{NLPD}}
\label{sec:Eval:subsec:RegMetrics:subsubsec:NLPD}
The negative log predictive density (\textit{NLPD}) evaluates a model's probabilistic predictive performance by jointly considering predictive accuracy and precision through an a proper probabilistic scoring rule that rewards both accurate predictions and appropriately scaled predictive uncertainty. Prediction errors are scaled by the predicted variance, while the logarithmic variance term penalizes excessively broad predictive distributions. This metric quantifies how well a probabilistic model's predictive densities fit the observed data and, for a Gaussian predictive distribution, is given by:
\begin{equation}\label{eqn:NLPD}
    \begin{aligned}
        NLPD = \frac{1}{N}&\sum_{n=1}^N \Bigg( \frac{1}{2} \Bigg[ \log(\sigma_{pred,n}^2) + \\
         & \frac{(y_{true,n}-y_{pred,n})^2}{\sigma_{pred,n}^2} + \log(2\pi) \Bigg] \Bigg)
    \end{aligned}
\end{equation}
\noindent where $N$ is the sample size, $\sigma^2_{pred,n}$ is the predicted variance for sample $n$, and $y_{true,n}$, $y_{pred,n}$ are the ground truth and predicted values respectively. This equation assumes that the data have been standardized or otherwise scaled to be dimensionless, ensuring that the argument of $\log(\sigma_{pred,n}^2)$ is dimensionless. For univariate evaluation in fixed units, rescaling shifts all \textit{NLPD} values by the same additive constant and therefore does not affect relative model comparisons. Standardization is particularly important when \textit{NLPD} values are compared or aggregated across output variables with different units or scales. This equation should be evaluated using dimensionless standardized quantities so that both the variance term and the log-variance term ($\log(\sigma_{pred,n}^2)$) are dimensionless. For each metabolite, the ground-truth and predicted values should be transformed using the same fixed location and scale parameters, and the predicted variance should be divided by the square of the corresponding scale factor. These standardization parameters should be estimated from the training or development data and subsequently held fixed during validation and testing. When \textit{NLPD} values are aggregated across metabolites, standardization should be performed separately for each metabolite before the metabolite-specific scores are combined.

\textit{NLPD} can be computed per output variable or summarized across the metabolite set. When only metabolite-specific marginal predictive distributions are available, each metabolite should first be evaluated separately and a model-level summary may then be obtained as the equal-weighted mean of the metabolite-specific \textit{NLPD} values. For this aggregate, fixed per-metabolite standardization derived from the training or development data can be used to place the marginal scores on comparable scales, with the predicted variances transformed by the corresponding squared scale factors. This quantity should be described as an aggregate marginal or macro-averaged \textit{NLPD}, rather than as a joint multivariate predictive score. If the model instead provides a full predictive covariance matrix across metabolites, a joint multivariate \textit{NLPD} can be calculated directly from the corresponding multivariate predictive density.

\textit{NLPD} should be reported alongside $\rho_C$ when evaluating probabilistic models or models with reported uncertainties. Where $\rho_C$ captures the accuracy and distributional concordance of point predictions, \textit{NLPD} scores the predictive mean and variance jointly and is sensitive to inappropriate uncertainty estimates. Together they provide complementary assessments of predictive accuracy and probabilistic performance. However, \textit{NLPD} is not a calibration diagnostic by itself, so uncertainty calibration should be assessed separately using dedicated calibration analyses. Two limitations are worth noting. First, the Gaussian formulation used here assumes a Gaussian predictive distribution, which may not hold for all \DL models. Second, it penalizes overconfident predictions severely, since a small $\sigma^2_{pred}$ combined with a large error produces a disproportionately large loss.

\subsubsection{Error Consistency}\label{sec:Eval:subsec:RegMetrics:subsubsec:error con}
While Lin's $\rho_C$ and \textit{NLPD} are valuable model-level metrics for assessing overall performance, they do not capture pairwise, sample-level differences between models. This limitation is particularly relevant when comparing updated versions of a model or evaluating architectural robustness. Ensuring that improvements arise from correcting prior errors rather than introducing new ones is critical. This concept, referred to as \textbf{churn} by Milani Fard et al.~\cite{MilaniFard2016}, quantifies prediction disagreements between models.

For example, consider a classifier achieving $80\%$ accuracy that is updated to $90\%$. In the best case, the additional $10\%$ of correct predictions comes entirely from the previously incorrect $20\%$. In the worst case, the model corrects all $20\%$ of prior errors but introduces new errors on $10\%$ of previously correct predictions, resulting in $30\%$ churn, i.e. predictions that changed between models, excluding consistently incorrect predictions.

As a second example, consider five identical networks, i.e., networks with the same architecture and training configuration, trained with different random seeds and each achieving $85\%$ accuracy. A robust model and training protocol would be expected to produce consistent predictions across runs. If each network correctly predicts a different subset of the data, the apparent accuracy is misleading and the conclusion of robustness would be unjustified.

Rahman et al.~\cite{Rahman2022} proposed error consistency metrics ($\rho_{EC}$) to quantify sample-wise difference in residual errors between models. Of those proposed variations, this work recommends: 
\hl{equation}
\begin{equation}\label{eqn:error consistency}
    \begin{gathered}
        \rho_{EC} = \frac{1}{N}\sum_{n=1}^N \textnormal{sign}(r_{i,n}\cdot r_{j,n}) \cdot \frac{||r_{i,n}|-|r_{j,n}||}{||r_{i,n}|+|r_{j,n}||} \\
        \textnormal{with\ } f(x) = \hat{y} \approx y %\\
            \ \textnormal{and} \ 
            r = \hat{y} - y
    \end{gathered}
\end{equation}
\noindent where $y_n$ is the ground truth, $\hat{y}_{i,n}$ and $\hat{y}_{j,n}$ are the predictions of models $i$ and $j$ for sample $n$, $r_{i,n}$ and $r_{j,n}$ are their respective residuals, and $N$ is the sample size. $\rho_{EC} \in [-1, 1]$. At the sample level, contributions near zero indicate similar residual magnitudes, irrespective of residual sign. Positive contributions indicate residuals with the same sign, whereas negative contributions indicate residuals with opposite signs. The magnitude of each contribution increases as the relative difference between the two residual magnitudes increases. Because $\rho_{EC}$ averages these signed contributions, values near zero can reflect either similar residual magnitudes across models or cancellation between positive and negative sample-wise contributions and should therefore not be interpreted as evidence of high error consistency in isolation. This metric is particularly useful for studying variability introduced by random seeding, weight initialization, or hyperparameter changes. Evaluating $\rho_{EC}$ across training runs or ensemble members can help identify whether changes in aggregate performance are accompanied by systematic changes in the models' sample-wise residual behavior rather than merely a redistribution of errors. Considering churn together with $\rho_{EC}$ helps distinguish genuine improvements in predictive capability from a redistribution of errors that produces similar aggregate performance.

\subsection{Quantifying Distributional Similarity}
\label{sec:Eval:subsec:DistSim}
A complementary approach to regression-based evaluation is to assess whether there is statistical evidence that the predicted distribution of metabolite quantities differs from the ground truth distribution. Hypothesis tests provide a principled framework for this, evaluating three primary characteristics: location, scale, and shape. \textit{Location} refers to the centroid and is described by the mean and median. \textit{Scale} refers to the dispersion of the distribution and is measured by, for example, the standard deviation or IQR. The \textit{shape} of the distribution refers to higher order features, such as skewness and kurtosis, that are beyond what is captured by location and scale alone.

To interpret these tests properly, it is essential to identify what each test evaluates. After calculating the test statistic, it is compared against a \textit{critical value}, i.e. a threshold determined by the sample size and chosen confidence level. If the test statistic does not exceed the critical value, the null hypothesis, that the two samples are drawn from the same distribution, is not rejected. Importantly, failing to reject the null does not confirm that the distributions are identical. It indicates that the evidence against the null is insufficient at the chosen confidence level. This distinction is particularly relevant for small sample sizes, where low statistical power may obscure genuine differences.

The binary nature of this decision provides a straightforward criterion for indicating evidence of distributional differences within a defined scope. However, these tests do not capture the linearity between ground truth and predictions, which is addressed by the regression metrics presented in Section~\ref{sec:Eval:subsec:RegMetrics}. Furthermore, strong association or agreement metrics, such as $\rho_C$, are consistent with distributional similarity but do not guarantee it, since two samples can be highly correlated while differing in their marginal distributions. Hypothesis tests therefore serve as a supplemental tool, particularly useful for identifying whether specific distributional characteristics, like location, scale, or shape, contribute to model failures that aggregate metrics do not expose.

The natural formulation is a nonparametric two-sample test for dependent (paired) data, comparing predicted values directly against ground truth. However, nonparametric tests designed for dependent data with no distributional assumptions are scarce. This work proposes reframing the problem as a comparison of independent samples: model outputs on the test set are compared against ground truth values from a sampled subset of the training data, provided both are generated under the same i.i.d. simulation parameters. Under this condition, the two samples share the same parent distribution by construction, and finite-sample variability is explicitly accounted for by the null distribution of the two-sample test. A statistically significant distributional discrepancy can therefore be attributed to model error at the chosen significance level rather than to differences in the underlying data-generating distribution. 

The second required condition is that the neural network model must be sufficiently trained. In supervised learning, model optimization explicitly minimizes the discrepancy between predictions and the corresponding training labels, and sufficiently expressive \DL models commonly fit, or memorize, their training data to very low error. Therefore, once convergence and low training error have been confirmed from the training curves of the loss function and evaluation metrics, the model should closely reproduce the ground truth values for the sampled training subset. This provides the mechanism that justifies using those ground truth values as an independent reference distribution. This reframing then enables the use of a broader class of nonparametric two-sample tests. Because this approach assumes i.i.d. simulation conditions, its validity in settings where training and test data are not drawn from the same generative distribution should be carefully assessed.

\subsubsection{Hypothesis Tests for \DL-Based Studies}
\label{sec:Eval:subsec:DistSim:subsubsec:hypothesis tests}
The following tests are nonparametric procedures for independent samples, making them applicable to \DL-based studies without requiring a specified parametric distribution. While asy\-mptotic approximations are often considered reasonable for moderate sample sizes (commonly $n \geq 30$), larger samples improve the accuracy and robustness of statistical inference \cite{Conover1999,Hollander2015}. This work strongly recommends $n \geq 500$ based on the increased robustness this affords in practice for the sample sizes typical of \DL evaluation sets \cite{Armstrong2019}.

K\"{o}ssler and Mukherjee \cite{Kossler2019} introduced the concept of \textit{versatile alternatives}, which are tests that jointly evaluate location, scale, and shape. They proposed the ECWAS statistic (Eqn.~\ref{eqn:ecwas}) as a composite metric based on the squared Euclidean distance between standardized versions of: the Wilcoxon-Mann-Whitney U test (WI) for location, the Ansari-Bradley test (AB) for scale, and the AntiSavage test (AS) for shape.
\begin{equation}\label{eqn:ecwas}
    \begin{aligned}
        T_{EC}& _{WAS} = \\&\left(\frac{WI-\mu_{WI}}{\sigma_{WI}}\right)^2\ +\ \left(\frac{AB-\mu_{AB}}{\sigma_{AB}}\right)^2+\ \\[0pt] 
        &\phantom{\left(\frac{WI-\mu.}{\sigma}\right.}\left(\frac{AS-\mu_{AS}}{\sigma_{AS}}\right)^2
    \end{aligned}
\end{equation}
\noindent where $WI$, $AB$, and $AS$ are the Wilcoxon-Mann-Whitney, Ansari-Bradley, and AntiSavage test statistics, evaluating location, scale, and shape respectively; $\mu_{WI}$, $\mu_{AB}$, and $\mu_{AS}$ are their respective means under the null hypothesis, derived analytically from the null distributions as functions of sample size; and $\sigma_{WI}$, $\sigma_{AB}$, and $\sigma_{AS}$ are their respective standard deviations under the null hypothesis. Each term is a standardized statistic, and $T_{ECWAS}$ is their squared Euclidean distance, combining evidence against the null across all three distributional characteristics into a single composite test statistic. Each component can be evaluated individually, enabling targeted diagnosis of whether distributional divergence arises from differences in location, scale, or shape, an important capability when characterizing specific failure modes of \DL models.

While composite metrics like Euclidean-type Combined Wilcoxon, Ansari-Bradley, and Savage-type statistic (ECWAS) leverage the strengths of their components, they also inherit their limitations. In this work, the primary limitation of the AntiSavage statistic is its sensitivity to distributional skewness: it is optimized for right-skewed distributions and may underperform on left-skewed or symmetric distributions. Funato et al.~\cite{Funato2024} addressed this by proposing a max-type variant that returns the maximum of the AntiSavage statistic $AS$ and the Savage statistic $S$, the latter being suited to left-skewed Gumbel distributions. This modification is recommended when the skewness of \DL model outputs cannot be assumed in advance, which is the typical situation in model fitting. Eqn.~\ref{eqn:ecwas} is then reformulated as follows:
\begin{equation}\label{eqn:funato}
    \begin{aligned}
        T_{EC}& _{WAS} = \\&\left(\frac{WI-\mu_{WI}}{\sigma_{WI}}\right)^2\ +\ 
        \left(\frac{AB-\mu_{AB}}{\sigma_{AB}}\right)^2+\ \\[1pt]
        &\phantom{\left(\frac{WI}{\sigma}\right.}
        \left(\frac{\max(AS,S)-\mu_{\max}}{\sigma_{\max}}\right)^2
    \end{aligned}
\end{equation}

\subsection{Extension to Multivariate Analyses}
\label{sec:Eval:subsec:multivariate}
The metrics defined above are most informative on a per-metabolite basis. Nevertheless, extending these analyses to the joint metabolite set provides complementary, model-level information that makes it easier to compare multiple methods. Na\"{\i}ve pooling of the data across metabolites could introduce Simpson's paradox \cite{Blyth1972, Simpson1951}, in which an association present within each subgroup is obscured or reversed upon aggregation, when inter-metabolite variance dominates the pooled statistic. The transformations below address this directly, each encoding distinct information about model behavior across the joint set. 

\subsubsection{Error Analysis}
\label{sec:Eval:subsec:multivariate:subsubsec:error}
Absolute prediction error should first be evaluated as a function of metabolite-level \mSNR. For each metabolite, absolute error can be plotted against both spectral and power-based \mSNR. The two definitions provide complementary perspectives: spectral \mSNR characterizes error relative to metabolite visibility in the spectrum, whereas power-based \mSNR characterizes error relative to the metabolite's contribution to the measured signal. Until their relative utility has been established empirically, both should be considered. These relationships may be visualized separately for each metabolite or jointly across the metabolite set. 

Because these analyses treat metabolites independently, aggregate error should be calculated in a metabolite-balanced manner. The error metric is first averaged within each metabolite and the resulting metabolite-specific values are then averaged across metabolites. This prevents metabolites with larger numbers of valid estimates, larger signal amplitudes, or prominent singlet resonances from dominating the aggregate statistic. It also reduces the influence of unequal effective sample sizes that can arise when weakly visible metabolite signals are not reliably distinguished from the fitted baseline. The resulting summary facilitates comparison between models, although it should be interpreted alongside the metabolite-specific error distributions rather than as a complete characterization of performance.

More detailed analyses may examine absolute error directly as a function of individual factors such as concentration, linewidth, or frequency offset. Because the relevance and range of these variables can differ substantially among metabolites, this level of analysis is most appropriately performed on a metabolite-specific basis rather than by pooling all metabolites into a common response curve.

\subsubsection{Regression Analysis}
\label{sec:Eval:subsec:multivariate:subsubsec:regression}
To isolate intra-metabolite predictive linearity in a pooled setting and suppress the inter-metabolite cluster structure responsible for Simpson's paradox, the reference and predicted values for each metabolite are standardized to zero mean and unit variance prior to pooling:
\begin{subequations}\label{eqn:standardized_pooled}
\begin{align}
    y^\prime_m = \frac{y_m - \mu_m}{\sigma_m}, \qquad \hat{y}^\prime_m = \frac{\hat{y}_m - \mu_m}{\sigma_m}\label{eqn:standardize normal data}\\
    y^\prime_m = \frac{y_m - \tilde{\mu}_m}{{IQR}_m}, \qquad \hat{y}^\prime_m = \frac{\hat{y}_m - \tilde{\mu}_m}{{IQR}_m}\label{eqn:standardize nonnormal data}
\end{align}
\end{subequations}
where $\mu_m$ and $\sigma_m$ are the mean and standard deviation of the reference values for metabolite $m$ in Eqn.~\ref{eqn:standardize normal data} for standardizing normal data. For nonnormal data, Eqn.~\ref{eqn:standardize nonnormal data} uses the median $\tilde{\mu}_m$ and the interquartile range ${IQR}_m$ instead of the mean and standard deviation. Transforming all metabolites to a common origin and scale before aggregation removes the inter-metabolite location and scale differences that would otherwise dominate the pooled correlation while preserving deviations of the predictions from their corresponding reference distributions. The resulting Pearson $r$ summarizes the pooled intra-metabolite correlation after suppressing inter-metabolite differences in location and scale and serves as a compact scalar summary for comparing models across a large metabolite set. It encodes no information beyond what the per-metabolite values already provide, but offers a convenient aggregate figure of merit.

\subsubsection{Agreement Analysis}
\label{sec:Eval:subsec:multivariate:subsubsec:agreement}
Lin's $\rho_C$ evaluates univariate agreement. To extend this concept to the joint metabolite set, the multivariate concordance correlation coefficient \cite{Choi2025} is defined for paired vector observations. Let $\boldsymbol{x}$ be the $d \times p$ matrix of fitting parameters, with each row in $d$ containing the results from one fitting method $i$. 
The multivariate concordance correlation coefficient is given by:
\begin{equation}\label{eqn:multivariate rho_c}
\begin{aligned}    
    \rho_C^{(M)}& = \\ &\frac{2\sum\limits_{i=1}^{d-1}\sum\limits_{j=i+1}^d \text{tr}\left[\text{Cov}(\boldsymbol{x}_i,\boldsymbol{x}_j) \right]}{(d-1)\sum\limits_{i=1}^d \text{tr}[\text{Var}(\boldsymbol{x}_i)]+\sum\limits_{i=1}^{d-1}\sum\limits_{j=i+1}^d||\boldsymbol{\mu}_i-\boldsymbol{\mu}_j||^2}
\end{aligned}
\end{equation}
\noindent where $d$ is the total number of fitting methods $i$, indexed into unique pairs by $i$ and $j$, $\text{tr}$ indicates the trace, $\text{Var}$ and $\text{Cov}$ represent the variance-covariance and the cross-covariance matrices, and $\boldsymbol{\mu}_i$ and $\boldsymbol{\mu}_j$ are the expected values. Choi et al.~\cite{Choi2025} provide detailed mathematical and theoretical derivations and proofs of this proposed metric, and a series of properties that $\rho_C^{(M)}$ satisfies which reinforce its interpretability as an analog of $\rho_C$.

Assuming that the correlation between fitting methods is independent of the variables and that the correlation between variables is independent of the fitting methods, inference on a dataset is given by:
\begin{equation}\label{eqn:cccm_estimator}
\begin{aligned}
    \widehat{\rho_C}&^{(M)} = \\
        &\frac{
            \mathbf{1}^\top \widehat{\mathbf{M}}_2 \mathbf{1} - \mathrm{tr}\left(\widehat{\mathbf{M}}_2\right)
            + \mathrm{tr}\left(\widehat{\mathbf{M}}\widehat{\mathbf{M}}^\top\right)
            - \mathbf{1}^\top \widehat{\mathbf{M}}\widehat{\mathbf{M}}^\top \mathbf{1}
        }{
            (d-1)\mathrm{tr}\left(\widehat{\mathbf{M}}_2\right)
            + \mathrm{tr}\left(\widehat{\mathbf{M}}\widehat{\mathbf{M}}^\top\right)
            - \mathbf{1}^\top \widehat{\mathbf{M}}\widehat{\mathbf{M}}^\top \mathbf{1}
        }
\end{aligned}
\end{equation}
\noindent where $\hat{\mathbf{M}} = \frac{1}{n}\sum_{i=1}^n Y_i$ denotes the sample mean matrix and $\hat{\mathbf{M}}_2 = \frac{1}{n}\sum_{i=1}^n Y_i Y_i^\top$ denotes the second-moment matrix. Here, $Y_i$ is the $d \times p$ matrix of measurements for subject $i$, $n$ is the number of samples, and $\mathbf{1}$ is a $d \times 1$ vector of ones. This expression is obtained by substituting the method-of-moments estimators into Eqn.~\ref{eqn:multivariate rho_c} and simplifying, yielding a form that depends only on the measurements. 

The second assumption above limits what soft-constraints can be used in the fitting methods. Instead of eliminating them, it is best to combine the affected metabolites during analysis. At 3\,T, this means using tNAA instead of NAA and NAAG separately. Similarly, tCr, tCho, and Glx should be used based on the achievable spectral separability. The joint set of moieties and metabolites that need to be modeled jointly must be identified for each field strength individually.

For two fitting methods, this formulation reduces to a direct multivariate extension of Lin's coefficient for the pooled data, combining correlation and mean agreement into a single scalar metric. In this setting, the coefficient quantifies agreement between two fitting methods over the full metabolite set. The trace operator summarizes the overall variation across the metabolite set, while the terms involving $\mathbf{1}$ incorporate the corresponding covariance between fitting methods. The mean difference terms penalize systematic disagreement. Thus, $\rho_C^{(M)}$ provides a compact measure of agreement that aggregates agreement across the joint metabolite output space while accounting for systematic deviations between methods.

\section{Analyses \& Visualizations}\label{sec:Reg}
Deep learning-based model fitting requires more comprehensive analysis than traditional methods. Unlike iterative fitting algorithms that refine predictions until convergence criteria are met, \DL models perform a single forward pass at inference and cannot correct inadequate predictions on a per-sample basis. Ideally, \DL models should produce accurate predictions for each output variable, avoid introducing spurious dependencies between variables beyond those present in the underlying data, and exhibit a highly linear relationship between predictions and ground truth values across the full parameter space. Prediction errors should be centered about zero with narrow dispersion, reflecting both accuracy and precision in the sense established in Sec.~\ref{sec:eval:subsec:multivariate:subsubsec:error}. The following sections present methodologies for characterizing regression performance and identifying model failure modes that may be obscured by aggregate metrics.

\subsection{Correlation Analysis}
Correlation analysis evaluates the degree of linear dependence between multivariate model outputs, with the goal of identifying spurious correlations introduced by the model rather than those present in the underlying data. Traditional model fitting methods based on linear combination modeling often impose soft- or hard-constraints to reduce the parameter space and simplify optimization. These constraints can introduce systematic bias that can improve or degrade performance depending on whether the underlying assumptions are well-founded for the given data. In \DL models, training on biased data can cause the model to reproduce that bias in its predictions, effectively learning spurious feature correlations from the training distribution that produce systematic prediction errors when the true signal deviates from those biases. Demographic data such as age, sex, ethnicity, and pathology can violate these assumptions and cause systematic prediction errors \cite{Banerjee2023}. Encouraging output independence encourages the model to identify and rely on features that are genuinely informative for for the predicted quantities rather than on covariance structures that are specific to the training distribution.

\begin{figure}%[!B]
    \centering
    \includegraphics[width=0.4\textwidth,keepaspectratio=True]{images/Shamaei2023.png}
    \caption{Heatmap used to visualize the correlation analysis of the predicted variables. Unbiased estimators should show minimal correlation between predicted variables. Adapted from Shamaei et al. \cite{Shamaei2023}}
    \label{fig:correlation matrix heatmap}
\end{figure}

Figure~\ref{fig:correlation matrix heatmap} shows a heatmap of the correlation matrix of model outputs on a test set. Ideally, the predicted correlation matrix should reproduce the correlation structure of the ground truth. For synthetic datasets in which the modeled parameters are sampled independently, this corresponds approximately to an identity matrix. In practice, small off-diagonal values are acceptable, but notable exceptions often reflect known or logical physiological or spectral relationships. For example, creatine (Cr) and phosphocreatine (PCr), and choline (Cho) and glycerophosphorylcholine (GPC), may exhibit strong inverse correlations due to overlapping spectral peaks and are typically reported as total creatine (tCr) or total choline (tCho) \textit{in vivo}, due to insufficient spectral resolution to separate them \cite{DeGraaf2007}. Similarly, alanine (Ala) and lactate (Lac) exhibit moderate positive correlation, attributable to their spectral proximity combined with their susceptibility to artifacts such as noise and line broadening \cite{DeGraaf2007}.

Strong off-diagonal correlations can indicate an ill-posed estimation problem, in which the same observation is consistent with multiple combinations of component values, i.e. a one-to-many mapping. Such ill-posedness reflects insufficient identifiability in the inverse problem, where the observed spectrum does not contain enough information to uniquely distinguish between multiple combinations of fitting parameters. When such an identifiability problem impedes the model from learning a unique mapping from input to output, it induces a non-convex optimization landscape which increases the difficulty of training and can increase the risk of poor generalization. In such cases, a \DL model may then rely on learned priors or correlations in the training data to select among plausible solutions, increasing the risk of systematic errors when those relationships change. Correlation analysis must therefore be interpreted in the context of both the \MRS acquisition and deep learning to be meaningful. Visualizations should use a consistent, symmetric color scale spanning $[-1, 1]$ centered at zero to facilitate comparison across models and datasets.

\subsection{Maps and Marginal Distributions}\label{sec:Analy&Visual:subsec:maps and marg}
Rizzo et al.~\cite{Rizzo2023} demonstrated that training neural networks on synthetic data with uniformly sampled simulation parameters does not guarantee that the network learns a balanced, uniform representation. The learned input-to-output mapping is shaped by the optimization dynamics and loss landscape, 
not solely by the underlying parameter distributions. Maps and marginal distributions, shown in Fig.~\ref{fig:maps and distributions}, provide a visualization of this by comparing the distribution of ground truth values against predicted outputs as a scatter plot, with marginal histograms illustrating the respective distributions of each axis. Each plot includes an overlay of both the identity line (representing perfect agreement) and the line of linear regression (representing the model's actual trend). The deviation between these two lines reveals systematic bias or scale discrepancies in the model's predictions.

Color-coding data points by metabolite-specific \SNR (\mSNR, defined in ec.~\ref{sec:dataModel:subsec:curation:subsubsec:SNR}) provides further diagnostic insight. Poor performance across all \mSNR values suggests a fundamental model failure, while performance degradation at low \mSNR indicates sensitivity to noise. Cumulative, model-level plots facilitate direct comparison across models (consistent with Sec.~\ref{sec:Eval:subsec:RegMetrics}), while per-metabolite breakdowns identify specific points of failure. Relevant metrics, including $r$, $r^2$, $\rho_C$, and the error metrics from Sec.~\ref{sec:eval:subsec:multivariate:subsubsec:error}, can be embedded in the plot for additional quantitative context.

\begin{figure*}
    \centering
    \includegraphics[width=0.9\textwidth,keepaspectratio=True]{images/Rizzo2023_01.png}
    \caption{Maps and marginal distributions of selected metabolites modeled by a neural network. The ground truth values, plotted along the x-axis, were sampled from uniform distributions, but the networks predict marginal distributions that range from unimodal Gaussian to approximately uniform distributions. Adapted from Rizzo et al.~\cite{Rizzo2023}.}
    \label{fig:maps and distributions}
\end{figure*}

\subsection{Model Calibration}\label{sec:Analy&Visual:subsec:calibration}
Model calibration measures the agreement between a model's nominal predictive probabilities and their empirical coverage in a given test set. A model is well-calibrated if, for all confidence levels $p \in [0,1]$, the corresponding predictive interval contains the true value with empirical frequency $p$ over a large number of samples~\cite{DeGroot1983, Walpole2012, Wasserstein2019}. While calibration is intuitive in binary classification, its extension to continuous outputs, such as metabolite concentrations, requires a predictive distribution or uncertainty representation from which a cumulative distribution function (\CDF) and predictive intervals can be obtained~\cite{Kuleshov2018}.

Calibration plots evaluate calibration quality by comparing, for each confidence level, the predicted fraction of values expected to fall within a given confidence interval against the empirically observed fraction. Confidence levels are defined into discrete intervals to enable quantitative comparison between predicted confidence and empirical accuracy. An interval is considered well-calibrated when the empirical frequency matches the nominal confidence level. The Expected Calibration Error (\textit{ECE}), Eqn.~\ref{eqn:ECE}, provides a global summary using a weighted average of per-bin calibration errors:
\begin{equation}\label{eqn:ECE}
    \begin{aligned}
        a_m &= \frac{1}{N}\sum_{n=1}^{N}
        \mathbb{1}\left[F_n(y_n)\leq c_m\right] \\
        ECE &= \sum_{m=1}^{M} w_m\left|a_m-c_m\right|,
        \qquad \sum_{m=1}^{M}w_m=1
    \end{aligned}
\end{equation}
\noindent where $M$ is the number of evaluated confidence levels, $c_m$ is the nominal confidence level, and $a_m$ is its empirical coverage (accuracy), calculated as the fraction of observations for which the true value lies below the corresponding predicted quantile \cite{Kuleshov2018}. The weight $w_m$ determines the contribution of each confidence level to the aggregate error and is normalized such that $\sum_{m=1}^{M}w_m=1$. Consequently, both $a_m$ and $c_m$ lie within $[0,1]$, and the resulting \textit{ECE} is also bounded to $[0,1]$, where $0$ represents perfect calibration. This formulation retains the absolute-deviation structure of \textit{ECE} while defining calibration according to the empirical-coverage framework for continuous regression proposed by Kuleshov et al.~\cite{Kuleshov2018}.

However, \textit{ECE} is known to have several limitations \cite{Nixon2019}. The first is that fixed equal-width bins create uneven bin populations due to the sharpness of \DL predictions. Then, the bias-variance tradeoff in bin number selection means \textit{ECE} is sensitive to this hyperparameter. And finally, cancellation effects between overconfident and underconfident predictions within the same bin can produce misleadingly low \textit{ECE} values for poorly calibrated models. These limitations do not preclude the use of \textit{ECE}, but require its value to be interpreted in the context of the underlying binning scheme and calibration curve. Because calibration-error estimates can depend on how the confidence range is discretized, adaptive equal-population binning can be used as a sensitivity analysis to reduce dependence on sparsely populated fixed-width bins \cite{Nixon2019}. Because \textit{ECE} reduces calibration performance to a single aggregate value, it should be interpreted jointly with the calibration curve, which preserves the direction and confidence-level dependence of the observed miscalibration. Calibration should additionally be evaluated across relevant model operating conditions, such as metabolite concentration and \mSNR, to determine whether apparently adequate global calibration conceals condition-specific failures.

Figure~\ref{fig:calibration plots} displays calibration plots and \textit{ECE} for four metabolites. NAA demonstrates excellent calibration, with empirical coverage closely matching nominal confidence across confidence levels. 
Asp and Lac exhibit underconfidence (empirical coverage exceeds predicted confidence) at lower concentration values (confidence $< 0.4$) and overconfidence (empirical coverage falls below predicted confidence) at higher concentration values (confidence $> 0.6$). This pattern is consistent with predictive uncertainty estimates that do not adequately capture the heteroscedasticity of the prediction errors, resulting in systematic miscalibration across confidence levels. The elevated \textit{ECE} values corroborate the overall miscalibration. Whether this behavior is systematic across metabolites or reflects metabolite-specific properties of the training distribution warrants further investigation.

\begin{figure}
    \centering
    \includegraphics[width=0.8\columnwidth,keepaspectratio=True]{images/Rizzo_02.png}
    \caption{Reliability diagrams with the expected calibration errors for selected metabolites predicted by a neural network. Adapted from Rizzo et al.~\cite{Rizzo2022}.}
    \label{fig:calibration plots}
\end{figure}

\subsubsection{Model Recalibration} 
Recalibration addresses two distinct but related problems that arise in \DL-based model fitting: point-estimate recalibration and confidence recalibration. 

\paragraph{Point-Estimate Recalibration.}
Point-estimate recalibration corrects systematic bias in predicted concentration values \cite{Rizzo2022}. As discussed in Sec.~\ref{sec:eval:subsec:errorMetrics}, \DL models trained on uniform distributions tend to exhibit bias at the edges of the parameter space, where predictions regress toward the center of the training distribution. The effect of any recalibration should always be evaluated jointly with the marginal distribution of point predictions across the test set to ensure that correcting aggregate bias does not distort individual prediction accuracy. Na\"{i}ve corrections, such as directly inverting the calibration line, should be avoided \cite{Rizzo2022}. While such methods enforce a statistical property on the aggregate distribution, they do so without regard for the correspondence between individual predicted values and their ground truths. That effectively redistributes predictions without improving their local accuracy. Effective point-estimate recalibration should incorporate domain-relevant covariates, such as spectral noise levels or metabolic patterns, and where possible should be embedded within the model's training process or architecture rather than applied as a post-hoc correction.

\begin{figure*}
    \centering
    \includegraphics[width=\textwidth,keepaspectratio=True]{images/Rizzo_03.png}
    \caption{shows an example of pitfalls for naive recalibration for lactate. On the left, the original model results are reported where lactate’s predictions are gaussianly skewed to the mean value (y-axis) in correspondence to a uniformly testing set (x-axis). The analysis of its calibration confirms the findings, reporting low concentrations that tend to be overestimated (system is overconfident for confidence < 0.4) and high concentrations that tend to be underestimated (system underconfident for confidence > 0.6). On the right, naively recalibrating leads to a perfect reliability diagram with the predictions that statistically display as uniformly distributed, in agreement with the statics of the ground truth. Unfortunately though,  estimates of individual values are far off from the 1:1 relationship between estimate vs. ground truth and are instead distributed across the labeled concentration range to yield a correct uniform distribution of predicted values. Figure adapted from citation.}
    \label{fig:recalibration}
\end{figure*}

\paragraph{Confidence Recalibration.}
Confidence recalibration adjusts a model's predicted uncertainty estimates so that predicted predictive intervals match empirical coverage \cite{Niculescu-Mizil2005, Kuleshov2018}. Post-hoc methods such as temperature scaling \cite{Guo2017} and isotonic regression \cite{Niculescu-Mizil2005} can correct aggregate miscalibration but do not guarantee local calibration across the full concentration range. Temperature-scaling approaches have been extended to continuous regression by calibrating prediction intervals \cite{Keren2018} or, for parametric predictive distributions, by rescaling the predicted standard deviation or variance using a single fitted parameter \cite{Levi2022}. More flexible extensions, such as bin-wise variance scaling, allow the correction to vary across the predicted uncertainty range \cite{Pernot2023}.

As with point-estimate recalibration, embedding confidence recalibration within the training process may be beneficial where feasible, ensuring that concentration estimates and their uncertainty representations are jointly optimized. Post-hoc recalibration additionally assumes that the relationship learned on the calibration data remains applicable to the target deployment distribution. This assumption warrants scrutiny in clinical \MRS applications where distributional shift between acquisition sites or patient populations is a practical concern.

\subsection{Visualizing Comparison Metrics}
Quantification performance should be evaluated across multiple complementary scales. Aggregate metrics provide an overall assessment of model behavior, whereas metabolite-specific analyses, \SNR-dependent evaluations, confidence bounds, and multivariate comparisons reveal failure modes that may be obscured by global summaries. The following sections describe recommended visualization and analysis strategies for each of these perspectives.

\subsubsection{Model-Level Metrics}
Proposed models must be benchmarked against gold-standard and state-of-the-art methods. Quantitative metrics allow for rapid model comparison, but can be effectively complemented by visualizations. While results are commonly presented in tables, heatmaps facilitate simultaneous comparison across both the metabolite and model dimensions, making patterns of consistent failure or improvement immediately apparent. Figure~\ref{fig:comparing heatmaps} illustrates two examples comparing MAPE and CV values for seven metabolites across 15 models. For metrics with a meaningful midpoint (such as percent error centered at zero), diverging color maps with high contrast between low, middle, and high values are recommended. Sequential color maps are more appropriate for metrics bounded at zero. Colorblind-safe palettes should be used to ensure accessibility in publications.

\begin{figure}
    \centering
    \begin{subfigure}{0.48\textwidth}
        \includegraphics[width=\textwidth,keepaspectratio]{images/Gupta_01.png}
    \end{subfigure}
    \begin{subfigure}{0.48\textwidth}
        \includegraphics[width=\textwidth,keepaspectratio]{images/Gupta_02.png}
    \end{subfigure}
    \caption{Heatmaps showing visualizations of comparison metrics for multiple variables across a large number of models. Adapted from Gupta et al.~\cite{Gupta2023}.}
    \label{fig:comparing heatmaps}
\end{figure}

\subsubsection{Component-Level Metrics}
As discussed in Section~\ref{sec:Analy&Visual:subsec:maps and marg}, scatter plots with marginal distributions compare single-model outputs against ground truth values. Overlaying predictions from multiple models within a single scatter plot often leads to visual clutter and can obscure differences between model behaviors.

A more effective approach is the violin plot, shown in Fig.~\ref{fig:violin plots}. Violin plots display a kernel density estimate (KDE) of the distribution mirrored about a central axis, combined with an inner box-and-whisker element showing the interquartile range (IQR) and median. Compared to standard box plots, violin plots convey skewness and multimodality, making them more informative for evaluating distributional shape and dispersion \cite{Rizzo2023}. Box plots can obscure distributional asymmetry, such as skewness, that is relevant when interpreting prediction bias because they do not display the shape of the underlying distribution. When sample sizes are small, however, the KDE in a violin plot can suggest artificial smoothness. The reliability of the estimated density depends on sample size, bandwidth selection, and the underlying distribution, so violin plots based on small samples should be interpreted cautiously and complemented by individual observations or robust summary statistics where appropriate. For readability when comparing many metabolites simultaneously, descriptive statistical metrics should be embedded directly within the violin plot rather than reported separately.

Results can be organized either by metabolite, as in Fig.~\ref{fig:violin plots}, or binned by \mSNR and plotted by bin. The latter is particularly informative for characterizing model sensitivity to noise across the concentration range. Both presentations are recommended where figure space permits. Examples are given by Rizzo et al. \cite{Rizzo2023} in Supporting Information Figures S11 and S12.

\begin{figure}
    \centering
    \includegraphics[width=0.95\columnwidth,keepaspectratio=True]{images/Gupta_03.png}
    \caption{Violin plots depicting box-and-whisker plots and the shape of the distributions to visualize how a model (purple) captures the underlying target distribution (green) compared to the reference method (orange). Adapted from Gupta et al.~\cite{Gupta2023}. Will replace with updated version from Jin X.’s thesis (2024)}
    \label{fig:violin plots}
\end{figure}

\subsection{Dependence of Prediction Quality on SNR}
\label{sec:Analy&Visual:subsec:SNR dependence}
The performance of \DL models for quantitative inverse problems depends heavily on sample noise levels, complicating concise characterization of model performance because both accuracy and precision can degrade with decreasing \SNR. In model fitting, this dependency is metabolite-specific: metabolites with lower intrinsic signal contributions can be disproportionately affected, meaning that a single global \SNR (gSNR) metric is insufficient and per-metabolite \SNR characterization is required. 

Figure~\ref{fig:SNR dependence} illustrates a method for visualizing this dependence. Test data are binned by either their \mSNR or the inverse of it, which expresses noise relative to the metabolite-specific signal. For each bin, precision metrics are evaluated over the aggregate distribution of predictions within the bin, while accuracy is reported for individual sample predictions. This asymmetry reflects the distributional nature of precision, which requires sufficient samples to characterize, and the sample-specific nature of accuracy. Following the metric justification in Section~\ref{sec:eval metrics}, this work recommends reporting the nonparametric percent error plotted against the inverse \mSNR to characterize accuracy, and the nonparametric coefficient of variation to characterize precision. Percent error is preferred for its scale-independence across metabolites, and the nonparametric CV for its robustness to the outliers that are common at low \SNR.

The requirements of this visualization differ between model development and deployment-oriented characterization. 
During model development and initial characterization, relative metrics such as percent error and CV are sufficient for comparing model behavior across \SNR regimes. When characterizing a model for clinical deployment, it additionally becomes necessary to evaluate each metabolite individually and, where feasible, to report absolute concentration differences in millimoles, since clinically relevant concentration changes may be more meaningfully expressed in absolute units than by relative differences alone. \hl{As with relative }\CRLB \hl{filtering of poor quality data, relative measures may bias the outcome }\cite{Kreis2016}.

Although \mSNR-dependent analysis characterizes an important source of performance degradation, \mSNR alone does not describe the full difficulty of the estimation problem. Spectral overlap, linewidth, concentration, frequency offsets, and interactions between fitting parameters can alter parameter identifiability even under similar \mSNR conditions. This motivates the more general information-standardized analysis introduced below.

\subsection{Information-Standardized Failure Mode Analysis}
\label{sec:Analy&Visual:subsec:infoStand}
The \mSNR-dependent analysis above and the error metrics introduced in Section~\ref{sec:eval metrics} quantify the discrepancy between the estimated and ground-truth metabolite concentrations thereby providing a direct characterization of model performance. However, the magnitude of an observed error cannot by itself indicate whether that error arises because the underlying spectral estimation problem is intrinsically difficult or because the estimator introduces additional error. This distinction is particularly relevant in \MRS, where the precision with which individual metabolite concentrations can be estimated depends strongly on acquisition and spectral properties such as \SNR, metabolite concentration, linewidth, frequency offsets, and spectral overlap.

To address this limitation, an information-standardized characterization framework is introduced here. Rather than normalizing prediction errors according to \textcolor{violet}{\hl{metabolite magnitude or signal power alone}}, the proposed analysis compares the empirical behavior of the \DL estimator with an information-based reference derived from the known synthetic signal-generation model. This provides a means of evaluating model performance relative to the amount of information theoretically available for estimating the underlying parameters under a given set of simulated acquisition and spectral conditions.

The purpose of this analysis is therefore not to replace conventional error metrics, but to complement them by addressing a different question: "Are the observed errors commensurate with the intrinsic difficulty of the estimation problem, or does the \DL estimator introduce additional bias or variability under particular conditions?" By systematically examining this relationship across different spectral regimes, the framework can be used to identify conditions under which quantitative estimates remain reliable as well as regimes in which model performance begins to degrade.

\subsubsection{Motivation and Scope}
\label{sec:Analy&Visual:subsec:infoStand:subsubsec:}
The difficulty of metabolite quantification is not uniform across spectra or metabolites. Even for the same underlying estimation method, the achievable precision depends on the information contained in the measured signal. A low-concentration metabolite with substantial spectral overlap may, for example, be intrinsically more difficult to estimate than an isolated metabolite with high signal amplitude. Consequently, identical absolute errors can have substantially different interpretations depending on the underlying spectral conditions.

This issue cannot be fully addressed through concen\-tration-based or signal-based error normalization. Such approaches characterize the magnitude of the prediction error relative to a selected scalar property but do not account for the joint effects of noise, spectral overlap, linewidth, parameter interactions, and covariance between metabolite estimates. An information-based reference provides a more general framework because these factors influence the sensitivity of the measured spectrum to changes in the parameters being estimated.

For synthetic data, the underlying signal model and ground-truth parameters are known. This makes it possible to calculate the Fisher information associated with each simulated condition and, from it, a Cram\'{e}r–-Rao lower bound (\CRLB) reference describing the parameter-estimation precision supported by the assumed signal and noise model. In the present framework, the \CRLB is therefore used as a characterization of the information available in the parameter-estimation problem, against which the empirical behavior of the \DL estimator can be compared.

Importantly, this does not imply that the \DL estimator itself is assumed to satisfy the classical conditions underlying the \CRLB. A trained neural network may be biased and may exploit prior information learned from the training distribution. The \CRLB should therefore not be interpreted here as an uncertainty estimate of the network predictions or as a strict lower bound that the empirical variance of the \DL estimator must necessarily satisfy. Instead, it serves as an information-based reference that permits performance to be interpreted in relation to the intrinsic difficulty of each simulated estimation condition.

This distinction allows two qualitatively different situations to be separated. If the uncertainty predicted by the information-based (\CRLB-derived) reference and the empirical variability of the \DL estimator increase together as the estimation problem becomes more difficult, the deterioration in performance may largely reflect a reduction in the information available in the spectrum. In contrast, if the empirical bias or variability of the \DL estimator increases disproportionately relative to the information-based reference, this may indicate a model-specific failure regime. The principal objective of the analysis is therefore not to obtain a single normalized performance score, but to characterize how estimator behavior changes as the underlying estimation problem becomes progressively more challenging.

\subsubsection{Fisher Information and CRLB Reference}
\label{sec:Analy&Visual:subsec:infoStand:subsubsec:fisher reference}
Consider a synthetic spectrum generated according to:
\begin{equation}
\begin{aligned}
    \mathbf{x}=f(\boldsymbol{\theta})+\boldsymbol{\epsilon}\textnormal{, where }&\boldsymbol{\epsilon}=\epsilon_R+i\epsilon_I\textnormal{ and } \\ &\epsilon_R, \epsilon_I \sim \mathcal{N}(0,\sigma^2),
\end{aligned}
\end{equation}
where $\mathbf{x}$ denotes the observed spectral data, $f(\boldsymbol{\theta})$ represents the signal-generation model parameterized by $\boldsymbol{\theta}$ with the real and imaginary components concatenated, and $\boldsymbol{\epsilon}$ denotes additive noise with covariance matrix $\boldsymbol{\Sigma}$. The parameter vector contains the metabolite concentrations of interest together with additional spectral or nuisance parameters included in the generative model. Assuming independent noise with constant variance $\sigma^2$, the noise covariance matrix is $\boldsymbol{\Sigma} = \sigma^2\mathbf{I}$. 

However, if zero-filling is used in the synthetic data to match an \textit{in vivo} protocol and the Fisher information is evaluated in the resulting frequency-domain representation, the noise samples are no longer statistically independent because the additional spectral points are derived from the same finite set of noisy time-domain samples. In this case, $\boldsymbol{\Sigma}$ is generally non-diagonal and should account for the covariance induced by zero-filling and Fourier transformation rather than being represented simply as $\sigma^2\mathbf{I}$. This situation is particularly relevant to the commonly used first $2\times$ zero-fill of complex half-echo \MRS data, which transfers the independent information in the imaginary spectrum into the real spectrum such that the two become a Hilbert pair. Accordingly, if the Fisher information is evaluated after this transformation, a nonredundant representation should be used. For example, the $2\times$ zero-filled real spectrum may be analyzed using its propagated covariance structure. Concatenating both the real and imaginary components after the zero-fill would retain redundant information and may produce a singular covariance representation. Alternatively, the Fisher information can be evaluated before zero-filling using the concatenated real and imaginary components of the original complex data, for which the independent-noise covariance remains diagonal under the assumptions above. 


For a differentiable signal model, the local sensitivity of the measured spectrum to changes in the parameters is described by the Jacobian:
\begin{equation}
    \mathbf{J} = \frac{\partial \tilde{f}(\boldsymbol{\theta})}{\partial \boldsymbol{\theta}}\textnormal{ where }
    \tilde{f}(\theta) =
    \begin{bmatrix}
        \textnormal{Re} \{f(\boldsymbol{\theta})\} \\
        \textnormal{Im} \{f(\boldsymbol{\theta})\}
    \end{bmatrix}.
\end{equation}

Under the assumed noise model, the corresponding Fisher information matrix can be expressed as:
\begin{equation}
    \mathbf{F} = \mathbf{J}^{\mathsf T} \boldsymbol{\Sigma}^{-1} \mathbf{J}.
\end{equation}

The Fisher information matrix describes how strong\-ly the observed data constrain the parameters of the signal model. Parameters that produce clearly distinguishable changes in the spectrum contribute greater information, whereas weak signals, increased noise, broad linewidths, or substantial spectral overlap can reduce identifiability and introduce strong dependencies between parameter estimates.

The corresponding covariance lower-bound matrix is obtained from the inverse Fisher information matrix, provided that $\mathbf{F}$ is nonsingular,
\begin{equation}
    \mathbf{C}_{\mathrm{CRLB}} = \mathbf{F}^{-1},
\end{equation}
with the relevant metabolite-concentration elements, or covariance block, extracted for subsequent analyses. The diagonal elements provide metabolite-specific variance references,
\begin{equation}\label{eqn:sigma crlb}
    \sigma_{\mathrm{CRLB},m}^{2} = \left[\mathbf{C}_{\mathrm{CRLB}}\right]_{mm},
\end{equation}
while the off-diagonal elements describe the covariance structure predicted by the signal model. The latter is particularly important for metabolites with overlapping spectral contributions, for which estimation errors cannot necessarily be interpreted independently.

Within this work, $\mathbf{C}_{\mathrm{CRLB}}$ is used exclusively as a reference derived from the known synthetic generative model. It characterizes the information available for estimating metabolite concentrations under the assumed spectral model, acquisition conditions, and noise statistics. It does not represent predictive uncertainty produced by the \DL model itself. The empirical behavior of the \DL estimator is instead determined independently through repeated noisy realizations, allowing its bias and variance to be compared with this information-based reference in the analyses described below.

\subsubsection{Univariate Information-Standardized Error Analysis}
\label{sec:Analy&Visual:subsec:infoStand:subsubsec:}
The univariate analysis compares the empirical variability of the \DL estimator with the metabolite-specific precision reference derived from the Fisher information. The comparison is performed using repeated noisy realizations of the same underlying synthetic spectrum. Because the ground-truth parameters remain fixed across these realizations, the resulting variability in the network predictions provides an empirical measure of the estimator's sensitivity to noise under a given spectral condition.

For metabolite $m$, let $\hat{c}_{m,k}$ denote the concentration estimated from realization $k$, with $k=1,\ldots,K$, and let $\overline{\hat{c}}_m$ denote the mean prediction across realizations. The empirical variance of the \DL estimator is estimated as:
\begin{equation}
    \hat{\sigma}_{\DL,m}^{2} = \frac{1}{K-1}\sum_{k=1}^{K}\left(\hat{c}_{m,k} - \overline{\hat{c}}_m \right)^{2}.
\end{equation}

The corresponding standard deviation is then:
\begin{equation}
    \hat{\sigma}_{\DL,m} = \sqrt{\hat{\sigma}_{\DL,m}^{2}}.
\end{equation}

For the same underlying parameter setting and acquisition condition, the metabolite-specific \CRLB reference is obtained from the corresponding diagonal element of the covariance lower-bound matrix, Eqn.~\ref{eqn:sigma crlb}, and
\begin{equation}
    \sigma_{\CRLB,m} = \sqrt{\left[\mathbf{C}_{\CRLB}\right]_{mm}}.
\end{equation}

A descriptive information-standardized variability ratio can then be defined as:
\begin{equation}
    R_m = \frac{\hat{\sigma}_{\DL,m}}{\sigma_{\CRLB,m}}.
\end{equation}

This ratio expresses the empirical variability of the \DL estimator relative to the precision reference associated with the information available in the synthetic spectrum. Values of $R_m$ greater than one indicate that the empirical variability of the \DL estimates exceeds the corresponding \CRLB reference. Increasing values therefore indicate progressively greater loss of precision relative to the information content of the estimation problem.

The ratio should not, however, be interpreted as a formal measure of statistical efficiency. In particular, the classical interpretation of the \CRLB assumes an unbiased estimator satisfying the relevant regularity conditions, whereas the \DL estimator may be biased and may incorporate information learned from the training distribution as a prior. Consequently, $R_m$ is used here as a descriptive characterization quantity rather than as evidence that the network does or does not attain a theoretical efficiency bound.

More importantly, the interpretation of $R_m$ should focus on how it changes across controlled experimental conditions rather than on its value in isolation. Consider, for example, a progressive reduction in \SNR. If both $\sigma_{\CRLB,m}$ and $\hat{\sigma}_{\DL,m}$ increase at similar rates, the increasing uncertainty is consistent with an estimation problem that is becoming intrinsically more difficult. If $\hat{\sigma}_{\DL,m}$ instead increases substantially faster than $\sigma_{\CRLB,m}$, this suggests that the estimator is losing precision beyond the deterioration expected from the reduced information available in the spectrum.

The same analysis can be performed as a function of metabolite concentration, linewidth, spectral overlap, frequency offset, and other controlled parameters of the synthetic signal model. The resulting characterization therefore provides a means of distinguishing conditions that are difficult because the data contain limited information from conditions under which the \DL estimator itself exhibits disproportionate degradation. Cases for which $R_m$ is smaller than one require a separate interpretation and are considered in Section~5.6.7.

\subsubsection{Multivariate Information-Standardized Error Analysis}
\label{sec:Analy&Visual:subsec:infoStand:subsubsec:}
Metabolite-wise comparisons provide an intuitive characterization of estimator performance, but they do not account for dependencies between metabolite estimates. This limitation is particularly relevant when spectral contributions overlap. Under such conditions, several combinations of metabolite concentrations may produce similar changes in the measured spectrum, leading to correlated parameter uncertainty.

The off-diagonal elements of $\mathbf{C}_{\CRLB}$ encode this covariance structure. Consequently, the full covariance matrix can be used to characterize the joint prediction error relative to the information geometry of the corresponding estimation problem.

For a vector containing the estimated metabolite concentrations,
\begin{equation}
    \hat{\mathbf{c}} = 
    \begin{bmatrix}
        \hat{c}_1 &
        \hat{c}_2 &
        \cdots &
        \hat{c}_M
    \end{bmatrix}^{\mathsf{T}},
\end{equation}
and the corresponding ground-truth concentration vector,
\begin{equation}
    \mathbf{c} = 
    \begin{bmatrix}
        c_1 &
        c_2 &
        \cdots &
        c_M
    \end{bmatrix}^{\mathsf{T}},
\end{equation}
the joint error vector is:
\begin{equation}
    \mathbf{e} = \mathbf{c} - \hat{\mathbf{c}}.
\end{equation}

An information-standardized squared error can then be defined as:
\begin{equation}
    D^{2} = \mathbf{e}^{\mathsf{T}}\mathbf{C}_{\CRLB}^{-1}\mathbf{e}.
\end{equation}

This quantity weights the observed error according to the covariance structure implied by the synthetic signal model. Errors occurring along parameter combinations that are well constrained by the spectrum contribute more strongly than equivalent Euclidean errors occurring along poorly identifiable directions. Conversely, correlated errors along metabolite combinations that are intrinsically difficult to distinguish are interpreted in the context of that reduced identifiability.

The multivariate formulation therefore addresses a different question from independent metabolite-wise error analysis. Rather than asking only how large the error is for each metabolite, it asks whether the joint pattern of concentration errors is large relative to the information structure of the estimation problem.

This distinction is important for spectrally overlapping metabolites. An apparent error in one metabolite may be coupled to an opposing error in another because the two signals are difficult to distinguish under the particular spectral condition. Treating the two deviations independently ignores this interaction. The covariance-aware formulation instead evaluates the resulting error vector jointly and therefore provides a more appropriate characterization of coupled estimation errors.

As with the univariate variance ratio, $D^{2}$ is interpreted as a descriptive information-standardized measure of how large the joint concentration error is relative to the covariance structure implied by the information available in the synthetic spectrum. Its interpretation does not require assuming that the \DL estimator is unbiased or that its errors follow the probability distribution implied by the classical \CRLB framework. The primary objective is to determine how the magnitude and direction of joint errors change as the information structure of the estimation problem becomes progressively less favorable.

\subsubsection{Monte Carlo Characterization of Failure Modes}
\label{sec:Analy&Visual:subsec:infoStand:subsubsec:MC}
The information-standardized framework is most informative when combined with controlled Monte Carlo (\MC) experiments. Because \DL model performance may exhibit parameter-dependent bias and heteroscedastic error variance across the parameter space, the \MC analysis should preserve both inter-condition and intra-condition variation. A hierarchical design can therefore be used in which an outer sampling procedure generates multiple underlying parameter configurations and an inner \MC procedure generates repeated noise realizations for each fixed configuration. This design separates variation attributable to measurement noise from systematic changes in estimator behavior across the parameter space.

Let $j=1,\ldots,J$ index the underlying synthetic parameter configurations. Each parameter vector is sampled from the prescribed simulation distribution:
\begin{equation}
    \boldsymbol{\theta}_j \sim p_{\mathrm{sim}}(\boldsymbol{\theta}).
\end{equation}

In the present framework, the simulation parameters are sampled from predefined uniform distributions over their respective characterization ranges. Each $\boldsymbol{\theta}_j$ therefore represents one location within the parameter space to be characterized rather than one realization from an assumed clinical population distribution.

For a fixed parameter vector $\boldsymbol{\theta}_j$, realization $k$ can be generated as:
\begin{equation}
    \mathbf{x}_{j,k} = f\left(\boldsymbol{\theta}_{j}\right) + \boldsymbol{\epsilon}_{j,k}.
\end{equation}
where the underlying signal parameters remain constant across realizations and only the noise realization $\boldsymbol{\epsilon}_{j,k}$ changes. All parameters defining experimental condition $j$, including the nominal \SNR, are held constant across the $K$ inner realizations. Consequently, $\boldsymbol{\epsilon}_{j,k}$ is repeatedly sampled from the same condition-specific noise distribution,
\begin{equation}
    \boldsymbol{\epsilon}_{j,k} \sim \mathcal{N}\left(\mathbf{0},\boldsymbol{\Sigma}_j\right),
\end{equation}
where only the random noise realization changes with $k$.

The \DL estimator is subsequently applied to each realization,
\begin{equation}
    \hat{\mathbf{c}}_{j,k} = g_{\DL}\left(\mathbf{x}_{j,k}\right),
\end{equation}
where $g_{\DL}$ denotes the trained concentration estimator.

Repeating this procedure $K$ times for each experimental condition provides an empirical distribution of predictions from which bias, variance, and total prediction error can be estimated. The procedure is then repeated across the $J$ sampled parameter configurations, producing condition-specific estimates of bias, variance, and total prediction error throughout the simulation parameter space. Relevant controlled factors include:
\begin{itemize}
    \item \SNR,
    \item metabolite concentration,
    \item linewidth,
    \item frequency offset,
    \item relaxation-related signal changes, and
    \item other nuisance parameters represented by the synthetic generative model.
\end{itemize}
Spectral overlap has been emphasized as a key factor but is noticeably absent from that list because it is not an independent variable, rather it depends upon the metabolite basis functions, their corresponding linewidths, and their individual frequency shifts. The degree of spectral overlap is determined by the basis functions, linewidth, and the moiety- and metabolite-level frequency shifts. Because these factors may jointly influence estimator performance, the primary characterization should retain their simultaneous variation rather than requiring one parameter to be varied while all others are fixed. Marginal analyses may subsequently be used to characterize the dependence of model behavior on individual parameters, while selected multidimensional analyses can be used to examine interactions between parameters when scientifically relevant.

For each metabolite and condition, the \MC predictions provide estimates of the mean prediction,
\begin{equation}
    \overline{\hat{c}}_{m,j} = \frac{1}{K}\sum_{k=1}^{K}\hat{c}_{m,j,k},
\end{equation}
the empirical bias,
\begin{equation}
    \widehat{\textnormal{Bias}}_{m,j} = \overline{\hat{c}}_{m,j} - c_{m,j},
\end{equation}
and the empirical variance,
\begin{equation}
    \hat{\sigma}_{\DL,m,j}^{2} = \frac{1}{K-1}\sum_{k=1}^{K}\left(\hat{c}_{m,j,k} - \overline{\hat{c}}_{m,j}\right)^{2}.
\end{equation}
The corresponding \RMSE can then be estimated as:
\begin{equation}
    \widehat{RMSE}_{m,j} = \sqrt{\frac{1}{K}\sum_{k=1}^{K}\left(\hat{c}_{m,j,k} - c_{m,j}\right)^{2}}.
\end{equation}

The additional index $j$ is essential because these quantities are estimated conditionally at each location in the parameter space. The resulting characterization therefore treats bias, variance, and total error as functions of the underlying simulation parameters:
\begin{subequations}
\begin{align}
    \textnormal{Bias}_m &= \textnormal{Bias}_m(\boldsymbol{\theta}), \\ \sigma_{DL,m}^{2} &= \sigma_{DL,m}^{2}(\boldsymbol{\theta}), \\ {RMSE}_m &= {RMSE}_m(\boldsymbol{\theta}).
\end{align}
\end{subequations}

This conditional formulation is particularly important when the estimator exhibits heteroscedastic error variance. Pooling all predictions across heterogeneous parameter configurations before estimating variance would conflate intra-condition noise sensitivity with inter-condition differences in systematic error. By the law of total variance \cite{Casella2002}, % Theorem 4.4.7 from Casella et al.
\begin{equation}
\textnormal{Var}(e_m) = \textnormal{E}_{\boldsymbol{\theta}}\left[\textnormal{Var}(e_m\mid\boldsymbol{\theta})\right] + \textnormal{Var}_{\boldsymbol{\theta}}\left(\textnormal{E}[e_m\mid\boldsymbol{\theta}]\right),
\end{equation}
where the first term represents the average conditional error variance and the second represents variation in conditional bias across the parameter space. A variance calculated after pooling all spectra would therefore not provide a direct estimate of noise-related precision.

The information-standardized quantities should likewise be calculated separately for each parameter configuration. For metabolite $m$ and condition $j$, the corresponding \CRLB variance reference is:
\begin{equation}
    \sigma_{\CRLB,m,j}^{2} = \left[\mathbf{C}_{\CRLB}(\boldsymbol{\theta}_j)\right]_{mm},
\end{equation}
from which the condition-specific variability ratio ($R_{m,j}$) and excess variance ($\Delta V_{m,j}$) are obtained as:
\begin{equation}
    R_{m,j} = \frac{\hat{\sigma}_{\DL,m,j}}{\sigma_{\CRLB,m,j}},
\end{equation}
and
\begin{equation}
    \Delta V_{m,j} = \hat{\sigma}_{\DL,m,j}^{2} - \sigma_{\CRLB,m,j}^{2}.
\end{equation}

The purpose of the \MC analysis is not merely to obtain more stable estimates of global model performance. Its primary role is to identify failure regimes. Bias, empirical variance, \RMSE, the univariate information-standardized ratio $R_{m,j}$, and the multivariate quantity $D^{2}$ can be examined as functions of the controlled experimental variables. For the multivariate analysis, the distribution of $D^{2}$ across repeated noise realizations should be summarized within each condition before comparisons across parameter configurations are made. Each simulated parameter condition should therefore constitute the primary unit of analysis for characterization across the parameter space, with metrics first estimated or summarized over repeated noise realizations within that condition. 
For each $\boldsymbol{\theta}_j$, the simulation parameters and corresponding estimates of bias, empirical variance, \RMSE, $\sigma_{\CRLB,m,j}$, $R_{m,j}$, and $\Delta V_{m,j}$ are retained before any aggregation across conditions is performed.

These condition-specific results can subsequently be summarized across the parameter space. Marginal summaries can characterize the distribution of performance over the prescribed simulation domain, while stratified analyses or smooth response functions can characterize how individual metrics vary with parameters such as \SNR, concentration, or linewidth. Because parameter interactions may produce localized failure modes that are obscured by one-dimensional marginalization, selected two-dimensional or multivariate analyses should additionally be considered for combinations of parameters with plausible physical or spectral interactions.

The use of uniform sampling distributions is advantageous for this characterization objective because it provides broad coverage of the predefined operating range without concentrating samples around a particular region of the parameter space. Simple summaries across $j$ therefore describe performance with respect to the uniformly sampled characterization distribution. They should not, however, be interpreted as estimates of expected performance in an \textit{in vivo} or clinical population unless the simulated parameter distribution is representative of that population. Using the recommendations from Sec.~\ref{sec:dataModel:subsec:curation:subsubsec:OOD} means that the ranges sampled by these uniform distributions should encompass the target \textit{in vivo} populations, but they are not representative of those populations.

These analyses make it possible to determine whether performance deteriorates gradually as spectral information decreases or whether particular conditions produce disproportionately large changes in estimator behavior. For example, decreasing \SNR may initially increase both the \CRLB reference and the empirical \DL variance in parallel. Below a particular range of \SNR, however, the empirical variance, or bias, may begin to increase substantially faster. Such behavior would suggest a transition from an intrinsically difficult estimation regime to a regime in which model-specific degradation becomes increasingly important.

The characterization should therefore emphasize response curves, transition regions, and combinations of experimental conditions associated with substantial changes in estimator behavior rather than reducing performance to a single aggregate value. In this way, the \MC experiments provide an empirical map between acquisition and spectral conditions and the reliability of the resulting quantitative estimates.

The resulting failure-mode characterization is particularly relevant for translation beyond the synthetic evaluation setting. A model may achieve low aggregate error while still exhibiting severe degradation within specific regions of the parameter space. Identifying such regions explicitly provides a foundation for defining operating regimes, developing quality-control criteria, and determining conditions under which model outputs should be interpreted with greater caution.

\subsubsection{Bias and Excess Variance}
\label{sec:Analy&Visual:subsec:infoStand:subsubsec:bias}
Variance alone is insufficient to characterize the behavior of a \DL estimator. A model may produce highly consistent predictions across repeated noise realizations while systematically estimating concentrations above or below their ground-truth values. Conversely, an estimator may be approximately unbiased while exhibiting substantial variability. Bias and variance should therefore be evaluated separately.

For metabolite $m$ and parameter condition $j$, the estimator bias at a fixed ground-truth condition is defined as:
\begin{equation}
    \textnormal{Bias}_{m,j} = \textnormal{E}\left[\hat{c}_{m,j}\right] - c_{m,j}.
\end{equation}

In the \MC experiment, this quantity is estimated from the average prediction across repeated realizations,
\begin{equation}
        \widehat{\textnormal{Bias}}_{m,j} = \overline{\hat{c}}_{m,j} - c_{m,j}.
\end{equation}

The corresponding empirical \DL variance is:
\begin{equation}
    \sigma_{\DL,m,j}^{2} = \textnormal{Var}\left(\hat{c}_{m,j}\right),
\end{equation}
with the sample estimate obtained from the repeated realizations described in Section~\ref{sec:Analy&Visual:subsec:infoStand:subsubsec:MC}.

These two quantities characterize qualitatively different failure modes. Large bias indicates a systematic displacement of the predictions from the ground truth, whereas large variance indicates sensitivity of the predictions to different realizations of the measurement noise. An estimator can therefore be stable but biased, variable but approximately unbiased, or both biased and variable.

Their combined contribution to total squared prediction error is described by the standard bias-variance decomposition:
\begin{equation}
    \MSE_{m,j} = \textnormal{Bias}_{m,j}^{2} + \sigma_{\DL,m,j}^{2},
\end{equation}
for a fixed ground-truth parameter setting, where the expectation is taken over repeated noise realizations.

The information-based reference can then be used to determine whether the empirical variance is elevated relative to the precision supported by the parameter-estimation problem. One direct measure of excess variance is:
\begin{equation}
    \Delta V_{m,j} = \sigma_{\DL,m,j}^{2} - \sigma_{\CRLB,m,j}^{2}.
\end{equation}
Positive values of $\Delta V_m$ indicate empirical variability above the corresponding \CRLB variance reference. The same comparison is expressed in relative form through the squared variance ratio,
\begin{equation}
    R_{m,j}^{2} = \frac{\sigma_{\DL,m,j}^{2}}{\sigma_{\CRLB,m,j}^{2}}.
\end{equation}

These comparisons should be interpreted alongside the estimator bias. A low empirical variance relative to the \CRLB reference does not necessarily indicate superior estimation if the predictions exhibit substantial systematic bias. Similarly, a large \RMSE may arise predominantly from bias, predominantly from variance, or from both.

The combined analysis therefore separates three conceptually different components of model behavior: the intrinsic precision limitations represented by the information-based reference, systematic errors introduced by the estimator, and empirical variability of the estimator beyond the information-based reference. Tracking these components across the controlled \MC experiments enables different modes of degradation to be distinguished rather than collapsing them into a single error statistic.

\subsubsection{Interpreting Sub-\CRLB \DL Variance}
\label{sec:Analy&Visual:subsec:infoStand:subsubsec:subcrlb var}
Because the \CRLB is used as a reference rather than as a strict bound on the learned estimator, cases in which the empirical \DL variance falls below the corresponding \CRLB reference should be expected to require careful interpretation. Such observations are defined by:
\begin{equation}
    \sigma_{\DL,m}^2 < \sigma_{\CRLB,m}^2,
\end{equation}
or, equivalently,
\begin{equation}
    R_m^2 < 1.
\end{equation}

These cases should not be described as violations of the \CRLB. The classical \CRLB establishes a lower variance bound under a specific set of assumptions, including the relevant conditions concerning estimator bias and the statistical model used to derive the Fisher information. A learned estimator is not guaranteed to satisfy those assumptions. In particular, a \DL estimator may trade variance for bias. A network can produce predictions that vary less across repeated noisy realizations because its outputs are constrained toward relationships learned from the training distribution. Reduced empirical variance therefore does not necessarily imply improved estimation accuracy. The corresponding bias must be examined simultaneously.

Learned prior information provides a second important consideration. The \DL model is trained on a distribution of spectra and parameter combinations and can therefore exploit statistical structure that is not represented by a conventional local Fisher information calculation for an individual spectrum. The information available to the trained estimator is consequently not necessarily identical to the information represented by the synthetic \CRLB reference.

A further consideration is the scope of the reference model itself. The Fisher information describes parameter sensitivity under the specific generative model, noise assumptions, and parameterization used for its calculation. Differences between this reference formulation and the behavior of the learned estimator may therefore contribute to apparent sub-\CRLB variance.

For these reasons, sub-\CRLB observations should be treated as diagnostic cases rather than removed or interpreted as evidence that either the \DL estimator or the \CRLB calculation is invalid. Their frequency can be reported for each metabolite and across experimental conditions. For $J$ evaluated conditions, the fraction of sub-\CRLB cases for metabolite $m$ may be summarized as:
\begin{equation}
    p_{\mathrm{sub},m} = \frac{1}{J}\sum_{j=1}^{J}\mathbf{1}\left\{\sigma_{\DL,m,j} < \sigma_{\CRLB,m,j} \right\},
\end{equation}
where $\mathbf{1}\{\cdot\}$ denotes the indicator function.

More importantly, these observations should be examined jointly with bias and total error. Sub-\CRLB variance accompanied by low bias and low total error represents a different estimator regime from sub-\CRLB variance accompanied by substantial systematic bias. Reporting these quantities together prevents low empirical variance from being interpreted in isolation as evidence of superior estimator performance.

Taken together, the univariate and multivariate information-standardized analyses, the controlled \MC experiments, and the decomposition of bias and variance provide complementary views of model behavior. The objective is not to establish that the \DL estimator obeys or attains a classical \CRLB. Rather, the \CRLB provides an information-based reference that makes it possible to determine whether changes in model performance broadly follow changes in the intrinsic difficulty of the parameter-estimation problem or whether specific acquisition and spectral conditions produce additional model-specific degradation. This distinction provides a basis for identifying reliable operating regimes and characterizing known failure modes before application to more complex or clinically representative data.

\subsubsection{Visualization}
\label{sec:Analy&Visual:subsec:infoStand:subsubsec:visualization}
The condition-specific quantities obtained from the \MC analysis should be visualized without first collapsing the multidimensional parameter space into a single global summary. Visualization should instead proceed hierarchically from one-dimensional marginal response curves to two-dimensional response surfaces and, where necessary, stratified representations of higher-dimensional interactions. This approach preserves parameter-dependent changes in bias, precision, and information-standardized performance while maintaining interpretable graphical representations.

For each metabolite, one-dimensional marginal response curves can be used to examine the dependence of $\widehat{\textnormal{Bias}}_{m,j}$, $\hat{\sigma}_{\DL,m,j}^{2}$, $\widehat{\RMSE}_{m,j}$, $R_{m,j}$, $\Delta V_{m,j}$, and, where appropriate, a condition-level summary of $D^{2}$ on individual simulation parameters. The remaining simulation parameters can be marginalized over their sampled distributions. 

Marginal response curves should display both the central tendency and the variation in performance across the marginalized parameters. A smoothed median or mean response may therefore be accompanied by an interquartile or other quantile interval. Displaying the conditional spread is particularly important because a stable marginal mean may otherwise conceal localized regions of poor performance caused by interactions with other simulation parameters.

The dependence of metabolite-level performance on signal strength should initially be evaluated using both the metabolite spectral \SNR and the metabolite power-based \SNR. The two definitions characterize signal strength differently and should therefore be retained as complementary characterization axes where they provide distinct information. Their relationships with model performance should be compared directly to determine where the resulting response patterns agree or diverge. One measure need only be omitted if empirical analysis demonstrates that it is effectively redundant for the characterization being performed.

Interactions between simulation parameters should be examined using two-dimensional response surfaces. Heatmaps or contour plots are preferable to perspective-based three-dimensional surfaces because they permit quantitative differences and transition regions to be compared more directly. For a selected pair of parameters, the remaining dimensions are marginalized over their sampled distributions. Particularly informative combinations are expected to include \SNR with metabolite concentration, linewidth, frequency offset, and other factors that plausibly modify spectral identifiability. Both spectral \SNR and power-based \SNR should initially be considered in these analyses.

The selection of parameter pairs should be hypothesis-driven rather than exhaustive. Pairwise visualizations are most informative when the parameters have a plausible physical or spectral interaction. For example, the effect of decreasing \SNR may depend strongly on metabolite concentration or linewidth, whereas other combinations may provide little additional information beyond their corresponding marginal response curves. Exhaustively plotting all possible parameter combinations should therefore be avoided unless exploratory analysis indicates an interaction that warrants further investigation.

Interactions involving three or more parameters should be represented through stratified or faceted two-dimensional response surfaces rather than direct three-dimensional or higher-dimensional plots. For example, an \SNR-by-concentration response surface may be displayed separately for selected linewidth intervals. This preserves the interpretability of the two-dimensional representation while allowing higher-order interactions and localized failure regimes to be identified.

Response curves and surfaces should be examined for gradual deterioration, nonlinear changes in performance, and localized transition regions. Where a predefined and scientifically justified acceptability criterion exists, the corresponding boundary can be superimposed on the response surface and used to define an operating or failure region. For a generic characterization metric $M_m(\boldsymbol{\theta})$ with an upper acceptability threshold $M_{\mathrm{crit}}$, the corresponding failure region may be defined as:
\begin{equation}
    \mathcal{F}_m = \left\{\boldsymbol{\theta}:M_m(\boldsymbol{\theta}) > M_{\mathrm{crit}}\right\}.
\end{equation}

The fraction of the sampled characterization domain satisfying this criterion can then be estimated as:
\begin{equation}
    p_{\mathrm{fail},m} = \frac{1}{J}\sum_{j=1}^{J}\mathbf{1}\left\{M_{m,j} > M_{\mathrm{crit}}\right\}.
\end{equation}

Because the underlying parameter space is sampled uniformly, $p_{\mathrm{fail},m}$ describes the fraction of the predefined characterization domain associated with the specified failure criterion. It should not be interpreted as the prevalence of model failure in an \textit{in vivo} population unless the sampled parameter distribution is representative of that population.

When no independently justified threshold is available, transition regions should not be converted into arbitrary binary failure boundaries. In such cases, the response should remain continuous and be characterized using the shape and slope of the response curve or surface, together with its conditional spread. Regions of accelerated deterioration may then be identified descriptively without assigning an unsupported clinical or statistical cutoff.

Aggregate summaries should complement rather than replace these visualizations. Central tendency, interquartile ranges, and selected upper quantiles can summarize model behavior across the characterization domain, while the response curves and surfaces identify where those aggregate values arise. Signed bias should be retained when visualizing systematic estimation direction, whereas absolute or squared bias is more appropriate when summarizing the magnitude of bias across heterogeneous conditions.

Taken together, these visualizations provide complementary levels of characterization. One-dimensional marginal curves describe the overall dependence of performance on individual parameters, two-dimensional response surfaces reveal interactions and localized failure regions, stratified surfaces extend these analyses to higher-dimensional effects, and threshold-based summaries quantify the extent of predefined failure regions when defensible criteria are available. The objective is therefore not to reduce the multidimensional characterization to a single visualization or aggregate statistic, but to identify where, how, and under which combinations of acquisition and spectral conditions model performance begins to degrade. 

\subsection{Nonparametric Confidence Bounds for Stratified Error Thresholds}
\label{sec:Reg:subsec:pred_intervals}
The preceding analyses characterize where model performance deteriorates and whether that deterioration broadly follows changes in the intrinsic difficulty of the spectral estimation problem. For clinical translation, however, identifying failure regimes is only one part of establishing trustworthiness. It is also necessary to translate this characterization into an interpretable statement about whether an individual reported concentration is likely to satisfy a predefined error requirement. The reliability of an individual concentration estimate can be characterized at two complementary levels. First, the empirical probability that a prediction satisfies a prespecified error tolerance can be estimated as a function of the concentration predicted by the model. Second, this analysis can be extended across the characterized range of intrinsic estimation difficulty to determine a conservative reliability bound that does not depend on the relative frequency with which those conditions occur in the target population.

\subsubsection{Predicted-Concentration-Stratified Error Attainment}
\label{sec:Reg:subsec:pred_intervals:subsubsec:conc strat error}
Aggregate error metrics describe model performance over an evaluation dataset, but they do not directly quantify the reliability of a particular concentration reported by the model. A key clinical question is whether a predicted metabolite concentration can be trusted: given that a model outputs a concentration of $\hat{y}$ mM for a particular metabolite, how confident can we be that the true prediction error lies within a prespecified tolerance? For synthetic data, this question can be framed in terms of \mSNR, since \mSNR is directly observable. For \textit{in vivo} data, however, \mSNR is not directly accessible because it depends jointly on the metabolite concentration and lineshape parameters, neither of which is known with certainty at inference time. The predicted concentration $\hat{y}$, on the other hand, is always available regardless of acquisition context. The predicted concentration can therefore be used to stratify the empirical probability that a future prediction satisfies a prespecified error criterion.

For metabolite $m$, let $B_{m,b}$ denote predicted-concen\-tration stratum $b$, and let $E_{n,m}$ denote the chosen error measure for sample $n$. Depending on the intended interpretation, $E_{n,m}$ may represent absolute concentration error or absolute percent error. For a prespecified error tolerance $\tau_{m,s}$, the conditional error-attainment probability is defined as:
\begin{equation}
    p_{m,b,s} = P\left( E_m \leq \tau_{m,s} \mid \hat{Y}_m \in B_{m,b} \right).
\end{equation}

On an independent synthetic evaluation set, let,
\begin{equation}
    N_{m,b} = \sum_{n=1}^{N} \mathbf{1}\left\{\hat{y}_{n,m} \in B_{m,b}\right\}
\end{equation}
denote the number of predictions in the stratum and
\begin{equation}
    S_{m,b,s} = \sum_{n=1}^{N} \mathbf{1}\left\{\hat{y}_{n,m} \in B_{m,b}\right\} \mathbf{1}\left\{E_{n,m}\leq\tau_{m,s}\right\}
\end{equation}
the number satisfying error tolerance $\tau_{m,s}$. The empirical probability is then
\begin{equation}
    \hat p_{m,b,s} = \frac{S_{m,b,s}}{N_{m,b}}.
\end{equation}

Because attainment of the error criterion is a binary event, uncertainty in $p_{m,b,s}$ can be quantified using an exact binomial confidence bound \cite{Clopper1934}. For a one-sided $100(1-\alpha)\%$ lower confidence bound,
\begin{equation}
    L_{m,b,s} = B^{-1} \left( \alpha; S_{m,b,s}, N_{m,b} - S_{m,b,s} + 1 \right),
\end{equation}%
where $B^{-1}$ denotes the beta-distribution quantile function, with $L_{m,b,s}=0$ when $S_{m,b,s}=0$. The Clopper--Pearson procedure guarantees coverage of at least the nominal confidence level, although its discrete construction can produce conservative bounds \cite{Clopper1934, Brown2001}. Repeating the analysis for multiple prespecified tolerances, for example $5\%$, $10\%$, $15\%$, and $20\%$ error, produces a stratified reliability profile. The number and magnitude of these tolerances may differ between metabolites when justified by analytical or clinical requirements, but should be specified independently of the observed evaluation-set performance.

This methodology requires a synthetic test set, allowing ground truth metabolite concentrations to be known while reflecting the relevant sources of variation observed in the \textit{in vivo} test set. Then for each metabolite, test-set predictions $\hat{y}_n$ are binned over the predicted concentration axis using adaptive equal-population binning based on the principle used by Nixon et al.~\cite{Nixon2019}, which helps maintain more balanced sample counts across bins when predicted concentrations are unevenly distributed. To preserve the inferential interpretation of the final evaluation, adaptive bin boundaries should preferably be determined using an independent development or calibration set and then held fixed when the confidence bounds are calculated on the evaluation set.

For a required reliability level $\gamma$, the resulting metabolite-specific operating envelope for error tolerance $\tau_{m,s}$ can be defined as:
\begin{equation}
    \mathcal O_{m,s,\gamma} = \bigcup_{b:L_{m,b,s} \geq \gamma} B_{m,b}.
\end{equation}
This identifies the predicted-concentration ranges for which the lower confidence bound on error-threshold attainment satisfies the required reliability criterion. A clinician can therefore use the model output itself to identify the corresponding concentration stratum and interpret both the observed error-attainment probability and its conservative lower confidence bound.

This analysis measures reliability relative to prespecified error tolerances and therefore directly addresses whether a prediction is likely to satisfy an analytically or clinically meaningful error criterion. It does not describe the complete distribution or direction of possible errors. Where this additional information is useful, data-driven nonparametric prediction intervals such as those proposed by Frey~\cite{Frey2013} can be reported as a complementary, threshold-independent characterization of the signed error distribution within a concentration stratum. Frey's method constructs a distribution-free prediction interval for one future observation and can provide shorter intervals than conventional equal-tailed nonparametric intervals for skewed distributions, but its confidence coefficient should not be interpreted as the probability of satisfying a separately prespecified error threshold.

The principal limitation of the concentration-stratif\-ied probability is that it is marginal over the range of estimation difficulties represented within each concentration stratum. Consequently, $\hat{p_{m,b,s}}$ depends not only on model behavior but also on the relative frequency with which easier and more difficult estimation conditions occur in the synthetic characterization distribution. Broad uniform simulation provides comprehensive domain coverage but does not make these frequencies representative of an unknown target population. The information-stratified analysis below therefore evaluates whether the concentration-specific reliability estimate can be supported by a conservative lower bound across the range of intrinsic estimation difficulties represented by the synthetic generative model.

\subsubsection{Robust Operating-Envelope Bounds Across Information Conditions}
\label{sec:Reg:subsec:pred_intervals:subsubsec:robust op env}
The concentration-stratified analysis above estimates reliability after averaging over the range of spectral conditions represented within each predicted-concentration stratum. Because the relative frequencies of those conditions in a uniformly sampled synthetic characterization set need not reflect their frequencies in an eventual \textit{in vivo} population, a second analysis can be used to determine whether the concentration-specific reliability estimate remains acceptable across different levels of intrinsic estimation difficulty. Rather than stratifying over the individual physical parameters of the generative model, the synthetic conditions are grouped according to a metabolite-specific oracle information reference derived from the Fisher information analysis in Section~\ref{sec:Analy&Visual:subsec:infoStand}.

For metabolite $m$ and synthetic parameter configuration $\boldsymbol{\theta}$, define the oracle precision reference as:
\begin{equation}
    \sigma_{\mathrm{CRLB,oracle},m}(\boldsymbol{\theta}) = \sqrt{ \left[ \mathbf{C}_{\mathrm{CRLB,oracle}} (\boldsymbol{\theta}) \right]_{mm} }.
\end{equation}%
Here, $\mathbf{C}_{\mathrm{CRLB,oracle}}$ is calculated from the known synthetic generative model and therefore characterizes the information available for estimating metabolite $m$ under the corresponding simulated condition. It is not a \CRLB for the \DL estimator and is not interpreted as predictive uncertainty produced by the network. Its role in the present analysis, as with the previous analyses, is solely to provide a compact, information-based description of the intrinsic difficulty of the estimation problem.

The synthetic characterization domain can then be divided into $G$ information strata for metabolite $m$:
\begin{equation}
\begin{aligned}
    G_{m,g} &= \left\{ \boldsymbol{\theta}: q_{m,g-1} < \sigma_{\mathrm{CRLB,oracle},m}(\boldsymbol{\theta}) \leq q_{m,g} \right\}, \\ g &= 1,\ldots,G,
\end{aligned}
\end{equation}
where the boundaries $q_{m,g}$ define increasing levels of oracle estimation difficulty. The boundaries may be selected to provide approximately balanced numbers of synthetic characterization samples across strata, but should be determined independently of the final evaluation set and then held fixed. Because $\sigma_{\mathrm{CRLB,oracle},m}$ is obtained from the inverse of the full Fisher information matrix, the resulting metabolite-specific reference already reflects the influence of spectral overlap and covariance with other fitted parameters on the information available for estimating metabolite $m$.

For predicted-concentration stratum $B_{m,b}$, error tolerance $\tau_{m,s}$, and oracle information stratum $G_{m,g}$, define the condition-specific error-attainment probability as:
\begin{equation}
    p_{m,b,s,g} = P\left( E_m \leq \tau_{m,s} \mid \hat{Y}_m \in B_{m,b}, \boldsymbol{\theta} \in G_{m,g} \right).
\end{equation}

On the independent synthetic evaluation set, let $N_{m,b,g}$ denote the number of predictions falling jointly within concentration stratum $B_{m,b}$ and oracle information stratum $G_{m,g}$, and let $S_{m,b,s,g}$ denote the number of these predictions satisfying error tolerance $\tau_{m,s}$. The empirical condition-specific reliability is then:
\begin{equation}
    \hat{p}_{m,b,s,g} = \frac{S_{m,b,s,g}}{N_{m,b,g}}.
\end{equation}

Consider an eventual target population whose predictions within concentration stratum $B_{m,b}$ arise from an unknown mixture of the characterized oracle information strata. Let $\pi_{m,b,g}$ denote the corresponding mixture weights, such that
\begin{equation}
    \pi_{m,b,g} \geq 0, \qquad \sum_{g=1}^{G}\pi_{m,b,g} = 1.
\end{equation}

The resulting target-population error-attainment probability can then be written as:
\begin{equation}
    p_{m,b,s}^{\mathrm{target}} = \sum_{g=1}^{G} \pi_{m,b,g} p_{m,b,s,g}.
\end{equation}

Because the mixture weights are nonnegative and sum to one, $p^{\mathrm{target}}_{m,b,s}$ is a weighted average of the stratum-specific error-attainment probabilities. A weighted average cannot be smaller than its smallest component, and therefore:
\begin{equation}
    p_{m,b,s}^{\mathrm{target}} \geq \min_g p_{m,b,s,g}.
\end{equation}

Accordingly, a conservative reliability over the characterized oracle information range can be defined as:
\begin{equation}
    p_{m,b,s}^{\mathrm{rob}} = \min_{g} p_{m,b,s,g}.
\end{equation}
This quantity is independent of the relative frequency with which the characterized information strata occur. It therefore avoids interpreting the uniform synthetic sampling distribution as though it represented the prevalence of estimation conditions in an \textit{in vivo} population.

A one-sided Clopper--Pearson lower confidence bou\-nd can be calculated within each oracle information stratum \cite{Clopper1934}:
\begin{equation}
    L_{m,b,s,g} = B^{-1} \left( \alpha_g; S_{m,b,s,g}, N_{m,b,g} - S_{m,b,s,g} + 1 \right),
\end{equation}
with $L_{m,b,s,g}=0$ when $S_{m,b,s,g}=0$. Because the robust conclusion depends simultaneously on all $G$ information strata, the individual error probabilities $\alpha_g$ should satisfy
\begin{equation}
    \sum_{g=1}^{G} \alpha_g \leq \alpha.
\end{equation}
A simple choice is the Bonferroni allocation,
\begin{equation}
    \alpha_g = \frac{\alpha}{G},
\end{equation}
which yields the simultaneous guarantee \cite{Dunn1961}:
\begin{equation}
    P\left( p_{m,b,s,g} \geq L_{m,b,s,g} \quad \textnormal{for all}\ g \right) \geq 1-\alpha.
\end{equation}

Under this simultaneous guarantee, every stratum-specific error-attainment probability is bounded below by its corresponding confidence bound. Consequently, the minimum reliability across the characterized information strata is itself bounded below by the minimum of the stratum-specific lower bounds. The robust lower confidence bound is therefore defined as:
\begin{equation}
    L_{m,b,s}^{\mathrm{rob}} = \min_g L_{m,b,s,g}.
\end{equation}
The corresponding one-sided robust confidence interval is therefore:
\begin{equation}
    \mathcal I_{m,b,s}^{\mathrm{rob}} = \left[ L_{m,b,s}^{\mathrm{rob}}, 1 \right].
\end{equation}

This construction has an important interpretation for translation. The oracle information strata are used only during synthetic characterization and do not need to be identified for an individual \textit{in vivo} spectrum at inference time. After taking the minimum across the characterized strata, the final reliability bound is indexed only by the predicted metabolite concentration and the prespecified error tolerance. A broad uniform simulation can therefore be used to characterize reliability across the predefined information range without requiring the simulation distribution itself to reproduce the prevalence of those conditions in the eventual target population.

For a required minimum reliability $\gamma$, the robust operating envelope is defined as:
\begin{equation}
    \mathcal{O}_{m,s,\gamma}^{\mathrm{rob}} = \bigcup_{ b: L_{m,b,s}^{\mathrm{rob}} \geq \gamma } B_{m,b}.
\end{equation}

Predicted concentrations within this envelope satisfy the required lower reliability bound across all characterized oracle information strata rather than only on average over the synthetic characterization distribution. The resulting lookup can therefore provide a conservative concentration-specific reliability statement without requiring the information stratum of the future subject to be known.

The robustness of this result remains limited to the information conditions represented by the synthetic characterization model. It does not establish that the same condition-specific error behavior will necessarily hold under unmodeled acquisition effects, pathological states, site differences, or other forms of distribution shift. The resulting envelope should therefore be interpreted as a robust analytical operating envelope over the characterized synthetic domain rather than as a universal clinical guarantee. Representative \textit{in vivo} validation remains necessary to establish transfer of these reliability bounds to the intended clinical population. However, the more the synthetic data can be made to approximate \textit{in vivo} data from a given scanner or site, the more relevant these analytical reliability bounds become to subsequent clinical validation. Increasing the number of oracle information strata can improve localization of reliability changes but reduces the number of evaluation samples available within each stratum, producing a resolution-precision tradeoff which can be minimized with additional synthetic data. Sparse or inadequately characterized strata should therefore be supplemented with additional data or reported as regions of insufficient evidence rather than assigned extrapolated reliability.

\begin{figure*}
    \centering
    \includegraphics[width=\textwidth,keepaspectratio=True]{images/Rizzo_04.png}
    \caption{Metabolite quantification errors $\Delta$ and the RMSE $\sigma$ as a function of the inverse spectral SNR. Lower SNR values yield wider error dispersions. Adapted from Rizzo et al.~\cite{Rizzo2023}}.
    \label{fig:SNR dependence}
\end{figure*}

\subsection{Evaluation of Additional Fitting Parameters}
The operating-envelope analysis above focuses on metabolite concentrations because these constitute the primary quantitative output of \MRS fitting. For models that predict a broader set of spectral fitting parameters, however, reliable concentration estimates alone do not establish that the remaining outputs are physically meaningful. Deep learning-based model fitting, when it predicts the full set of spectral parameters rather than concentrations alone, can provide a more complete and interpretable account of the fitting problem. Predicted lineshape parameters enable fitted spectral reconstructions and residual analysis, and can allow sources of variance that would otherwise inflate concentration errors, such as linewidth variations and frequency shifts, to be explicitly modeled rather than absorbed into the concentration estimates. This transition from concentration-only predictions is arguably a step towards explainable \AI in \MRS. Accurately modeling these additional sources of variance may improve the reliability of the predicted concentrations, since residual unexplained variance in the spectrum can then be attributed to a more complete set of sources of variability.

While metabolite quantification is the primary objective of model fitting and has been the focus of this work, these additional fitting parameters must also be evaluated to confirm that the model learns physically meaningful relationships. Lineshape parameters directly affect peak amplitudes and therefore concentration estimates. Dependencies between predicted lineshape parameters and concentrations should therefore be compared with the corresponding ground-truth relationships or with relationships expected from the signal model. Pearson's $r$ can characterize linear dependencies between parameters, but correlation between predicted outputs alone does not establish that the learned relationship is physically correct, since it may also reflect covariance present in the training distribution or parameter ambiguity. Accuracy and agreement with ground truth should therefore remain the primary evaluation criteria for each predicted fitting parameter, with cross-parameter correlation providing complementary information about the learned dependency structure. Plotting marginal distributions for each parameter identifies regions of the parameter space that are learned well and flags systematic failures that aggregate metrics may not expose.

\section{Uncertainty \& Interpretability}
\label{sec:Interpretability}
Many \DL models remain difficult to interpret because the basis for an individual prediction is not directly observable. In safety-critical applications, such as medicine, this lack of transparency poses a barrier to large-scale clinical adoption of these tools. Interpretability remains an active field in \AI research. The preceding analyses characterize trustworthiness from observed model behavior, including the conditions under which prediction errors increase and the operating ranges within which predefined reliability requirements are satisfied. Clinical translation also benefits from information that can help a user assess the plausibility or uncertainty of an individual output at inference time. In \MRS, explicit modeling fitting and uncertainty quantification provide complementary approaches for increasing this interpretability and supporting the translation of \DL fitting models from research to clinical workflows.

\subsection{Spectral Modeling versus Metabolite Quantification}
The central objective of many model fitting protocols is accurate metabolite quantification. Consequently, many deep learning approaches have focused primarily on predicting metabolite concentrations directly, rather than modeling the full observed spectrum and its constituent signal components. A major drawback of this approach is the difficulty of constructing a spectral fit without the remaining fitting parameters. This prevents clinicians and researchers from visually evaluating the quality of the predicted concentrations. While some disciplines accept outputs from a well-characterized model, spectroscopists still typically insist on comparing spectral fits with the measured data. This allows the analysis of residuals and the calculation of the fit quality number (FQN), which is recommended by the consensus papers \cite{Near2021}. Furthermore, model hallucination is challenging to detect by quantitative metrics, making visual inspection an important safeguard for identifying implausible spectral reconstructions. 

Explicit spectral modeling can improve interpretability by allowing quantitative predictions to be evaluated in the context of the measured spectrum, fitted signal components, and residual structure. Where this form of inspection is required by the intended workflow, this work therefore recommends favoring \DL approaches that support reconstruction of the spectral fit over concentration-only prediction. In cases where models only predict metabolite quantities, it is suggested to use modified linear combination modeling (LCM) techniques to regress the remaining fit parameters. Such a post-hoc fit should be interpreted as a compatibility and residual diagnostic rather than as independent validation of the \DL concentration estimates, since adjustment of the remaining fitting parameters may partially compensate for errors in the predicted concentrations.

\subsection{Uncertainty}
Uncertainty estimation provides a complementary approach to failure detection by quantifying uncertainty associated with individual predictions \cite{Abdar2021}. Unlike the operating-envelope analyses above, which characterize empirical reliability over an evaluation domain, an input-dependent uncertainty estimate can potentially provide a prediction-specific warning signal at inference time. When appropriately validated, this information can therefore be used to support the decision-making process for a user (e.g., how much to trust a given prediction) or to enable optimization of the deployed algorithm (e.g., adversarial training). Although many uncertainty estimation methods have been proposed for deep learning, (1) little is known about their benefits and challenges for various medical applications \cite{Jungo2019}, (2) general acceptance and agreement on methodology are far from established \cite{Psaros2023}, and (3) human-level readability and interpretability remain open issues \cite{Rizzo2022}.

There is widespread agreement on classifying two main types of uncertainty: aleatoric and epistemic uncertainty \cite{Kendall2017}. The irreducible uncertainty in data that gives rise to uncertainty in predictions is called \textit{aleatoric} uncertainty, or data uncertainty. This type is not a property of the model but rather an inherent property of the data distribution. In contrast, \textit{epistemic} uncertainty, or knowledge uncertainty, occurs due to insufficient model knowledge. Contrary to aleatoric uncertainty, epistemic uncertainty can in principle be reduced with additional training data. A toy example that visualizes the effect of uncertainty is given in Fig.~\ref{fig:uncertainty toy example}.

\begin{figure*}
    \centering
    % \includegraphics[width=0.85\textwidth,keepaspectratio=True]{images/ours_01.png}
    \includegraphics[width=0.85\textwidth,keepaspectratio=True]{images/toy_example2.pdf}
    \caption{(Top) A toy example of a data set (training set, blue) is distributed in a 1D domain D = [−4, +4] following a mostly sinusoidal pattern that matches ground truth (yellow, extended to a bigger domain D-GT- defined in [-8,8]). A DL model (not displayed) predicts the sinusoidal wave (red). Good performances are reported within the training data domain D and noise divergent from the sinusoidal way characteristic is seen out of the training data domain. Visualization of aleatoric (bottom-left) and epistemic (bottom-right) uncertainties is reported in shade of red. Despite perfect predictions, a measure of uncertainty relates to the interval [1, 3] where training data suffered a bigger variance (aleatoric). A second measure of uncertainty relates to training distribution, displaying higher uncertainty when predictions are asked for data outside the training domain D (epistemic). Modified from Xu et al. \cite{Xu2021}}
    \label{fig:uncertainty toy example}
\end{figure*}

Various approaches can be used to construct or characterize uncertainty estimates for deep learning models, with different methods providing information about different sources or representations of uncertainty. A schematic representation comparing these estimators with a baseline vanilla architecture is presented in Fig.~\ref{fig:uncertainty estimation types} and is organized in this framework into measured, predicted, learned, calculated, and calibrated uncertainty. A detailed description of each type is provided below.

\paragraph{Measured Uncertainties} (Fig.~\ref{fig:uncertainty estimation types}-B) Measured uncertainties are derived from external evaluations and are not part of the model itself. These uncertainties typically stem from a diverse array of pretrained models (Fig.~\ref{fig:uncertainty estimation types}-B1). They include ensembles of deep learning models, most notably deep ensembles \cite{Lakshminarayanan2017}, in which multiple independently initialized networks are trained and their predictions are aggregated. They can also arise from extensive variations in testing data, such as data ablation, augmentation, test-time augmentation (TTA) \cite{Wang2019}, or bootstrapping, as shown in Fig.~\ref{fig:uncertainty estimation types}-B2. Mathematically, these uncertainties aggregate the centroids and spreads of predictions from point-estimate models. This approach facilitates robust assessment by leveraging model diversity or sensitivity to different data features. However, these uncertainties cannot be estimated during training and thus cannot be leveraged during training.

\paragraph{Predicted Uncertainties} (Fig.~\ref{fig:uncertainty estimation types}-C) Predicted uncertainties are integrated within the model architecture and vary with input data. The model is trained to produce uncertainty-related parameters alongside its primary prediction, which can represent aleatoric uncertainty or total uncertainty (aleatoric and epistemic together), depending on how the model is implemented. In the heteroscedastic formulation \cite{Kendall2017}, the network produces both a mean estimate and a variance estimate, commonly parameterized as $\log \sigma^2$, yielding a parametric predictive distribution whose spread reflects the input-dependent aleatoric uncertainty of the data. Dropout techniques can generally enhance network robustness by randomly deactivating nodes during training. Monte Carlo Dropout (MC Dropout) \cite{Gal2016} extends this principle to stochastic inference by retaining dropout at test time, such that repeated forward passes generate a predictive distribution that approximates epistemic model uncertainty. In this respect, MC Dropout bridges aspects of externally measured prediction variability and model-derived uncertainty estimation.

\paragraph{Learned Uncertainties} (Fig.~\ref{fig:uncertainty estimation types}-D) Learned uncertainty approaches can represent uncertainty by placing probability distributions over model parameters or latent variables rather than relying only on point estimates. In Bayesian Neural Networks (BNNs) \cite{Blundell2015}, network weights are treated as random variables and approximate posterior inference over those parameters provides a probabilistic representation of model predictions.This approach introduces complexity and demands significant computational resources, particularly for implementing variational inference over the parameter posterior during training.

\paragraph{Calculated Uncertainties} (Fig.~\ref{fig:uncertainty estimation types}-E) Calculated uncertainties are derived post-hoc from physics-based or model-based architectures grounded in analytical descriptions of parameter sensitivity or local curvature, including Fisher-information- and Hessian-based approaches. Calculated uncertainties can be obtained post-hoc from approximations to the model's parameter or predictive distribution rather than being emitted directly by a learned uncertainty head. The Laplace approximation \cite{Daxberger2021}, for example, fits a Gaussian around the maximum a posteriori (MAP) estimate using the local curvature of the negative log posterior, with scalable implementations allowing the Hessian to be replaced by structured curvature approximations such as the Fisher information matrix or generalized Gauss-Newton matrix, yielding scalable post-hoc uncertainty estimates without retraining. In this context, the Fisher information describes sensitivity with respect to the neural-network parameters and is therefore distinct from the signal-model Fisher information used in Section~\ref{sec&Visual:subsec} to derive the synthetic oracle \CRLB reference. Both use information or curvature to characterize parameter sensitivity, but they operate on different parameter spaces and serve different inferential purposes. However, these uncertainties may only partially capture intrinsic uncertainties because they primarily describe the uncertainty of the model parameters rather than prediction uncertainty, although this parameter uncertainty can subsequently be propagated to obtain a predictive distribution, as in the Laplace approximation.

\paragraph{Calibrated Uncertainties} Quantile regression \cite{Koenker1978} directly estimates conditional quantiles of the target distribution and can therefore be used to construct input-dependent prediction intervals from estimated lower and upper quantiles. If the conditional quantiles are estimated accurately, these intervals can provide the corresponding nominal conditional coverage, but finite-sample calibration is not guaranteed by quantile regression alone. Quantile models are trained using the asymmetric pinball loss, which independently optimizes each target quantile without imposing parametric constraints on the output distribution, allowing asymmetric and non-Gaussian prediction intervals to emerge naturally. Calibration and sharpness are complementary desiderata for any interval estimator: a calibrated model achieves its declared coverage, while a sharp model produces the narrowest intervals consistent with that coverage, since arbitrarily wide intervals can be perfectly calibrated while conveying no useful predictive information.

\paragraph{Conformal Predictions} Where an explicit coverage guarantee is required, conformal methods provide a complementary calibration framework. Conformal prediction is a distribution-free, post-hoc calibration framework that constructs valid prediction intervals at a user-specified confidence level. It guarantees finite-sample marg\-inal coverage without requiring a parametric model of the data-generating distribution or a particular model architecture, provided the calibration and future observations satisfy the required distributional exchangeability conditions \cite{Angelopoulos2023}. Crucially, conformal prediction should be interpreted as a tool for obtaining calibrated coverage rather than as an intrinsic uncertainty estimator. Large prediction intervals may reflect either genuinely uncertain input regions or samples drawn from outside the calibration distribution. Standard conformal prediction can produce inefficient, overly wide intervals under heteroscedastic noise. Conformalized quantile regression addresses this by combining quantile regression with conformal calibration to yield locally adaptive interval widths \cite{Romano2019}.

\begin{figure*}
    \centering
    \includegraphics[width=\textwidth,keepaspectratio=True]{images/ours_02.png}
    \caption{Diagram representing various estimators for uncertainties in DL. (A) Vanilla DL network, alpha input, Beta output and $\theta_i$ are different deterministic network nodes. (B) Measured Uncertainties. B.1 Multiple pre-trained models either on similar or different dataset are tested for the same testset alpha and produce different predictions Beta. B.2 The same pretrained network M1 is tested on different test sets producing different output Beta. For both cases, the evaluation of the single model predicted output as an aggregate gives a measure of uncertainty reflected in the estimate’s distribution. (C) Predicted uncertainties include specific nodes in the network that are dedicated to embed features reflected as output into a measure of the output variability (sigma). (D) Learned uncertainties shift the network paradox introducing stochastic nodes (rho) instead of deterministic ones (Theta). A stochastic node activation varies according to a distribution. The output will inherit stochastic properties from the network. (E) Calculated uncertainties relies on the re-definition of part of the network by encoding a specific  physic-driven set of nodes (pi). The mathematical model definition defines a proxy to Fisher information calculation and generates inference to CRLBs.}
    \label{fig:uncertainty estimation types}
\end{figure*}

Each of these uncertainty measures provides valuable insights, enabling researchers to better understand and improve the reliability of deep learning models across various scientific fields. Regardless of the chosen method, uncertainty measures need to be both calibrated and validated. This is important in \DL applications as a whole, but especially for predictive uncertainties, which can exhibit overconfident or otherwise unreliable behavior under conditions including distribution shift and non-stationarity \cite{Meronen2020, Meronen2021}, with ReLU-based networks additionally exhibiting documented confidence and uncertainty pathologies outside the training distribution \cite{Hein2019a, Ulmer2021}. \DL models should accurately model their predictive confidence for a given estimate and identify when a sample is likely out-of-distribution (\OOD). Their evaluation should therefore establish whether the reported uncertainty is associated with prediction error, whether its probabilistic interpretation is calibrated, and whether uncertainty responds appropriately to \OOD inputs and other departures from the characterized data distribution. Samples identified as potentially \OOD can then be flagged for additional review or alternative fitting using established methods.

\subsection{Predictive Uncertainty Analysis}
Producing an uncertainty estimate does not itself establish that the model knows when its predictions are unreliable. Beyond estimating predictive uncertainty, it is equally important to evaluate the quality and reliability of the resulting uncertainty estimates. Predictive uncertainty analysis seeks to characterize the relationship between predicted uncertainties and model performance, determine whether confidence estimates are well-calibrated, and assess the predictive distributions generated by stochastic inference methods. The following sections describe complementary approaches for validating the relationship between uncertainty and error, evaluating calibration and sensitivity to distribution shift, and characterizing Monte Carlo prediction distributions.

\subsubsection{Validation: Correlation Between Error and Uncertainty, \textit{cobeau}}
When probabilistic models are used to predict uncertainty estimates, the uncertainty predictions must be validated. The first step is to identify what these uncertainty measures are meant to represent. In some deep learning applications, multiple forward passes can be used to create a distribution of predictions for each sample, and the uncertainty is calculated as the standard deviation of the predictions. In contrast, other models predict uncertainty parameters directly alongside the primary model output. In both cases, the resulting uncertainty estimate should be empirically validated, while uncertainty estimates intended to have a probabilistic interpretation additionally require calibration. Depending on the implementation, these values can represent aleatoric uncertainty, epistemic uncertainty, or both.

Rost \cite{Rost2019} introduced $cobeau$, which measures the linearity of the mapping between the predicted uncertainty and any relevant calculated error or uncertainty metric. Its definition, given in Eqn.~\ref{eqn:cobeau}, shows that $cobeau$ is Pearson’s correlation coefficient between the predicted uncertainty and a relevant nonnegative error metric, such as absolute, squared, or root squared error. A positive $cobeau$ therefore indicates that larger predicted uncertainties tend to occur alongside larger observed errors, providing evidence that the uncertainty estimate is informative about relative prediction reliability. It does not by itself establish probabilistic calibration. 
\begin{equation}\label{eqn:cobeau}
    cobeau = \frac{\sum_{i=1}^I(\epsilon_i - \bar{\epsilon})(\sigma_i-\bar{\sigma})}{\sqrt{\sum_{i=1}^I(\epsilon_i - \bar{\epsilon})^2}\sqrt{\sum_{i=1}^I(\sigma_i - \bar{\sigma)^2}}}
\end{equation}
\noindent where $\bar{\epsilon}$ and $\bar{\sigma}$ denote the means of the error metric $\epsilon$ and predicted uncertainty $\sigma$, respectively, and $I$ is the number of observations. This can be calculated individually or collectively for multivariate outputs. Section~\ref{subsubsec:Lin's rho} proposed replacing the mean and variance with the median and nonparametric variance to more accurately accommodate non-normal distributions.

As with the standard Pearson’s $r$, this metric is bounded to the range $[-1,+1]$, with $-1$ indicating a perfect negative linear correlation, $0$ indicating no linear correlation, and $+1$ indicating a perfect positive linear correlation. The average composite value can be used to compare different implementations of uncertainty modeling for a given task. Because this is based on Pearson’s $r$, its statistical significance can be assessed using the same inferential framework described in Section~\ref{sec:perf:subsec:awareness:subsubsed:p-values}. For interpretation, however, the magnitude of \textit{cobeau} and its confidence interval should remain the primary measures of the strength and precision of the uncertainty-error association.

The second validation step is to assess the inter-variable dependence structure of the uncertainty estimates. This can be studied through a correlation analysis, as described above, using a heatmap of the correlation matrix between the predicted uncertainties and the percent errors. Similar to analysis for metabolite concentrations, strong unexplained correlations between otherwise unrelated metabolites are undesirable, whereas correlations may be expected between metabolites with significant spectral overlap as well as with a metabolite's lineshape parameters. If additional correlations exist between one component and a group of unrelated components, this could be indicative of an ill-posed problem, as described above, or they could identify or confirm the existence of important correlations between spectral components. This is a helpful tool to further one’s understanding of the problem.

\paragraph{Multivariate Extensions}
\textit{Cobeau} is built upon the Pearson correlation coefficient, which means that it inherits the same strengths and limitations of $r$. While there is no direct multivariate form of the statistic, the following alternatives can be used to formulate aggregate, model-level summaries. First, as described in Section~\ref{sec:Eval:subsec:multivariate:subsubsec}, values can be calculated per-metabolite and then averaged across the metabolites for an aggregate summary statistic. However, for $r$-derived values specifically, this work recommends averaging via the Fisher transformation \cite{Fisher1921}. For metabolite $m$,
\begin{equation}
    z_m = \textnormal{arctanh}(cobeau_m) = \frac{1}{2} \ln \left(\frac{1+cobeau_m}{1-cobeau_m}\right).
\end{equation}
The transformed values can then be combined across the $M$ metabolites and subsequently transformed back to the correlation scale. When the number of observations differs between metabolites, the Fisher-transformed coefficients can be weighted according to their approximate inverse sampling variance. Because $\textnormal{Var}(z_m)\approx 1/(N_m-3)$, the conventional weight is $w_m=N_m-3$, giving:
\begin{equation}
    \\overline{z} = \frac{\sum_{m=1}^{M} w_m z_m}{\sum_{m=1}^{M} w_m}, \qquad w_m=N_m-3.
\end{equation}
\begin{equation}
    cobeau_{\mathrm{Fisher}} = \tanh\left(\overline{z}\right).
\end{equation}
When all metabolites are evaluated using the same number of observations, the weights are identical and $\overline{z}$ reduces to the arithmetic mean of the Fisher-transformed coefficients. In the present application, however, the metabolite-specific coefficients are generally calculated from the same spectra and are therefore statistically dependent. The conventional Fisher weights can account for differences in metabolite-specific sample size when forming the aggregate point estimate, but they do not account for this dependence. Confidence intervals for the aggregate statistic should therefore preferably be obtained by jointly resampling spectra, or subjects for \textit{in vivo} data, so that the dependence between metabolite-specific coefficients is preserved.

This provides a compact model-level summary while retaining the metabolite-specific \textit{cobeau} values as the primary interpretable quantities. Because the metabolite-specific coefficients are generally calculated from the same spectra and are therefore statistically dependent, uncertainty in the aggregate statistic should preferably be estimated by jointly resampling spectra, or subjects for \textit{in vivo} data, rather than treating the individual coefficients as independent observations.

An alternative pooled representation can be obtained by standardizing the errors and uncertainties separately per-metabolite and then calculating the \textit{cobeau} statistic:
\begin{equation}\label{eqn:standardization for cobeau}
    \epsilon^\prime_{n,m} = \frac{\epsilon_{n,m} - \mu_{\epsilon,m}}{s_{\epsilon,m}}\textnormal{, and }
    \sigma^\prime_{n,m} = \frac{\sigma_{n,m} - \mu_{\sigma,m}}{s_{\sigma,m}}
\end{equation}
The standardized values can subsequently be pooled across metabolites and the Pearson correlation calculated between $\epsilon^\prime$ and $\sigma^\prime$. This suppresses inter-metabolite differences in the location and scale of the error and uncertainty distributions, so that the resulting statistic summarizes whether relatively large uncertainty values are associated with relatively large errors across the joint metabolite set. However, when the same number of observations is available for every metabolite, this pooled correlation is algebraically equivalent to the arithmetic mean of the metabolite-specific \textit{cobeau} coefficients. It should therefore be interpreted as an alternative pooled representation rather than an independent multivariate association measure.

A complementary analysis can determine whether the relationship between uncertainty and error persists after accounting for differences in intrinsic estimation difficulty. For metabolite $m$, \textit{cobeau} can be evaluated separately within \mSNR or oracle-information strata $G_{m,g}$:
\begin{equation}
    cobeau_{m,g} = \textnormal{Corr} \left( \epsilon_m, \sigma_m \mid G_{m,g} \right).
\end{equation}
The unstratified coefficient characterizes whether uncertainty tracks prediction error across the full characterization domain, whereas the stratified coefficient evaluates whether the same relationship remains present among spectra of comparable estimation difficulty. This distinction can help determine whether the uncertainty estimate provides prediction-specific information beyond its shared dependence with error on factors such as \SNR or intrinsic spectral identifiability.

\subsubsection{Calibration and \OOD Sensitivity: Predictive Entropy \& Predictive Marginal Variance}
An important function of uncertainty metrics in \DL is to identify when a model does not know. The previous section detailed how to validate the relevance of predicted uncertainties. The next step is to assess uncertainty calibration and its sensitivity to distribution shift to evaluate the model’s ability to identify its own uncertainty. These properties are complementary: calibration characterizes whether the magnitude of the reported uncertainty is consistent with observed prediction behavior, while \OOD sensitivity characterizes whether uncertainty increases appropriately when inputs depart from the distribution on which the model was characterized. Meronen et al. \cite{Meronen2021} showed that the predictive marginal variance is a useful metric for detecting out-of-distribution (\OOD) data in their experiments, where it provided stronger \OOD discrimination than predictive entropy. When comparing in-distribution (\ID) and \OOD datasets, the median of the predicted variances for all outputs of a sample can be displayed in a histogram. Fig.~\ref{fig:predictive marginal variance} shows that \ID data should generally produce low predictive marginal variance, while \OOD data, for a well-calibrated uncertainty model that is appropriately sensitive to distribution shift, should exhibit increased predictive variance and become more readily distinguishable. Although Meronen et al. \cite{Meronen2021} demonstrated these measures in a classification setting, predictive marginal variance extends naturally to regression. For stochastic prediction methods, let $\hat{y}_{n,m}^{(t)}$ denote stochastic prediction $t=1,\ldots,T$ for metabolite $m$ and spectrum $n$. The metabolite-specific predictive marginal variance can be estimated as:
\begin{equation}
    v_{n,m} = \frac{1}{T-1} \sum_{t=1}^{T} \left(\hat{y}{n,m}^{(t)} - \overline{\hat{y}}{n,m}\right)^2.
\end{equation}
For heteroscedastic models that additionally predict an aleatoric variance $s_{n,m}^{2(t)}$ during each stochastic forward pass, the total predictive marginal variance can instead be estimated using the law of total variance:
\begin{equation}
    v_{n,m}^{\mathrm{total}} = \frac{1}{T}\sum_{t=1}^{T} s_{n,m}^{2(t)} + \frac{1}{T-1}\sum_{t=1}^{T} \left(\mu_{n,m}^{(t)} - \overline{\mu}_{n,m}\right)^2.
\end{equation}
The first term represents the expected predicted aleatoric variance, while the second represents variability between stochastic model predictions. Because raw predictive variances are not directly comparable across metabolites with different concentration scales, a model-level \OOD score can be obtained after normalization by a fixed metabolite-specific scale $a_m$ determined from the training or development data:
\begin{equation}
    v_{n,m}^{\mathrm{std}} = \frac{v_{n,m}^{\mathrm{total}}}{a_m^2}.
\end{equation}

A robust spectrum-level summary analogous to the aggregation used by Meronen et al. can then be defined as:
\begin{equation}
    U_{n}^{\mathrm{var}} = \textnormal{median}{m} \left(v{n,m}^{\mathrm{std}}\right).
\end{equation}
The fixed scale $a_m$ may, for example, correspond to the metabolite-specific standard deviation in an independent development set. Using a fixed reference prevents the normalization itself from changing with the uncertainty or predicted concentration of the test sample.

\begin{figure*}[h!]
    \centering
    % \includegraphics{}
    \begin{subfigure}{0.32\textwidth}
        \includegraphics[width=\textwidth,keepaspectratio]{images/Meronen_01.png}
    \end{subfigure}
    \begin{subfigure}{0.32\textwidth}
        \includegraphics[width=\textwidth,keepaspectratio]{images/Meronen_02.png}
    \end{subfigure}
    \begin{subfigure}{0.32\textwidth}
        \includegraphics[width=\textwidth,keepaspectratio]{images/Meronen_03.png}
    \end{subfigure}
    \caption{Histograms of sample-level variances of different datasets. From left to right, these plots show the OOD detection capability of models with increasingly well-calibrated uncertainty predictions. Here, green is the in-distribution dataset, blue is in the middle of ID and OOD, and red is strongly OOD. Adapted from Meronen et al. \cite{Meronen2021}.}
    \label{fig:predictive marginal variance}
\end{figure*}

For models that define or permit estimation of a predictive probability density, predictive entropy provides a complementary uncertainty measure (Fig.~\ref{fig:predictive entropy}), although the differences between the datasets may be less pronounced than those observed with the predictive marginal variance. In the regression setting, this requires the entropy of a continuous predictive distribution rather than the categorical predictive entropy used in the original classification implementation.

For an approximately Gaussian predictive distribution for metabolite $m$, a standardized differential-entropy score can be calculated from the standardized predictive variance:
\begin{equation}
    H_{n,m}^{\mathrm{std}} = \frac{1}{2}\log\left(2\pi e,v_{n,m}^{\mathrm{std}}\right).
\end{equation}
For a multivariate Gaussian predictive distribution with covariance matrix $\widehat{\boldsymbol{\Sigma}}_n$, the corresponding joint entropy can be calculated after metabolite-wise scale standardization. Let $\mathbf{A}=\textnormal{diag}(a_1,\ldots,a_M)$ and
\begin{equation}
    \widehat{\boldsymbol{\Sigma}}_n^{\mathrm{std}} = \mathbf{A}^{-1} \widehat{\boldsymbol{\Sigma}}_n \mathbf{A}^{-1}.
\end{equation}
The standardized joint differential entropy is then:
\begin{equation}
    H_n^{\mathrm{std}} = \frac{1}{2}\log\left[(2\pi e)^M \det\left(\widehat{\boldsymbol{\Sigma}}_n^{\mathrm{std}} \right)\right].
\end{equation}
When the predictive distribution is substantially non-Gaussian, its differential entropy may instead be estimated from the stochastic prediction samples, although this introduces an additional density-estimation problem and should therefore be validated separately.

\begin{figure*}[ht!]
    \centering
    % \includegraphics{}
    \begin{subfigure}{0.32\textwidth}
        \includegraphics[width=\textwidth,keepaspectratio]{images/Meronen_04.png}
    \end{subfigure}
    \begin{subfigure}{0.32\textwidth}
        \includegraphics[width=\textwidth,keepaspectratio]{images/Meronen_05.png}
    \end{subfigure}
    \begin{subfigure}{0.32\textwidth}
        \includegraphics[width=\textwidth,keepaspectratio]{images/Meronen_06.png}
    \end{subfigure}
    \caption{Histograms of the entropies of the mean predicted values. The models correspond to those in Fig. \ref{fig:predictive marginal variance} above. Here, green is the in-distribution dataset, blue is in the middle of ID and OOD, and red is strongly OOD. Adapted from Meronen et al.~\cite{Meronen2021}.}
    \label{fig:predictive entropy}
\end{figure*}

For controlled evaluation datasets with known \ID and \OOD status, the discriminatory value of $U_n^{\mathrm{var}}$, $H_n^{\mathrm{std}}$, or another prespecified uncertainty score can additionally be quantified by varying an \OOD decision threshold and reporting the area under the receiver operating characteristic curve and precision-recall curve, following the evaluation strategy of Meronen et al.~\cite{Meronen2021}. These evaluations are helpful to confirm that low variance estimates correspond well to \ID data and that predictive uncertainty increases appropriately under controlled distribution shifts. Given an uncertainty model that is well-calibrated on \ID data and appropriately sensitive to distribution shift, this will also show whether the model assigns systematically greater uncertainty to different types of \textit{in vivo} data that depart from its characterized distribution. The resulting uncertainty scores can therefore complement the preceding error and operating-envelope analyses by identifying inputs that warrant additional scrutiny and further characterizing the boundaries of the model's reliable operating domain.

\subsubsection{Monte Carlo Prediction Distributions}
Monte Carlo (MC) prediction distributions are useful for characterizing uncertainty, but they should not be confused with the distribution of the dataset \cite{Rubinstein2016}. A dataset distribution describes the variability of parameters across a population of samples, whereas an MC prediction distribution describes the variability of repeated predictions for a single sample under a specific source of randomness.

Historically, Monte Carlo studies have been useful for the evaluation of MRS fitting and quantification algorithms because ground truth metabolite concentrations are typically not available \textit{in vivo}. Instead, synthetic spectra with known ground truth values were repeatedly perturbed with noise realizations to characterize estimator precision and uncertainty \cite{Instrella2024, Landheer2021b}. This approach remains valuable when the performance of a model must be evaluated against a known ground truth.

In \DL model fitting, MC prediction distributions arise from repeated stochastic evaluations of the same spectrum, with respect to a source of randomness. The source of randomness may be generated from techniques such as MC dropout, variational inference, an ensemble of methods, or repeated perturbations of the input data. These distributions characterize prediction variability associated with the specified source of randomness and can therefore be used to characterize corresponding components of predictive uncertainty rather than population variability and therefore answer fundamentally different questions than dataset-level analyses. The interpretation of such distributions therefore depends on both the mechanism used to generate the samples and the number of MC samples used. As the number of independent, or sufficiently weakly dependent, MC samples increases, empirical estimates of distributional quantities such as the mean, variance, and quantiles become increasingly stable and, under the relevant regularity conditions, converge toward the corresponding quantities of the predictive distribution.

Practical implementations therefore require consideration of convergence and the number of MC samples needed to obtain stable uncertainty estimates. Alves et al. provide a practical methodology for selecting the number of Monte Carlo realizations required to achieve a prespecified confidence level and confidence-interval precision for an MC-averaged concentration estimate and demonstrate that the required number of simulations depends on factors such as metabolite concentration and SNR \cite{Alves2024}. Their approach derives the required number of realizations from the sampling variability of the Monte Carlo mean and a prespecified maximum acceptable deviation. This provides a useful general principle for \DL uncertainty analysis: the number of stochastic forward passes should be selected according to the precision required for the quantity being estimated rather than fixed arbitrarily. For uncertainty quantities other than the predictive mean, such as predictive variance, quantiles, or entropy, the convergence criterion should be defined for that specific statistic because its sampling variability is not described by the confidence interval for the mean used by Alves et al.

\subsection{Decomposed Uncertainties}
\subsubsection{Validating the Decomposition}
With predicted values from deep learning, it is especially important to confirm expected model behavior. When predicting aleatoric and epistemic uncertainty, or decomposing uncertainty into those components, certain patterns are expected. Consider a plot of aleatoric uncertainty versus prediction error. One would expect that low prediction errors correspond to low aleatoric uncertainty. As Ghoshal et al. \cite{Ghoshal2021} did in Fig.~\ref{fig:decomposed uncertainties}, the data markers can be color-coded according to the other type of uncertainty. In the case above, the epistemic uncertainty should not exhibit any particular patterns when shown in the aleatoric plot. Similarly, when plotting the epistemic uncertainty against the prediction errors, the example shown by Ghoshal et al. exhibits no clear direct relationship, consistent with the epistemic component capturing a different source of uncertainty from the aleatoric error-related structure. However, when coloring the data markers according to their aleatoric uncertainty, a pattern emerges: while the epistemic uncertainty forms a blob-like pattern, the aleatoric uncertainty diffuses left to right from low to high, which should be the same pattern observed when plotting the aleatoric uncertainty versus prediction error in the first plot. The primary objective here is to confirm that the model produces the expected behavior.

This diagnostic should be interpreted with respect to the uncertainty formulation being evaluated rather than as a requirement that epistemic uncertainty remain constant across samples. Epistemic uncertainty originates from uncertainty in the learned model or its parameters, but its prediction-specific magnitude can vary with the input, for example when training support changes or a sample lies outside the characterized distribution.

\begin{figure*}[bh]
    \centering
    \includegraphics[width=0.8\textwidth,keepaspectratio]{images/Goshal_01.png}
    \caption{Expected patterns of aleatoric and epistemic uncertainty. Accuracy increases as aleatoric uncertainty decreases (shown on the left) while the epistemic uncertainty has no discernible pattern. On the right, the epistemic uncertainty still has no pattern, but the aleatoric uncertainty decreases from left to right as the prediction accuracy increases. Adapted from Ghoshal et al. \cite{Ghoshal2021}.}
    \label{fig:decomposed uncertainties}
\end{figure*}

\subsubsection{Characterizing the Uncertainty of Spectral Components}
The underlying model used for model fitting, in traditional and deep learning-based approaches, can be further studied using aleatoric and epistemic uncertainties. Figure~\ref{fig:component-level uncertainty} shows how model confidence can be characterized across the parameter space for each spectral component. In this case, the uncertainties can be plotted against the spectral components’ values. If the parameter ranges are different for the training and testing datasets, the ranges for the training dataset can be indicated in the plots, as in Fig.~\ref{fig:component-level uncertainty}, which helps study how the model performs for parameter values outside of the training ranges. When considering components such as \SNR and concentration, it would be more concise to use a single plot of uncertainty versus \mSNR. When characterizing a model for deployment, it would be useful to separate the data by metabolite. Another application of these plots is to study how various spectral components, such as the linewidths, affect the uncertainties for metabolite quantification. For this approach, the metabolite uncertainties should be plotted against the spectral components’ parameter values.

\begin{figure*}[!b]
    \centering
    \includegraphics[width=0.8\textwidth,keepaspectratio]{images/Lee_01.png}
    \caption{Breakdown of the aleatoric, epistemic, and total uncertainy for individual predicted variables across their parameter spaces highlighting the heteroscedasticity. Adapted from Lee and Kim~\cite{Lee2022}}
    \label{fig:component-level uncertainty}
\end{figure*}

\subsection{Reliability and the Cram{\'e}r--Rao Lower Bounds}
\label{sec:Laying:subsec:Reliability} % revised JTL 13.12.24
Incorporating uncertainty quantification into \DL models for metabolite quantification and model fitting is important. Ideally, uncertainty predictions should be decomposed into their aleatoric and epistemic components\footnote{Defined in the appendix.}, providing insights into model- and sample-level limitations. This is particularly relevant in medical applications, where accurate characterization and calibration of \AI models are needed. While research into uncertainty prediction methodologies is ongoing, significant advancements have been made in recent years. Future work should aim to integrate these advancements into \MRS, particularly as \DL applications continue to expand in the field. An overview of relevant trends is provided in Section~\ref{sec:Interpretability}.

Calculating the Cram\'{e}r--Rao Lower Bounds (CRLBs) for \DL-based model fitting results is not recommended. Within the \MRS community, CRLBs are traditionally used in conjunction with mean-based metrics as a proxy for the standard deviation of the error \cite{Landheer2021b}. This practice stems from the logistical challenges of obtaining multiple experimental repetitions because typical \MRS scan durations are constrained by clinical time. In traditional model fitting, the classical \CRLB provides a lower bound on estimator covariance under the specified statistical signal model, together with the relevant regularity conditions and, in its conventional form, an unbiased estimator. It is derived from the inverse of the Fisher information matrix for that model. The diagonal elements of the resulting covariance bound correspond to lower bounds on the parameter variances, and their square roots therefore provide the corresponding standard-deviation bounds conventionally reported in \MRS. However, this classical interpretation is not automatically applicable to learned \DL estimators, which may violate the unbiasedness or other estimator assumptions required by the conventional \CRLB.

This standard-deviation interpretation also means that the \CRLB-derived reference should be compared with an empirical metric that characterizes the same type of dispersion. As discussed in Section~\ref{sec:eval:subsec:error:subsubsection}, standard deviation and the median-based robust dispersion measures adopted in this work characterize variability using different estimands. A conventional \CRLB-derived standard-deviation reference should therefore not be treated as a direct proxy for a median-based or other robust dispersion metric. The reference quantity and empirical metric should therefore be selected and interpreted according to the statistical properties they measure.

A more fundamental consideration for \DL-based estimators is estimator bias. The conventional \CRLB assumes an unbiased estimator, whereas learned estimators may introduce bias through constraints, prior information, or other features of the estimation procedure. Such estimators may also incorporate regularization terms to reduce variance and prevent parameter inflation. For \DL estimators, additional bias may arise from the learned mapping, training distribution, architectural constraints, or learned prior information. Consequently, the classical unbiased \CRLB should not be applied directly as a lower variance bound for such estimators when these sources of bias are relevant \cite{Fessler1996}. Generalized \CRLB formulations for biased estimators exist, but require characterization of the estimator bias and therefore do not provide the same simple estimator-independent reference as the conventional unbiased form.

Furthermore, the calculation of CRLBs depends on the mathematical formulation of the fitting model. For physics-derived or model-driven \DL approaches, a \CRLB may be calculated from the embedded forward signal model when that model defines an appropriate likelihood. Such a calculation characterizes the information available in the modeled signal but does not automatically characterize the uncertainty of the learned estimator, since it does not incorporate the complete training procedure, learned parameter mapping, training distribution, or stochastic behavior of the network. This distinction is important because Fisher information can also be defined for neural-network parameters or predictive models, but such a Fisher information matrix is a different statistical object from the signal-model Fisher information used to derive conventional \MRS CRLBs. Additionally, assumptions regarding the noise distribution in NMR signals must be considered more rigorously in \DL-based post-processing contexts. The likelihood and noise model used to construct the Fisher information should reflect the relevant data-generating process. White Gaussian noise provides a convenient and common formulation in \MRS, but is not a general requirement of \CRLB theory.

This distinction does not preclude the use of CRLBs as external information-based references. As described in Section~\ref{sec:Analy&Visual:subsec:infoStand}, an oracle CRLB derived from the known synthetic generative model can be used to characterize the intrinsic information available under a controlled simulated condition. In that role, the \CRLB is not interpreted as the uncertainty of the \DL estimator or as a strict bound on its empirical variance. Instead, it provides a common reference against which model-specific bias and variability can be characterized.

\section{Method Comparison Studies}
\label{sec:Reg:subsec:comparison}
The preceding analyses characterize predictive performance, identify failure regimes, and establish the conditions under which model outputs can be considered reliable. Method comparison addresses the subsequent translational question of whether those measurements agree sufficiently with an established reference for the intended application. Method comparison analysis evaluates whether the outputs of a model or method are in sufficient agreement with those of an established reference to be considered interchangeable for a given clinical or scientific purpose. This is distinct from Sec.~\ref{sec:eval:subsec:benchmarking} which presents model-level summaries to characterize aggregate predictive performance across the full metabolite set. Method comparisons characterize the \textit{structure} of disagreement on a per-metabolite basis, decomposing it into systematic bias, proportional bias, and residual variability as a function of concentration. The methodology is organized around two comparison regimes with structurally different statistical requirements: an \textit{in silico} regime in which simulated ground truth values are available, and an \textit{in vivo} regime in which a conventional parametric fitting algorithm serves as the reference. Each regime is subject to distinct assumptions regarding measurement error structure, sample independence, and the appropriate choice of reference axis, as detailed in the following sections.

\subsection{Study Design and Error Structure}
\label{sec:Reg:subsec:comparison:subsubsec:design}
The statistical behavior of method comparison procedures depends critically on the relative measurement error of the two methods under comparison. In the \textit{in silico} regime, the reference consists of the known simulation parameters, which are exact by construction and carry negligible error. When one method is effectively error-free, the mean of the two methods is dominated by the error-bearing method. Plotting the differences against the mean introduces a spurious correlation between the difference and the average that inflates or deflates the observed agreement~\cite{Krouwer2008}. In this setting, differences must be plotted against the reference value directly, as proposed by Krouwer~\cite{Krouwer2008}, or regression-based limits of agreement must be used. 

In the \textit{in vivo} regime, both the \DL model and the conventional fitting algorithm carry non-negligible measurement uncertainty. Here, ordinary least squares (OLS) regression is inappropriate because its assumption of an error-free predictor is violated, leading to attenuation of the estimated slope toward zero~\cite{Riu1997}. An errors-in-variables framework is therefore required, as discussed in Section~\ref{sec:Reg:subsec:comparison:subsubsec:regression}.

Consistent with the framework established in Section~\ref{sec:eval:subsec:errorMetrics}, normality of the differences should not be assumed for \DL model outputs. Robust and nonparametric estimators are therefore applied throughout. With respect to the independence of observations, the two regimes differ. In the \textit{in silico} regime, simulation parameters are sampled i.i.d. by construction, so each spectrum constitutes a fully independent observation. In the \textit{in vivo} regime, multiple spectra acquired from the same subject share physiological and acquisition-related sources of variation, violating the independence assumption. This dependence must therefore be incorporated when estimating confidence intervals for bias and limits of agreement, as discussed in Section~\ref{sec:Reg:subsec:comparison:subsubsec:samplesize}.

\subsection{Test Set Stratification Across the Patho-Physiological Parameter Space}
\label{sec:Reg:subsec:comparison:subsubsec:stratification}
The preceding characterization framework identifies whether model reliability changes across the parameter space. Method comparison must additionally determine whether agreement with the reference remains acceptable within clinically relevant physiological and pathological regimes rather than only across the pooled evaluation set. Three test sets should therefore be employed: 
\begin{itemize}
    \item A \textit{physiological} test set which restricts concentrations and relaxation values to the normal population range for each metabolite.
    \item A \textit{pathological} test set that covers clinically relevant extremes of the concentration range, including, for example, both elevated and reduced concentrations, as observed in neoplastic or inflammatory pathologies and in degenerative or metabolic disorders.
    \item A \textit{combined} test set pooling both subsets and spanning the full concentration and relaxation rang\-es, enabling regression-based characterization of the concentration-dependent agreement structure.
\end{itemize}

Regression-based limits of agreement should be estimated on the combined set. Stratified results are then examined for consistency between the physiological and pathological subsets. A concentration-dependent breakdown in agreement within a stratum constitutes a specific, clinically actionable failure mode distinct from global model failure, and should be reported separately.

In the \textit{in silico} setting, stratification is operationally straightforward because the simulation parameters are controlled. Physiological and pathological subsets can be defined prospectively by partitioning the sampled parameter ranges. In the \textit{in vivo} setting, sufficient numbers of spectra from pathological cases must be obtained to support stratified inference, which is a practical constraint that may limit the resolution of the pathological evaluation and should be acknowledged accordingly.

\subsection{Bland-Altman (BA) Agreement Analysis}
\label{sec:Reg:subsec:comparison:subsubsec:ba}
The BA framework~\cite{Bland1986,Mansournia2021} characterizes the distribution of pairwise differences $d_n = x_n - y_n$ between the \DL model output $x_n$ and the reference $y_n$ for each spectrum $n$. This provides estimates of systematic bias and the interval within which a specified proportion of individual differences are expected to lie. Under the sign convention adopted in Section~\ref{sec:eval:subsec:error:subsubsection:errors}, $d_n < 0$ indicates underestimation by the \DL model relative to the reference and $d_n > 0$ indicates overestimation. The analysis is performed per metabolite and per test set stratum.

\subsubsection{Choice of Reference Axis}
As established in Section~\ref{sec:Reg:subsec:comparison:subsubsec:design}, the appropriate x-axis for the BA plot depends on the error structure of the comparison. In the \textit{in silico} regime, differences are plotted against the reference (ground truth) values directly, following the Krouwer modification~\cite{Krouwer2008}. In the \textit{in vivo} regime, differences are plotted against the mean of the two methods, consistent with the standard formulation of Bland and Altman~\cite{Bland1986}. These are not interchangeable choices: applying the mean-axis convention when the reference is error-free introduces a mathematical dependency between the plotted quantities that induces apparent proportional bias where none actually exists~\cite{Cesana2024}.

\subsubsection{Regression-Based Limits of Agreement}
Given the assumed heteroscedastic error structure described in Section~\ref{sec:Reg:subsec:comparison:subsubsec:design}, the standard constant-variance limits of agreement $\overline{d} \pm 1.96\,\hat{\sigma}_d$ are not appropriate. Instead, both the bias and the standard deviation of differences are modeled as functions of the reference values $c$, yielding value-dependent limits of agreement, $\widehat{\mathrm{LoA}}$:
\begin{equation}\label{eqn:ba_loa_heteroscedastic}
    \widehat{\mathrm{LoA}}(c) = \hat{\beta}(c) \pm 1.96\,\hat{\sigma}_d(c)
\end{equation}
\noindent where $\hat{\beta}(c)$ is the estimated bias function and $\hat{\sigma}_d(c)$ is the estimated standard deviation of differences, each modeled by regression on $c$. The factor 1.96 corresponds to the central 95\% limits under the assumption of conditional normality of the differences. This formulation also follows the regression-based extension of the BA method and its connection to tolerance intervals described by Francq and Govaerts~\cite{Francq2016}. In practice, the functional form of $\hat{\beta}(c)$ and $\hat{\sigma}_d(c)$ is determined empirically. A linear model in $c$ is used as the default, with residual diagnostics used to assess adequacy. Confidence bands on both functions are reported, providing a complete characterization of uncertainty in the estimated limits across the concentration ranges.

\subsubsection{Robust Estimation of Bias.}
Consistent with the nonparametric framework applied throughout this work (Sections~\ref{sec:eval:subsec:error:subsubsection:centroids} and~\ref{sec:eval:subsec:error:subsubsection:dispersion}), the bias is estimated using the median difference $\widetilde{d}$ in place of the mean difference $\overline{d}$, especially when the distribution of $d_n$ is non-normal or exhibits heavy tails. The dispersion of differences is correspondingly estimated using the nonparametric standard deviation $\widetilde{\sigma}_d$ defined in Section~\ref{sec:eval:subsec:error:subsubsection:dispersion}, and the IQR of the differences is reported as a supplementary robust characterization of spread. Where results are reported under both parametric and nonparametric estimators, notable discrepancies between them are diagnostic of outlier influence on the parametric estimates, and the robust results are treated as primary.

\subsubsection{Confidence Intervals and Clinical Acceptability}
Confidence intervals are reported on both the estimated bias and the limits of agreement at each concentration, reflecting the uncertainty arising from finite sample size~\cite{Lu2016}. The precision of these intervals is a direct function of sample size. See Section~\ref{sec:Reg:subsec:comparison:subsubsec:samplesize} for sample size considerations.

The BA analysis is descriptive by design: it characterizes the magnitude and concentration-dependence of disagreement but does not itself determine whether that disagreement is clinically acceptable~\cite{Silveira2024}. Acceptable limits of agreement must therefore be specified \textit{a priori}, grounded in clinical or biological criteria for each metabolite. For example, the known inter-subject biological variability of the metabolite, or the minimum concentration difference that would alter a clinical decision. Results are reported as the estimated limits of agreement with their confidence bands alongside the pre-specified acceptable limits, enabling a direct assessment of whether the \DL model meets the required standard. The BA analysis is complementary to $\rho_C$ (Section~\ref{subsubsec:Lin's rho}): where $\rho_C$ provides a single scalar summary of accuracy and precision, the BA plot provides a distributional and concentration-dependent characterization of the structure of disagreement that a scalar summary cannot illuminate.

\begingroup
\renewcommand{\arraystretch}{1.5}
\begin{table*}[h]
    \centering
    \begin{tabular}{|
    >{\raggedright\arraybackslash}p{89.3601pt}|| % Column 1: Left
    >{\centering\arraybackslash}p{110.1879pt}|  % Column 2: Center
    >{\centering\arraybackslash}p{112.1489pt}|  % Column 3: Center
    >{\centering\arraybackslash}p{132.08827pt}| % Column 4: Right
}\hline
         & Test 1: Accuracy & Test 2: Precision & Test 3: Agreement \\\hline\hline
        Structural H\textsubscript{0,[1,2,3]}
            & $\mathcal{E}[X]=\mathcal{E}[Y]$
            & $\lambda=\text{V}[\epsilon]/\text{V}[\delta]=1$
            & \makecell[c]{$X=Y$, or equivalently \\ $X-Y=\alpha+(\beta-1)Y$ \\ given $\alpha=0$ and $\beta=1$} \\\arrayrulecolor{lightgray}\hline\
        Functional H\textsubscript{0,[1,2,3]}
            & $\alpha=0$
            & $\rho(x-y,\,x+y)=0$
            & $\mathcal{E}[x]=y$ \\\arrayrulecolor{lightgray}\hline
        Regression
            & $x_i - y_i = \alpha + \beta(y_i - \bar{y}) + \nu_i$
            & $x_i - y_i = \alpha + \beta(y_i + x_i) + \theta_i$
            & $x_i - y_i = \alpha + (\beta - 1)y_i + \epsilon_i - \beta\delta_i$
            \\\arrayrulecolor{lightgray}\hline
        Error term(s)
            & $\nu_i$
            & $\theta_i$
            & $\delta_i,\,\epsilon_i$ \\\arrayrulecolor{black}\hline\hline
        \makecell[l]{Mixed-Effects \\ \ \ Regression}
            & \makecell[c]{$x_{i,n} - y_{i,n} = (\alpha + u_i) +$ \\
                           $\, \beta(y_{i,n} - \bar{y}_i) + \nu_{i,n}$}
            & \makecell[c]{$x_{i,n} - y_{i,n} = (\alpha + u_i) +$ \\
                           $\,(\beta + v_i)(x_{i,n} + y_{i,n}) + \theta_{i,n}$}
            & \makecell[c]{$x_{i,n} - y_{i,n} = (\alpha + u_i) +$ \\
                           $\,(\beta - 1 + v_i)\,y_{i,n} +$ \\
                           $\;\epsilon_{i,n} - (\beta+v_i)\delta_{i,n}$} \\\arrayrulecolor{lightgray}\hline
        Fixed-Effect term(s)
            & $\alpha,\,\beta$
            & $\alpha,\,\beta$
            & $\alpha,\,\beta$ \\\arrayrulecolor{lightgray}\hline
        \makecell[l]{Fixed-Effect Null}
            & $\alpha = 0$
            & $\beta = 0$
            & $\alpha=0$ and $\beta=1$\textsuperscript{\dagger} \\\arrayrulecolor{lightgray}\hline
        Random-Effect term(s)
            & $u_i \sim \mathcal{N}(0,\tau^2_\alpha)$
            & \makecell[c]{$u_i \sim \mathcal{N}(0,\tau^2_\alpha)$, \\
                           $v_i \sim \mathcal{N}(0,\tau^2_\beta)$}
            & \makecell[c]{$u_i \sim \mathcal{N}(0,\tau^2_\alpha)$, \\
                           $v_i \sim \mathcal{N}(0,\tau^2_\beta)$} \\\arrayrulecolor{black}\hline
    \end{tabular}
    \caption{Protocol from Silveira~\cite{Silveira2024} to supplement Bland-Altman plots with inferential statistics to confirm equivalence (top block), and its proposed mixed-effects extension for model-level inference over the joint metabolite set (bottom block). Index $i$ denotes metabolite and $n$ denotes spectrum. $x$ and $X$ denote the observed and structural values for the \DL-based candidate method, whereas $y$ and $Y$ similarly correspond to the target or reference values. The measurement errors of $x$ and $y$ are denoted by $\epsilon$ and $\delta$, respectively. Random effects $u_i$ and $v_i$ capture between-metabolite heterogeneity in constant and proportional bias respectively. Their variances $\tau^2_\alpha$ and $\tau^2_\beta$ quantify the consistency of model performance across the metabolite profile. $^{\dagger}$Bonferroni correction applied to the joint intercept and slope assessment, giving the effective significance level $\alpha=0.025$ for each parameter. Standard errors on all fixed-effect estimates are cluster-robust, clustered on spectrum index $n$, to account for intra-spectrum cross-metabolite estimation correlation.}
    \label{tab:Silveira_table}
\end{table*}
\endgroup
\subsubsection{Equivalence Testing}
\label{subsubsec:silveira}
Silveira et al.~\cite{Silveira2024} proposed a three-part inferential protocol to accompany Bland--Altman plots, providing statistical support for claims of equivalence between measurement methods that the BA framework alone cannot supply. The protocol decomposes equivalence into three nested conditions (accuracy, precision, and agreement) each addressed by a dedicated hypothesis test. The tests are described as nested in the sense that the result of each step conditions the interpretation of the subsequent tests. A difference in accuracy can be accommodated through the translation correction described by Silveira et al., whereas disagreement in precision prevents a meaningful claim of full bisector agreement. The three tests, together with their structural and functional null hypotheses and their corresponding regression formulations, are summarized in Table~\ref{tab:Silveira_table}.

The structural null hypotheses are stated in terms of the true, unobservable values $X$ and $Y$, which are related to the observed measurements $x$ and $y$ through additive, zero-mean error terms $\epsilon$ and $\delta$ respectively. Because these error terms have null expectation by assumption, the observed measurements have the same expectations as their structural counterparts. Therefore, their sample means provide unbiased estimates of the corresponding structural means, and the structural hypotheses can be evaluated through their functional counterparts using standard regression.

\textcolor{black}{\textbf{Test~1 (Accuracy)}} assesses the equivalence of structural means by regressing the pairwise differences $y_i - x_i$ on the mean-centered reference values $x_i - \bar{x}$. Centering ensures that the regression intercept $\alpha$ corresponds directly to the mean difference $\bar{y} - \bar{x}$, so that $H_{0,1}: \alpha = 0$ is equivalent to the structural null $\text{E}[X] = \text{E}[Y]$. The slope $\beta$ in this regression characterizes the concentration-dependence of the bias and is not itself the subject of the null hypothesis, but a non-zero slope is indicative of proportional bias that, if present, should be noted before proceeding to the subsequent tests.

\textcolor{black}{\textbf{Test~2 (Precision)}} assesses the equivalence of structural error variances via the test of Shukla~\cite{Shukla1973}, which demonstrated that, under the specified structural measurement-error model, the slope of the regression of $y_i - x_i$ on $x_i + y_i$ is exactly zero if and only if the ratio of error variances $\lambda = \text{V}[\delta] / \text{V}[\epsilon] = 1$. The null hypothesis $H_{0,2}: \rho(x - y,\, x + y) = 0$ is therefore a test of equal measurement precision between the two methods. Notably, the axes of this regression are identical to those of the standard Bland--Altman plot, which implies that the original BA method, under optimal conditions, evaluates precision rather than full equivalence, and does so without the inferential support that Test~2 provides. In the presence of a non-zero bias established by Test~1, a translation correction is applied before evaluating the Test~2 null, as described in Silveira et al.~\cite{Silveira2024}.

\textcolor{black}{\textbf{Test~3 (Agreement)}} assesses concordance with the structural bisector line using Deming regression \cite{Linnet1998}, which accounts for measurement error in both variables by incorporating the error variance ratio $\lambda$ estimated in Test~2. The null hypothesis $H_{0,3}: \alpha = 0,\, \beta = 1$ requires the regression line to coincide with the identity, confirming that the two methods produce equal measurements for the same subject. Because the intercept and slope are estimated jointly, a Bonferroni correction is applied, reducing the effective significance level to $2.5\%$ per parameter to preserve the overall type~I error rate. As with Test~2, a translation correction for residual bias is applied when Test~1 is not satisfied. Section~\ref{sec:Reg:subsec:comparison:subsubsec:regression} describes the corresponding regression methods and their differing requirements for the \textit{in silico} and \textit{in vivo} comparison regimes.

Both analytical $p$-values and bootstrap confidence bands are reported for each test. Discordance between the two approaches can arise when a non-zero bias shifts regression lines outside the bootstrap confidence band. Silveira et al.~\cite{Silveira2024} recommend interpreting the analytical and graphical bootstrap results jointly, particularly when bias produces discordant conclusions. Consistent with the nonparametric framework applied throughout this work, the bootstrap is the preferred basis for inference given that normality of the differences cannot be assumed for \DL model outputs (Section~\ref{sec:eval:subsec:errorMetrics}).

The Silveira et al.\ protocol is applied on a per-metabolite basis. A more nuanced model-level extension of this framework, operating over the joint metabolite set, is described in Section~\ref{subsubsec:mixed_effects_silveira}.

\subsection{Regression-Based Comparisons}
\label{sec:Reg:subsec:comparison:subsubsec:regression}
Regression-based comparisons evaluate whether the relationship between \DL model outputs and the reference values is consistent with the identity line, testing jointly for constant bias (a non-zero intercept) and proportional bias (a slope different from unity). In both regimes, the \DL model output $x_i$ is treated as the response variable. In the \textit{in silico} regime, the ground truth concentration $y_i$ is the error-free predictor. In the \textit{in vivo} regime, the conventional fitting algorithm output serves as the predictor, with neither variable treated as error-free. The consequences of this distinction for the choice of regression estimator are detailed below. This distinction must also be accounted for when adapting Silveira's third test in Section~\ref{subsubsec:silveira}. 
Separate regression analyses are performed for each metabolite and each test set stratum defined in Section~\ref{sec:Reg:subsec:comparison:subsubsec:stratification}. 

\subsubsection{Passing-Bablok Regression}
Passing-Bablok regression~\cite{Passing1983,Bablok1984} is a nonparametric, rank-based estimator that does not require normally distributed measurement errors and is robust to outliers, making it well-matched to the nonparametric framework adopted throughout this work. It nevertheless assumes a linear relationship between the two methods and continuously distributed measurements. The estimator begins by computing all $M = N(N-1)/2$ pairwise slopes over ordered pairs $(i,j)$ with $i < j$:
\begin{equation}\label{eqn:pb_slopes}
    S_{ij} = \frac{x_j - x_i}{y_j - y_i}, \qquad i < j, \quad y_i \neq y_j
\end{equation}
\noindent where $y_i$ and $x_i$ are the reference and \DL model values for the $i$-th observation, $y_j \neq y_i$ is required, and pairs with equal reference values are discarded. The slope estimate is then the shifted median of the sorted values, $S_{ij}$:
\begin{equation}\label{eqn:pb_slope}
    \hat{\beta}_1 = S_{\left(\lfloor M/2 \rfloor + K + 1\right)}
\end{equation}
\noindent where $S_{(m)}$ denotes the $m$-th order statistic of the sorted slope set, $\lfloor \cdot \rfloor$ is the floor operator, and $K$ is the number of pairwise slopes satisfying $S_{ij} < -1$. The shift by $K$ is what distinguishes the Passing-Bablok estimator from a naive median-of-slopes: it ensures consistency when the true slope is negative by correcting for the asymmetric distribution of slopes about $-1$. The intercept is estimated as the median of the per-observation intercept residuals:
\begin{equation}\label{eqn:pb_intercept}
    \hat{\beta}_0 = \widetilde{\mu}\!\left(x_i - \hat{\beta}_1\, y_i\right)
\end{equation}
\noindent where $\widetilde{\mu}(\cdot)$ denotes the median operator as defined in Section~\ref{sec:eval:subsec:error:subsubsection:centroids}, and $\hat{\beta}_1$ is the slope estimate from Eqn.~\ref{eqn:pb_slope}.

The reference values of interest are $\beta_0 = 0$ (no constant bias) and $\beta_1 = 1$ (no proportional bias), 
with confidence intervals used to characterize whether the estimated intercept and slope are compatible with these values. The \textit{cumulative sum} linearity test should be reported alongside the regression results \cite{Passing1983,Bilic-Zulle2011}. A detected departure from linearity indicates a concentration-dependent relationship between the two methods that cannot be adequately characterized by a single slope and intercept, in which case the regression-based limits of agreement in Section~\ref{sec:Reg:subsec:comparison:subsubsec:ba} provide the primary characterization. Confidence intervals are constructed at the $95\%$ level and reported alongside point estimates for all parameters.

\subsubsection{Deming Regression}
When both methods carry non-negligible measurement error, as in the \textit{in vivo} regime, Passing-Bablok regression remains valid but Deming regression~\cite{Riu1997} provides a fully parametric alternative that incorporates the ratio of measurement error variances as a weighting term. In the \textit{in silico} regime, the reference is exact by construction and carries negligible error, so the error-in-variables framework is not applicable and Deming regression is therefore not required in that setting, and Passing-Bablok regression is retained as the primary robust estimator. The Deming estimator is obtained by minimizing a weighted sum of squared orthogonal residuals over both the model parameters and the latent true predictor values $\{\xi_i\}$:
\begin{equation}\label{eqn:deming_objective}
    \underset{\beta_0,\,\beta_1,\,\{\xi_i\}}{\mathrm{min}}
    \sum_{i=1}^{N}
    \left[
        \frac{\left(y_i - \xi_i\right)^2}{\sigma_y^2}
        +
        \frac{\left(x_i - \beta_0 - \beta_1\,\xi_i\right)^2}{\sigma_x^2}
    \right]
\end{equation}
\noindent where $y_i$ and $x_i$ are the observed values for the reference and \DL model respectively, $\xi_i$ are the unknown true values of the predictor, $\beta_0$ and $\beta_1$ are the intercept and slope, and $\sigma_y^2$ and $\sigma_x^2$ are the measurement error variances of the reference and \DL model, respectively. When $\lambda = \sigma_x^2 / \sigma_y^2 = 1$, Eqn.~\ref{eqn:deming_objective} reduces to ordinary orthogonal regression and as $\lambda \to 0$, it converges to ordinary least squares (OLS). The minimization of Eqn.~\ref{eqn:deming_objective} yields the closed-form slope estimate:
\begin{equation}\label{eqn:deming_slope}
    \hat{\beta}_1 =
    \frac{
        S_{yy} - \lambda\, S_{xx}
        + \sqrt{\left(S_{yy} - \lambda\, S_{xx}\right)^2
                + 4\,\lambda\, S_{xy}^{\,2}}
    }{
        2\, S_{xy}
    }
\end{equation}
\noindent with the corresponding intercept:
\begin{equation}\label{eqn:deming_intercept}
    \hat{\beta}_0 = \bar{x} - \hat{\beta}_1\,\bar{y}
\end{equation}
\noindent where $\bar{y}$ and $\bar{x}$ are the sample means of the reference and \DL model outputs, and $S_{xx}$, $S_{yy}$, and $S_{xy}$ are the corrected sample variances and covariance:
\begin{equation}\label{eqn:deming_variances}
    \begin{aligned}
        S_{xx} &= \frac{1}{N-1}\sum_{i=1}^{N}\left(x_i - \bar{x}\right)^2,
        \qquad \\
        S_{yy} &= \frac{1}{N-1}\sum_{i=1}^{N}\left(y_i - \bar{y}\right)^2,
        \qquad \\
        S_{xy} &= \frac{1}{N-1}\sum_{i=1}^{N}\left(y_i - \bar{y}\right)\!\left(x_i - \bar{x}\right)
    \end{aligned}
\end{equation}
\noindent where $N$ is the number of paired observations. The positive root in Eqn.~\ref{eqn:deming_slope} is taken throughout. 

As before, the error variance ratio is defined as:
\begin{equation}\label{eqn:deming_lambda}
    \lambda = \frac{\sigma_y^2}{\sigma_x^2}
\end{equation}
\noindent where $\sigma_y^2$ and $\sigma_x^2$ are the per-method measurement error variances.

The error variance ratio should be estimated from independent information about the measurement precision of both methods. For conventional fitting algorithms, a \CRLB-derived variance may provide a model-based precision proxy only under the assumptions discussed in Section~\ref{sec:Laying:subsec:Reliability} and should not be interpreted as directly observed measurement-error variance. For probabilistic \DL models, a calibrated predictive variance may provide a candidate estimate of the measurement-error variance associated with the \DL output, \(\sigma_x^2\). The validity of this use depends on the calibration of the predicted uncertainty, which should be established independently, for example using the analyses described in Section~\ref{sec:Eval:subsec:RegMetrics:subsubsec:NLPD}. Where the measurement-error variances vary substantially across spectra or concentration, a weighted Deming formulation that incorporates observation-specific precision estimates is preferable to a single constant error variance ratio. When sufficiently reliable estimates of the measurement-error variances are unavailable, Passing-Bablok regression remains the preferred robust comparison.

\subsection{Model-Level Multivariate Method Comparisons}
\label{sec:Reg:subsec:comparison:subsubsec:model-level}
The per-metabolite analyses described in Sections \ref{sec:Reg:subsec:comparison:subsubsec:ba} and~\ref{sec:Reg:subsec:comparison:subsubsec:regression} provide a detailed characterization of method agreement for each metabolite individually, and their results should be reported in full in the appendices. However, when a fitting model produces estimates for 20 to 30 metabolites simultaneously, a per-metabolite report is impractical as a primary basis for comparing two models. A metabolite profile produced by one fitting method constitutes a single multivariate output, and a model-level comparison requires summary quantities that characterize the agreement between two such profiles as a whole, rather than as a collection of independent scalar agreements. The analyses described in this section address this requirement directly. They operate over the joint set of metabolite outputs, using shared indexing and information-based references, and yield both a model-level scalar summary and model-level inferential statistics that are interpretable without reference to individual metabolites.

\subsubsection{SNR as a Common Visualization Axis}
\label{subsubsec:snr_axis}
Metabolite-specific \SNR (\mSNR) provides a physically motivated common axis on which the outputs of all metabolites can be visualized jointly. \mSNR characterizes the visibility of a metabolite-specific spectral component relative to the noise floor and therefore provides a physically interpretable common axis across metabolites with different absolute concentration scales. As established in the information-standardized analysis, however, \mSNR does not capture all factors contributing to intrinsic estimation difficulty, such as spectral overlap and parameter covariance.

This property is what enables joint visualization. In the per-metabolite BA plots, each metabolite occupies a separate panel because its concentration axis is incommensurable with those of other metabolites. When differences are instead plotted against \mSNR, all metabolites can be overlaid on a single set of axes, with each metabolite contributing a cloud of points along the shared \SNR dimension. The regression-based limits of agreement (Equation~\ref{eqn:ba_loa_heteroscedastic}) may then be visualized as functions of \mSNR rather than concentration to provide a shared descriptive representation of heteroscedastic disagreement across the metabolite set. This pooled representation complements, rather than replaces, the per-metabolite concentration-based analyses and the information-standardized characterization of intrinsic estimation difficulty.

In the \textit{in silico} regime, \mSNR (Section~\ref{sec:dataModel:subsec:curation:subsubsec:SNR}) is computed directly from the simulation parameters during data generation, making it exact by construction and making this regime the most reliable for method comparison. In the \textit{in vivo} regime, \mSNR is not a native output of standard fitting pipelines and must be derived post-hoc. The fitted parameters are used to modulate the basis functions, scaling and broadening each metabolite component to match the estimated lineshape, after which the power of each metabolite-specific basis function is measured directly from the modulated waveform. The model-specific \mSNR recoverability threshold identified during the characterization procedure described in Section~\ref{sec:dataModel:subsec:curation:subsubsec:Annotations} should be overlaid on the error-versus-\mSNR plots to indicate the regime below which reliable recovery is no longer supported. This threshold may additionally be used as a prespecified quality-control criterion for analyses restricted to the model's characterized operating regime, provided that results from the full dataset are also retained and the proportion of excluded observations is reported. Because the \mSNR values derived in this way depend on the fitted parameters rather than on ground truth concentrations, they carry the estimation bias of the fitting software. This limits the absolute accuracy of the \textit{in vivo} \mSNR estimates but does not preclude their use as a relative indexing axis for method comparison, provided that a single prespecified derivation procedure is used to assign the indexing value for each paired observation. 

\subsubsection{Information-Standardized Error and Model-Level Disagreement}
\label{subsubsec:efficiency_index}
Plotting raw differences against \SNR characterizes the \SNR-dependent agreement structure, but does not itself provide a model-level scalar suitable for comparing two fitting methods. In the \textit{in silico} regime, a dimensionless information-standardized disagreement measure can be obtained by normalizing each absolute difference by the oracle \CRLB-derived standard-deviation reference for the corresponding metabolite and simulated condition:
\begin{equation}\label{eqn:crlb_normalized_error}
    \varepsilon_{i,n} = \frac{|\hat{x}_{i,n} - y_{i,n}|}{\sqrt{\mathrm{CRLB}_{i,n}}}    
\end{equation}
\noindent where $i$ indexes the metabolites, $n$ indexes the spectra, $\hat{x}_{i,n}$ is the \DL model output, $y_{i,n}$ is the reference value, and $\mathrm{CRLB}_{i,n}$ is the Cram{\'e}r--Rao lower bound on the estimation variance for metabolite $i$ in spectrum $n$, as reported by the reference fitting algorithm, \hl{which can any validated or properly characterized algorithm}. The model-level standardized disagreement summary is then defined as the mean of $\varepsilon_{i,n}$ over all metabolites and spectra above a minimum \SNR threshold $\rho_{\min}$:
\begin{equation}\label{eqn:efficiency_index}
\begin{aligned}
    \bar{\varepsilon} &= \frac{1}{|\mathcal{S}|}
        \sum_{(i,n) \in \mathcal{S}} \varepsilon_{i,n}, \\
    \mathcal{S} &= \{(i,n) : \mathrm{SNR}_{i,n} \geq \rho_{\min}\}    
\end{aligned}
\end{equation}
\noindent where $|\mathcal{S}|$ is the cardinality of the set of qualifying observations. The threshold $\rho_{\min}$ excludes the low-SNR regime in which the \CRLB approximation degrades and the normalizing quantity becomes unreliable~\cite{Cavassila2001}. 

Because $\varepsilon_{i,n}$ is an absolute error standardized by a condition-specific precision reference, $\bar{\varepsilon}$ should not be interpreted as an estimator-efficiency ratio or as having a theoretical lower bound of unity. Even when an unbiased Gaussian estimator attains the corresponding \CRLB variance, the expected absolute standardized error is $\sqrt{2/\pi}$ rather than one. The distribution of $\varepsilon_{i,n}$ can nevertheless provide a descriptive summary of disagreement relative to the intrinsic precision scale. Claims concerning variance efficiency or empirical sub-\CRLB behavior should instead be based on the variance-standardized quantities defined in Section~\ref{sec:Analy&Visual:subsec:infoStand:subsubsec:subcrlb var}.

In the \textit{in silico} regime, $\sigma{\CRLB,i,n}$ is derived from the Fisher information evaluated at the known ground-truth parameters and therefore provides the same oracle information reference used in the preceding information-standardized analyses. In the \textit{in vivo} regime, an analogous normalization may be performed using an appropriately characterized precision estimate from the reference fitting method. In that setting, however, the resulting quantity is a reference-standardized disagreement measure rather than an oracle information-standardized quantity and should not be interpreted as evidence that the \DL estimator exceeds a theoretical bound. This distinction should be kept in mind when comparing $\bar{\varepsilon}$ values between the two regimes.

The standardized disagreement summary $\bar{\varepsilon}$ is reported as a descriptive model-level scalar for method comparison, alongside the sum of absolute biases across all metabolites established as the community standard by the ISMRM fitting challenge~\cite{Marjanska2022}, which is included to support direct comparison with published fitting benchmark results.

\subsubsection{Mixed-Effects Equivalence Testing}
\label{subsubsec:mixed_effects_silveira}
The three-test inferential protocol of Silveira et al.~\cite{Silveira2024} introduced in Section~\ref{subsubsec:silveira} is extended here to produce model-level equivalence statistics operating over the joint metabolite set. The extension replaces the per-metabolite regression with a mixed-effects regression in which metabolites are treated as a clustering unit, with random effects absorbing inter-metabolite heterogeneity so that the fixed-effect parameters estimate model-level quantities. The regression equations and hypothesis structure are summarized in the lower block of Table~\ref{tab:Silveira_table}.

In the pooled dataset, index $i$ denotes metabolite and index $n$ denotes the spectrum, so that each observation $x_{i,n} - y_{i,n}$ is a scalar difference for metabolite $i$ from spectrum $n$. When differences from all metabolites are pooled, each spectrum contributes $K$ observations, one per metabolite, that are correlated by the shared spectral noise and by peak overlap between metabolites in the same fit. This intra-spectrum cross-metabolite correlation is not modeled explicitly, because doing so would require identifying the full $K \times K$ cross-metabolite error covariance matrix, which is not identifiable from single measurements without extensive replication. 

Instead, all fixed-effect standard errors are computed using cluster-robust sandwich estimators~\cite{Liang1986,MacKinnon1985}, clustered on the spectrum index $n$. The sandwich estimator takes the form: 
\begin{subequations}\label{eqn:cluster robust}
\begin{align}
    \hat{V} &= (\mathbf{X}^\top \mathbf{X})^{-1} \hat{B} (\mathbf{X}^\top \mathbf{X})^{-1}\textnormal{, where}\label{eqn:sandwich estimator}\\[1pt]
    \hat{B} &= \sum_{n} \mathbf{X}_n^\top \hat{e}_n \hat{e}_n^\top \mathbf{X}_n\label{eqn:sandwich mean}
\end{align}
\end{subequations}%
\noindent where $\mathbf{X}$ is the regression design matrix, $\mathbf{X}_n$ is its submatrix for spectrum $n$, and the mean $\hat{B}$ is accumulated over spectra rather than over individual observations, with $\hat{e}_n$ denoting the vector of residuals from spectrum $n$ across all $K$ metabolites. This construction allows the $K$ residuals within a spectrum to be arbitrarily correlated without requiring their covariance structure to be modeled or estimated. The estimator requires independence across clustering units. In the \textit{in silico} regime, clustering on spectrum index $n$ is appropriate when simulated spectra are generated independently. In the \textit{in vivo} regime, multiple spectra from the same subject may be correlated and should therefore be clustered at the subject level, or at another justified independent acquisition unit, rather than treated as independent spectrum-level clusters. In order for the estimator to be reliable, a sufficient number of clusters is required. The adequacy of cluster-robust inference should therefore be assessed using the number of independent clusters, which corresponds to the number of independent spectra in the \textit{in silico} regime but may be closer to the number of subjects in the \textit{in vivo} regime~\cite{Cameron2008}. Calculating effective sample sizes is discussed in detail in Section~\ref{sec:Reg:subsec:comparison:subsubsec:samplesize}.

\textbf{Test~1 (Model-Level Accuracy).} The mixed-effects accuracy regression is:
\begin{equation}\label{eqn:me_test1}
    x_{i,n} - y_{i,n} = (\alpha + u_i)
        + \beta\,(y_{i,n} - \bar{y}_i) + \nu_{i,n}
\end{equation}
\noindent where $u_i \sim \mathcal{N}(0, \tau^2_\alpha)$ is a metabolite-level random intercept, $\bar{y}_i$ is the mean reference value for metabolite $i$, and $\nu_{i,n}$ is the residual. The fixed-effect intercept $\alpha$ estimates the model-level mean bias, averaged over the metabolite distribution, and the null $H_{0,1}: \alpha = 0$ is the model-level accuracy test. The fixed slope $\beta$ characterizes how mean bias changes across the reference measurement range at the model level, and a non-zero $\beta$ is indicative of proportional bias that is consistent across the metabolite profile. The random intercept variance $\tau^2_\alpha$ quantifies inter-metabolite heterogeneity in constant bias. A model may satisfy $H_{0,1}$ globally while exhibiting large $\tau^2_\alpha$ if positive biases in some metabolites cancel negative biases in others. Both $\alpha$ and $\tau^2_\alpha$ should therefore be reported.

\textbf{Test~2 (Model-Level Precision).} The mixed-effects precision regression is:
\begin{equation}\label{eqn:me_test2}
    x_{i,n} - y_{i,n} = (\alpha + u_i)
        + (\beta + v_i)(y_{i,n} + x_{i,n}) + \theta_{i,n}
\end{equation}
\noindent where $u_i \sim \mathcal{N}(0, \tau^2_\alpha)$ and $v_i \sim \mathcal{N}(0, \tau^2_\beta)$ are metabolite-level random intercept and slope deviations respectively, and $\theta_{i,n}$ is the residual. By the Shukla result~\cite{Shukla1973} underlying Test~2, the fixed slope $\beta$ is zero if and only if the model-level error variance ratio $\lambda = 1$, and the null $H_{0,2}: \beta = 0$ is the model-level precision test. The random slope variance $\tau^2_\beta$ quantifies inter-metabolite heterogeneity in the precision ratio. A model may exhibit equal average precision across the metabolite set while still showing substantial metabolite-to-metabolite variation in relative precision, a form of heterogeneity captured by $\tau^2_\beta$. As in the per-metabolite case, a translation correction for residual mean bias is applied when Test~1 is not satisfied prior to evaluating the Test~2 null.

\textbf{Test~3 (Model-Level Agreement).} The mixed-effects Deming regression is:
\begin{equation}\label{eqn:me_test3}
\begin{aligned}
    x_{i,n} - y_{i,n} =\ &(\alpha + u_i)
        + (\beta - 1 + v_i)\,y_{i,n}
        + \\&\epsilon{i,n} - (\beta + v_i)\,\delta_{i,n}
\end{aligned}
\end{equation}
\noindent where $u_i \sim \mathcal{N}(0,\tau^2_\alpha)$ and $v_i \sim \mathcal{N}(0,\tau^2_\beta)$ are metabolite-level random effects, and $\delta_{i,n}$ and $\epsilon_{i,n}$ are the per-method error terms. The combined residual $\epsilon_{i,n} - (\beta + v_i)\,\delta_{i,n}$ is heteroscedastic with respect to $v_i$, which complicates analytical inference and supports the use of bootstrap confidence intervals over analytical intervals for this test at the model level. The joint null $H_{0,3}: \alpha = 0$ and $\beta = 1$ is evaluated with a Bonferroni correction at an effective significance level of $2.5\%$ per parameter, consistent with the per-metabolite implementation. The error variance ratio $\lambda_{i,n} = \sigma^2_{\mathrm{DL},i,n} / \sigma^2_{\mathrm{ref},i,n}$ is specified per observation using method-specific precision estimates appropriate to the two measurements being compared. For the \DL output, $\sigma_{\mathrm{DL},i,n}^2$ should be obtained from calibrated predictive uncertainty or empirical repeated-measurement variability rather than from the oracle \CRLB reference. In the \textit{in vivo} regime, this requires a probabilistic or uncertainty-aware \DL model. For non-probabilistic \DL models, the proposed mixed-effects Passing-Bablok extension can provide an alternative that does not require specification of $\lambda_{i,n}$. Its inferential properties should, however, be established for that formulation rather than assumed to be identical to those of the mixed-effects Deming model.

The random effect variances $\tau^2_\alpha$ and $\tau^2_\beta$ from Test~3 provide the most complete characterization of inter-metabolite heterogeneity available within this framework. A model that passes the joint fixed-effect null but exhibits large $\tau^2_\beta$ achieves average bisector agreement while producing proportional relationships with the reference that vary systematically across the metabolite profile. This constitutes a clinically meaningful failure mode distinct from global model failure, and should be reported separately regardless of the outcome of the joint null test.

\section{Reporting and Reproducibility}
\label{sec:comm reporting}
Reproducible research in \DL-based \MRS requires standardized dissemination of models, code, and evaluation pipelines. Containerized implementations, such as Docker images, hosted in persistent, publicly accessible repositories help ensure consistent execution across computing environments and long-term usability. To further facilitate reuse, models should provide command-line interfaces with clearly documented input and output specifications, accompanied by example data loaders and minimal working examples.

Repositories should, at minimum, include: i) scripts for training models from scratch and, when available, loading pretrained weights; ii) the configuration and scripts required to reproduce training and, where supported, adapt the model to new datasets via customizable input and output dimensions; and iii) complete inference pipelines for deployment. Architectural modifications and dataset revisions should be tracked through version control to ensure transparency across development stages. Inclusion of evaluation scripts for metric computation and figure generation further strengthens methodological clarity and enables independent verification of reported results.

Whenever permitted by ethical and legal constraints, datasets should be shared alongside comprehensive metadata detailing acquisition parameters, preprocessing steps, and annotation protocols. For \textit{in vivo} datasets, publication of accompanying MRSinMRS datasheets \cite{Lin2021} is recommended, whereas synthetic datasets should include MRSsynMRS dataset sheets \cite{LaMaster2026}. Whether \textit{in vivo} or synthetic data is used, explicit reporting of the parameter ranges and empirical or sampling distributions represented in the training, validation, and test sets is essential, as these define the characterized data domain and provide the context required to interpret model generalization and failure outside that domain.

\section{Conclusion}
This work lays a unified methodological foundation for the rigorous evaluation of deep learning-based model fitting and metabolite quantification in magnetic resonance spectroscopy. By bringing together data curation principles, sample size considerations, statistical theory, characterization metrics, uncertainty analysis, and method comparison strategies, it addresses a major barrier to progress in the field: the lack of consistent and well-justified validation practices. The foundational framework is intended to improve reproducibility, comparability, and interpretability across studies, while also clarifying the assumptions and limitations that shape model performance. In this sense, the present work is not merely a catalog of techniques, but a structured validation pathway for determining where a model performs reliably, whether its uncertainty estimates are informative, and whether its outputs agree sufficiently with reference methods for the intended research or clinical application.

By establishing a shared vocabulary and a coherent evaluation structure, this work aims to support more transparent research and to strengthen the connection between methodological development and clinical translation. The long-term value of such a framework lies in its potential to improve how models are validated, compared, and reported across the \MRS community. In doing so, it offers a foundation on which future methodological advances can be judged more consistently and interpreted more reliably, supporting a more evidence-based progression toward clinical deployment. Ultimately, the goal is to establish the level of rigor needed for them to become trustworthy tools in research and clinical workflows.